\documentclass[11pt,a4paper]{article}

% -------------------- PACKAGES --------------------

\usepackage[utf8]{inputenc}
\usepackage[T1]{fontenc}
\usepackage{lmodern}
\usepackage[margin=1in]{geometry}
\usepackage{setspace}
\usepackage{titlesec}
\usepackage{enumitem}
\usepackage{booktabs}
\usepackage{array}
\usepackage{xcolor}
\usepackage{fancyhdr}
\usepackage{microtype}
\usepackage{xurl}
\usepackage{hyperref}

% -------------------- COLORS --------------------

\definecolor{accent}{HTML}{1F3A5F}

% -------------------- PAGE SETTINGS --------------------

\setstretch{1.15}
\setlength{\parindent}{0pt}
\setlength{\parskip}{4pt}
\setlength{\headheight}{14pt}

\hypersetup{
    colorlinks=true,
    linkcolor=accent,
    citecolor=accent,
    urlcolor=blue!60!black,
    pdftitle={Literature Review},
    pdfauthor={}
}

% -------------------- SECTION STYLE --------------------

\titleformat{\section}
  {\Large\bfseries\color{accent}}
  {}
  {0em}
  {}
  [\vspace{0.5em}\color{accent}\titlerule]
\titlespacing*{\section}{0pt}{0pt}{10pt}

% -------------------- LIST STYLE --------------------

\setlist{nosep, itemsep=2pt, topsep=3pt, parsep=0pt}
\setlist[itemize,1]{label=\textbullet, leftmargin=1.3em}
\setlist[itemize,2]{label=\textendash, leftmargin=1.3em}
\setlist[itemize,3]{label=\textasteriskcentered, leftmargin=1.3em}
\setlist[enumerate,1]{
    leftmargin=*,
    label=\textbf{\color{accent}\arabic*.},
    itemsep=10pt,
    topsep=4pt
}

% Heading for each numbered part of an entry
\newcommand{\entrypart}[1]{\textbf{\color{accent}#1}\par\vspace{2pt}}

% -------------------- HEADER & FOOTER --------------------

\pagestyle{fancy}
\fancyhf{}
\fancyhead[L]{\small Minor Project}
\fancyhead[R]{\small Literature Review}
\fancyfoot[C]{\small\thepage}
\renewcommand{\headrulewidth}{0.4pt}

% -------------------- DOCUMENT --------------------

\begin{document}

\section{1 Your Firmware Has Arrived: A Study of Firmware Update Vulnerabilities}

\begin{enumerate}

    % ---------- 1 ----------
    \item \entrypart{Paper Details}

    \begin{tabular}{@{}>{\bfseries}p{3.2cm} p{11.5cm}@{}}
        Authors          & Yuhao Wu, Jinwen Wang, Yujie Wang, Shixuan Zhai, Zihan Li, Yi He, Kun Sun, Qi Li, Ning Zhang \\[2pt]
        Publication Year & 2024 (USENIX Security '24, Philadelphia, Aug 14--16) \\[2pt]
        DOI/Link         & \url{https://www.usenix.org/conference/usenixsecurity24/presentation/wu-yuhao}
    \end{tabular}

    % ---------- 2 ----------
    \item \entrypart{Research Problem}

    \begin{itemize}
        \item Firmware updates fix bugs, but the update process is also an attack surface.
        \item Many update flaws are logic flaws: missing or weak verification (no version check, MD5 for integrity).
        \item These are not memory bugs, so existing firmware tools do not find them.
        \item Three challenges:
        \begin{itemize}
            \item \textbf{C1:} update paths cross front-end code, scripts, and binaries, with different languages and IPC types.
            \item \textbf{C2:} no standard for verification code, so it is hard to recognize.
            \item \textbf{C3:} static alerts have false positives. Testing a later check needs an input that passes the earlier checks.
        \end{itemize}
    \end{itemize}

    % ---------- 3 ----------
    \item \entrypart{Proposed Approach}

    \textbf{Systematization first}
    \begin{itemize}
        \item 4-phase update workflow: generation, delivery, verification, installation.
        \item 381 update-related CVEs mapped to these phases.
        \item Conclusion: focus on the verification phase, since good verification blocks most issues from other phases.
    \end{itemize}

    \textbf{ChkUp} (binary-only tool for OS-based firmware, mainly Linux, with file systems). Four modules:

    \begin{itemize}
        \item \textbf{Execution Path Recovery}
        \begin{itemize}
            \item Find entry program by pattern matching (prompt messages, variable/function names, common code).
            \item Build Update Flow Graph (UFG): basic-block level, covers HTML/JS, shell scripts, binaries. Edge types: intra-process flow, IPC, program invocation.
            \item Backward slicing from reboot calls gives update paths.
            \item Path explosion handled by: skip std libs, skip built-in utilities, timeout, path filtering by prompts, path merging.
        \end{itemize}
        \item \textbf{Verification Procedure Recognition}
        \begin{itemize}
            \item Stage 1: rank functions by similarity to a corpus. Uses numeric syntactic and structural features. Threshold 0.5.
            \item Stage 2: DFG subgraph isomorphism (Ullmann). Also checks that return values or by-reference args feed a conditional.
            \item Corpus: 129 functions (76 from libraries, 53 proprietary).
        \end{itemize}
        \item \textbf{Vulnerability Discovery}
        \begin{itemize}
            \item Checks authenticity, integrity, freshness, compatibility.
            \item Missing verification is reported.
            \item Improper verification:
            \begin{itemize}
                \item authenticity: symmetric algorithms (HMAC, CMAC, Poly1305)
                \item integrity: CRC, SHA1, MD5
                \item freshness/compatibility: info from unprotected sources (e.g., filename)
            \end{itemize}
            \item Multiple paths: pick the one with the most proper verification.
        \end{itemize}
        \item \textbf{Vulnerability Validation}
        \begin{itemize}
            \item Only for emulatable firmware. Build a PoC image per alert type.
            \item Patch other checks (invert branches, or binary rewriting with recorded values) and repack.
            \item Vulnerable if the update still reaches reboot.
            \item Tools: Firmadyne, FirmAE, angr, Ghidra, NetworkX, Firmware-mod-kit.
            \item Non-emulatable images: manual validation.
        \end{itemize}
    \end{itemize}

    % ---------- 4 ----------
    \item \entrypart{Dataset and Evaluation}

    \begin{itemize}
        \item \textbf{Collection:} 157,141 images from 204 vendors. 111,958 unpacked.
        \item \textbf{$D_L$:} 12,000 random images, 8 vendors (Netgear, TP-Link, D-Link, TRENDnet, Asus, Ubiquiti, Zyxel, Linksys).
        \item \textbf{$D_G$ (ground truth):} 150 images from $D_L$, weighted random sampling, 33 device families. Labeled by 4 security researchers with cross-validation.
        \item \textbf{Setup:} Ubuntu 18.04, AMD EPYC 7302P, 64GB RAM.
        \item \textbf{Metrics:}
        \begin{itemize}
            \item Entry finding: accuracy, with ablation of pattern types P1/P2/P3.
            \item UFG: soundness and completeness.
            \item Procedure recognition: TP, FN, FP.
            \item Validation: patch, emulation, and PoC success rates.
            \item Overall: TPR, FPR.
            \item Analysis time.
        \end{itemize}
        \item Validation scalability tested on 1,200 more images from $D_L$.
        \item \textbf{Baseline comparison:} Not Specified.
    \end{itemize}

    % ---------- 5 ----------
    \item \entrypart{Key Findings}

    \begin{itemize}
        \item \textbf{CVE study}
        \begin{itemize}
            \item CVE count doubled from 2020 to 2021.
            \item High and critical CVEs are nearly 4x low and medium.
            \item Top 8 CWE types (47.67\%) are mostly missing or improper verification.
        \end{itemize}
        \item \textbf{Entry finding:} best with all patterns. P2 (variable/function names) matters most. 9 Netgear images had wrong entries.
        \item \textbf{UFGs (150 total)}
        \begin{itemize}
            \item 122 sound and complete. 136 sound, 133 complete.
            \item Unsound: 7 Asus, 4 D-Link, 3 Zyxel.
            \item Incomplete: 9 Netgear (wrong entry), 3 Zyxel (entry/back-end mismatch), 5 TP-Link (timeout).
        \end{itemize}
        \item \textbf{Recognition:} 461 TP, 45 FN, 17 FP. Integrity checks recognized best. Bad paths caused 23 FN and 6 FP.
        \item \textbf{Time per image:} 126.0 s (path recovery), 216.1 s (recognition).
        \item \textbf{Validation on $D_G$}
        \begin{itemize}
            \item 271 alerts. 119 on 72 emulatable images. 90 needed patching.
            \item 69 patched images built. 88 test environments built.
            \item 73 PoCs worked (TP). 6 were FP. 9 failed but should be TP.
            \item 9 of the 15 failures: wrong verification procedure chosen for bypass.
        \end{itemize}
        \item \textbf{Scalability (1,200 images):} 44.1\% emulatable. 72.2\% of those repacked. 82.7\% of repacked ones emulated. Main failure cause: runtime signature checks.
        \item \textbf{Large-scale run on $D_L$}
        \begin{itemize}
            \item 93.4\% done within 600 s.
            \item Paths resolved for 10,670 images. 15,132 alerts. 1,330 images unresolved.
            \item Most common issue: improper integrity verification (39.0\% CRC, 23.8\% MD5).
            \item Most alerts: network devices and cameras.
        \end{itemize}
        \item \textbf{Overall on $D_G$:} TPR 86.7\%, FPR 5.3\%.
        \item \textbf{Disclosure:} 25 CVE IDs, 1 PSV ID. Mostly CWE-345 and CWE-20. Asus results were n-day (CVE-2014-2718, CVE-2020-15498, CVE-2021-3166). 29 other device families were zero-day.
        \item \textbf{Case study (Netgear WNR)}
        \begin{itemize}
            \item Freshness and compatibility checked via filename on front end.
            \item Integrity: CRC. Authenticity: missing.
            \item Built a downgrade attack and a modification attack.
            \item Both work over the same network or by MITM, since update uses plain HTTP.
        \end{itemize}
    \end{itemize}

    % ---------- 6 ----------
    \item \entrypart{Contribution and Limitations}

    \textbf{Contributions (authors)}
    \begin{itemize}
        \item 4-phase systematization of update security from 381 CVEs.
        \item ChkUp: handles diverse update components, recognizes verification procedures, validates alerts.
        \item Real-world results: 12,000 images analyzed, 150 validated, 25 CVEs and 1 PSV ID.
        \item Two working exploits on a Netgear router.
    \end{itemize}

    \vspace{4pt}
    \textbf{Limitations (authors)}
    \begin{itemize}
        \item Entry and path heuristics can fail on customized firmware. Analysis then cannot start.
        \item Verification heuristic fails on unconventional designs.
        \item Ground truth is only an approximation (manual work).
        \item Static analysis: false alarms.
        \item Dynamic validation: one path at a time, needs emulation, needs correct patching.
        \item Only 72 of 150 images emulatable. Even lower for bare-metal.
        \item Only OS-based firmware with file systems. Bare-metal not supported.
        \item Unsupported: PHP in update steps, unusual IPC.
        \item Repacking fails on uncommon file systems.
        \item Validation is slow and manual, so limited to 150 images.
    \end{itemize}

    % ---------- 7 ----------
    \item \entrypart{My Analysis}

    \textbf{Good}
    \begin{itemize}
        \item Targets logic bugs in update verification, a gap in firmware analysis.
        \item Large-scale run, real CVEs, working exploits.
    \end{itemize}

    \vspace{4pt}
    \textbf{Weak points}
    \begin{itemize}
        \item Only 150 of 12,000 images validated (about 1\%). The 15,132 alerts were not checked. TPR/FPR only hold for $D_G$.
        \item Entry patterns, corpus, and outlier exceptions came from manual analysis of the same eight vendors. Not stated whether $D_G$ was kept separate from this design work. Results may be optimistic.
        \item Only 8 router/camera vendors. Other vendors and device types: Not Specified.
        \item Heavy heuristic dependence (entry patterns, value-to-condition rule, corpus). Manual exceptions do not scale. Obfuscated or encrypted update code is not discussed.
        \item Multi-path images: only the ``best verified'' path is kept. May miss weaker paths (recovery mode, other interfaces). Effect on FN not measured.
        \item Covers only the verification phase. Delivery (HTTP, hard-coded keys), installation bugs, and bootloader checks are not detected, though the paper lists them as real issues.
        \item Exploits shown on one device family only.
        \item TPR/FPR formulas and denominators not clearly given. No comparison with other tools.
    \end{itemize}

    \vspace{4pt}
    \textbf{Research gaps}
    \begin{itemize}
        \item Bare-metal/RTOS support, and languages like PHP.
        \item Better emulation and repacking (uncommon file systems, runtime signature checks).
        \item Automatic discovery of entry points and verification functions, with less manual work.
        \item Full update chain analysis: delivery, companion apps, bootloader, and verification together.
        \item Bigger, independent ground truth and baseline comparison.
        \item Checking whether vendor patches actually fix the issues. Not studied.
    \end{itemize}

\end{enumerate}

\newpage

\section{Paper 2: Unveiling the Global Landscape of Android Security Updates}

\begin{enumerate}

    \item \entrypart{Paper Details}

    \begin{tabular}{@{}>{\bfseries}p{3.2cm} p{11.5cm}@{}}
        Authors          & Haiyun Deng, G\"{u}liz Seray Tuncay, Abbas Acar, Esteban Luques, Harun Oz, Ahmet Aris, Selcuk Uluagac \\[2pt]
        Publication Year & 2026 (IEEE Transactions on Dependable and Secure Computing, Vol.\ 23, No.\ 3, May/June 2026; received 23 November 2024, revised 23 November 2025, accepted 30 December 2025, published 3 February 2026) \\[2pt]
        DOI/Link         & \url{https://doi.org/10.1109/TDSC.2026.3660642}
    \end{tabular}

    \item \entrypart{Research Problem}

    \begin{itemize}[leftmargin=*,itemsep=1pt,topsep=2pt]
        \item \textbf{Problem:} Android security update speed and availability are still a problem, even though AOSP made updates more transparent through Security Patch Levels (SPL).
        \item \textbf{Why it matters:}
        \begin{itemize}[itemsep=0pt]
            \item Android has over three billion active devices.
            \item Google (2022): half of the zero-day vulnerabilities discovered in the wild were variations of vulnerabilities already patched.
            \item As of 2020, over one billion Android devices no longer receive security patches (cited by the paper).
        \end{itemize}
        \item \textbf{Challenges listed by the paper:}
        \begin{itemize}[itemsep=0pt]
            \item Rollout is complex: Google/ASB $\rightarrow$ OEM customization and testing $\rightarrow$ carrier approval for branded devices $\rightarrow$ OTA.
            \item Samsung example: 1728 models in 423 devices, 97 countries, 109 carriers $\rightarrow$ 25 K security update variations to customize and test regularly.
            \item Inconsistent support: delays or complete failures for some models.
            \item OEMs communicate support policies unclearly; users complain in forums (region-, model-, or carrier-specific).
        \end{itemize}
        \item \textbf{Gap in prior work (authors' view):}
        \begin{itemize}[itemsep=0pt]
            \item Earlier studies looked at one vulnerability type (e.g., kernel) or at stakeholder roles.
            \item They overlooked the full device life-cycle and non-flagship devices.
            \item They mostly blamed fragmentation, ignoring support type, device model, region.
            \item Jones et al.\ report a 24-day delay; the paper says this holds only for monthly-supported devices (41 and 63 days for quarterly or biannual).
        \end{itemize}
        \item \textbf{Research questions:}
        \begin{itemize}[itemsep=0pt]
            \item RQ1: Maintenance chronicle of devices over their life-cycle.
            \item RQ2: Factors that affect update distribution during the support period.
            \item RQ3: Risks of unpatched devices and when they become unsafe.
            \item RQ4: Issues behind user concerns and exemplary OEM practices.
            \item RQ5: How far each OEM follows Android security guidelines.
            \item RQ6: Prevalent vulnerabilities and weaknesses on Android devices.
        \end{itemize}
        \item \textbf{Relation to prior work:} Extends the authors' earlier work [20] (RQ1--RQ4, four OEMs) with four more OEMs (Huawei, Vivo, Motorola, LG), manual firmware inspection, OEM compliance analysis, and CWE analysis (RQ5, RQ6).
    \end{itemize}

    \item \entrypart{Proposed Approach}

    \begin{itemize}[leftmargin=*,itemsep=1pt,topsep=2pt]
        \item \textbf{Main idea:} Device-centric measurement study of Android security updates over the full device life-cycle, including the period after support ends. Uses official OEM data only.
        \item \textbf{Stage 1: Data collection}
        \begin{itemize}[itemsep=0pt]
            \item Sources: ASBs (August 2015 to May 2024), OEM update pages, OEM security bulletins, OEM support lists.
            \item Selection criteria: data must be official, contain complete device history, and be large enough to be representative.
            \item Samsung: URLs built from 1728 official models and 314 community-sourced CSC values; over half a million pairs queried.
            \item Xiaomi: official API, plus a third-party tracker site for history (fetches updates every six hours since 2018).
            \item Oppo: regional websites; all firmware downloaded to extract SPL from config files.
            \item Google: OTA images from Google's factory-image site; SPL via Google build numbers.
            \item Vivo, Huawei: firmware downloaded; SPL and Android version extracted.
            \item Motorola, LG: limited to CVE records and partial update history.
            \item Support lists: Wayback Machine snapshots.
        \end{itemize}
        \item \textbf{Stage 2: Verification}
        \begin{itemize}[itemsep=0pt]
            \item Samsung, Google: extracted SPLs cross-validated against official records.
            \item Oppo, Vivo, Xiaomi: manual inspection of sample firmware files to confirm SPL and Android version; CVEs checked against OEM change logs.
            \item Manual spot checks across device families and regions.
            \item Dataset and analysis code released in a public repository.
        \end{itemize}
        \item \textbf{Stage 3: Support-period metrics}
        \begin{itemize}[itemsep=0pt]
            \item Update count, support duration, update frequency (average interval), release delay (AOSP SPL release to OEM firmware).
            \item Support type (monthly, quarterly, biannual) derived from support-list snapshots.
        \end{itemize}
        \item \textbf{Stage 4: Impacting-factor analysis}
        \begin{itemize}[itemsep=0pt]
            \item Factors: support type, geolocation, device type, carrier association, partnership agreements and platform solutions (Android Enterprise Recommended, Samsung Knox).
            \item Geolocation: country/region extracted from Samsung CSC via ISO definitions; top five vs.\ bottom five countries compared.
        \end{itemize}
        \item \textbf{Stage 5: Unsupported-period risk analysis}
        \begin{itemize}[itemsep=0pt]
            \item EOL date taken as the last security update received (OEMs do not announce EOL dates).
            \item CVEs matched to each device by last Android version and chipset; CVEs from AOSP, Samsung, and NIST merged by CVE ID.
            \item CVEs published before the device's EOL excluded.
            \item CVEs grouped into 3-month ``exploit-quarters'' by age relative to EOL.
            \item Cumulative weighted average: $CWA = \sum_{j=1}^{n} \frac{\sum_{i=1}^{k_j} m_{ij}}{k_j}$, with $m_{ij}$ the metric (e.g., CVE count) for device $i$ in period $j$, $k_j$ the number of active devices with unpatched vulnerabilities in period $j$.
            \item Severity, user interaction, attack complexity, attack vector taken from CVSS 3.x labels.
            \item CVEs linked to CWEs via the NIST database.
        \end{itemize}
        \item \textbf{Stage 6: OEM practice evaluation}
        \begin{itemize}[itemsep=0pt]
            \item Compare ASB CVEs with OEM-published CVEs: \textit{intersection CVEs} vs.\ \textit{missing CVEs}.
            \item CVEs filtered by OEM-supported Android versions and chipsets (third-party sources).
            \item Metrics: Average Date Difference (ASB date vs.\ OEM date) and Average Delay.
        \end{itemize}
        \item \textbf{Stage 7:} Key issues and exemplary practices derived from the findings.
    \end{itemize}

    \item \entrypart{Dataset and Evaluation}

    \begin{itemize}[leftmargin=*,itemsep=1pt,topsep=2pt]
        \item \textbf{Overall:} $\sim$567 K security updates (Table I total: 567728), 904 devices, 8 OEMs, 98 countries, 109 carriers, 2014 to 2024.
        \item \textbf{Update histories (Table I):}
        \begin{itemize}[itemsep=0pt]
            \item Samsung: 542592 updates, 423 devices, 97 countries, 109 carriers, Apr 2015--May 2024.
            \item Xiaomi: 4800, 251 devices, 11 regions, Dec 2014--May 2024.
            \item Oppo: 18836, 77 devices, 35 countries, Jan 2018--Aug 2022.
            \item Google: 1015, 45 devices, Nov 2015--May 2024.
            \item Vivo: 288, 72 devices, 40 countries, Mar 2021--Apr 2024.
            \item LG: 197, 36 devices, 1 country, Dec 2016--Aug 2022.
            \item Huawei, Motorola: no full update history (Motorola update-history analysis omitted).
        \end{itemize}
        \item \textbf{ASB data:} 5901 CVEs, 5150 unique CVE IDs.
        \item \textbf{CVE/CWE data:} 9659 unique CVEs across all datasets, associated with 170 unique CWEs; 837 CWEs queried across 24980 CVEs.
        \item \textbf{Support lists:} e.g., Samsung 849 snapshots (Oct 2017--May 2024), 340 unique devices; Oppo 18 snapshots, 329 devices; Xiaomi 8 snapshots, 52 devices; Huawei 57 snapshots, 257 devices; Vivo 9 snapshots, 173 devices.
        \item \textbf{Unsupported-period sample:} 27349 pairs from six OEMs; 6787 unique CVEs assigned by matching Android versions and chipsets.
        \item \textbf{Ground truth:} Not Specified (no labeled ground truth; accuracy checked through cross-validation and manual spot checks, sample sizes Not Specified).
        \item \textbf{Train/test splits:} Not Specified (descriptive measurement study).
        \item \textbf{Setup/metrics:}
        \begin{itemize}[itemsep=0pt]
            \item Update count, support duration, update frequency, release delay.
            \item Average Date Difference, Average Delay (CVE publication, ASB vs.\ OEM).
            \item Accumulated CVEs per device per exploit-quarter; CWE cumulative ratio.
            \item Statistical significance tests: Not Specified.
        \end{itemize}
        \item \textbf{Baseline comparison:} Not Specified. (Only textual comparison with prior studies, e.g., Jones et al.\ 24-day delay, and with the authors' earlier work [20].)
    \end{itemize}

    \item \entrypart{Key Findings}

    \begin{itemize}[leftmargin=*,itemsep=1pt,topsep=2pt]
        \item \textbf{Samsung (support period):}
        \begin{itemize}[itemsep=0pt]
            \item About 21 security updates per device, including about 2.57 major OS upgrades.
            \item Average support duration about 858 days; average update frequency about 46.35 days.
            \item Average release delay 125 days (24398 pairs); 409 of 495 devices under 70 days, the other 86 range from 71 to 2147 days.
            \item Only 29 of 423 devices received updates for over four years; 18 pairs over six years (13 South Korea, four Taiwan, one Hong Kong).
        \end{itemize}
        \item \textbf{Xiaomi:} 902 pairs; average five updates; support duration 298 days (China 355); update every 52 days; release delay 32 days.
        \item \textbf{Oppo:} 1183 pairs, 77 devices; average 16 updates; support duration 584 days; update frequency about 47 days; release delay 35 days. Different Reno4 models got no updates for over a year despite being on the quarterly list.
        \item \textbf{Google:}
        \begin{itemize}[itemsep=0pt]
            \item Monthly updates, no missed SPLs until March 2023.
            \item After that, updates stopped for all but six newer models (Pixel 7a, Tablet, Fold, 8, 8 Pro, 8a).
        \end{itemize}
        \item \textbf{Support type (Samsung):}
        \begin{itemize}[itemsep=0pt]
            \item Devices move from monthly to quarterly to biannual lists.
            \item Time on list: 1062 (monthly), 767 (quarterly), 442 (biannual) days; total 2271 days ($\sim$6 years).
            \item Average security updates: 35, 18, four.
            \item Median support duration: 1044, 675, 553 days.
            \item Median update frequency: nearly 30, 45, 90 days (expected 30, 90, 180).
            \item Average release delay: 25, 38, 57 days; up to 300 days for some models.
            \item Oppo: most devices received only three updates in quarterly support.
        \end{itemize}
        \item \textbf{Geolocation:}
        \begin{itemize}[itemsep=0pt]
            \item Samsung monthly support: top five countries get twice as many updates; bottom five (four in Latin America and the Caribbean) face delays three times longer.
            \item Bottom five: only 7--8 updates, release delay around 48 days vs.\ about 18 days for the top five; support duration similar ($\sim$3 years).
            \item Xiaomi: China support duration 358 days vs.\ 255 elsewhere.
            \item Oppo: Europe 18 updates every 37 days; other regions 13 and 16 updates every 47--48 days.
        \end{itemize}
        \item \textbf{Device type:}
        \begin{itemize}[itemsep=0pt]
            \item No Samsung tablet has monthly support.
            \item Smartphones supported longer: 2140 vs.\ 1114 days for tablets (quarterly plus biannual).
            \item Not the same for Google or Xiaomi.
        \end{itemize}
        \item \textbf{Carrier:} Samsung carrier-branded vs.\ non-carrier: 34 vs.\ 37 monthly updates; other differences under 10\% or within 7\%. Google names the carrier in 32\% of updates (326). No significant differences observed.
        \item \textbf{Partnerships:}
        \begin{itemize}[itemsep=0pt]
            \item AER (Samsung): of 87 AER devices, 70 in support lists, only 15 completed monthly and quarterly support, none completed biannual.
            \item AER devices stay on monthly and quarterly lists longer: 1191 vs.\ 1002 days and 1223 vs.\ 763 days.
            \item Xiaomi AER: no significant differences vs.\ non-AER.
            \item Knox (Samsung): 218 Knox, 33 Android Others, four Android Go devices. Android Go gets 96 days of biannual support on average, Knox 291 days.
        \end{itemize}
        \item \textbf{Unpatched devices:}
        \begin{itemize}[itemsep=0pt]
            \item 13974 ($\sim$51\%) of 27349 pairs received at least one CVE affecting a device without security updates in the last three months.
            \item Average CVEs per unpatched device: about 114 in the first quarter, 664 by the eighth exploit-quarter (two years), 956 by the 20th (five years).
            \item About 89.7\% exploitable without user interaction; high-complexity CVEs under 6\%.
            \item About 74.99\% local, around 21.15\% network.
            \item Unpatched devices expose users to 167 unique CWEs.
            \item Takeaway-2: critical CVEs continue to arrive, remotely exploitable without user interaction; emerge right after support ends and decline after two years.
        \end{itemize}
        \item \textbf{Case study (Galaxy S10 5G):}
        \begin{itemize}[itemsep=0pt]
            \item Five model--CSC pairs with no update for nearly 15 months; last update Android 12, Exynos 9820 chipset.
            \item 873 Android-version CVEs and eight chipset CVEs for two pairs; 1077 for the other three.
            \item 11 critical and 32 high-severity CVEs from network attack vectors, remotely exploitable without user interaction (e.g., CVE-2023-21287, code execution).
        \end{itemize}
        \item \textbf{OEM compliance (missing CVEs, ASB vs.\ OEM site):}
        \begin{itemize}[itemsep=0pt]
            \item Samsung 75, Huawei 1950, Xiaomi 18, Oppo 577, Vivo 164, Motorola 105, LG 780. Google omitted (Pixel bulletins are supplementary).
            \item Average Date Difference (days): Samsung 46.61, Huawei 43.35, Xiaomi 54.61, Oppo $\sim$157.19, Vivo about 121.88, Motorola about 32.87, LG about 74.
            \item Average Delay (days): Samsung 23.21, Huawei $\sim$13.85, Xiaomi 1.067, Oppo 1.66, Motorola 19.53, LG 52.56.
            \item Vivo: no delays found.
        \end{itemize}
        \item \textbf{Key issues (examples):}
        \begin{itemize}[itemsep=0pt]
            \item Galaxy S20 FE 5G: \texttt{SM-G781B/XEH} last updated October 2020 (Android 10); \texttt{SM-G781B/BGL} last updated May 2024 (Android 13).
            \item Galaxy M40 (\texttt{SM-M405F/SLK}): largest gap between last update and biannual support list date, 1492 days; 177 devices received only one update, 25 still listed with last update over 1000 days ago.
            \item Samsung AER: of 8220 pairs, 63.76\% (5241) updated in the last two months; 317 pairs had none in the past year, 129 in two years, 38 in three years.
            \item Samsung's February 2021 announcement: 85 devices, 7275 pairs; 493 pairs had no updates in the past year.
        \end{itemize}
        \item \textbf{Exemplary practices named:} guaranteed support dates (only Motorola and Pixel), end-of-support lists (Motorola, Xiaomi), model-level support lists (only Oppo), full patch lists with reasons for exclusions.
    \end{itemize}

    \item \entrypart{Contribution and Limitations}

    \textbf{Contributions (authors)}
    \begin{itemize}[leftmargin=*,itemsep=1pt,topsep=2pt]
        \item Largest official dataset in the literature: $\sim$567 K security updates, 904 devices, 98 countries, 109 carriers.
        \item Identified factors affecting update availability and timeliness: support type, geolocation, device type, OEM partnerships.
        \item Studied common vulnerabilities and weaknesses (CVEs, CWEs) across Android devices, mainly smartphones.
        \item First analysis of the security posture of unpatched Android devices.
        \item Compared OEM adherence to Android security best practices using public ASBs and OEM websites.
        \item Identified key issues for users and highlighted best practices from some OEMs.
        \item Released dataset and code for reproduction.
    \end{itemize}

    \vspace{4pt}
    \textbf{Limitations (authors)}
    \begin{itemize}[leftmargin=*,itemsep=1pt,topsep=2pt]
        \item Older models or discontinued lines may lack complete archival data.
        \item Huawei, Vivo, Motorola, LG data are sparse, fragmented, or incomplete; used for overarching patterns, not detailed quantitative conclusions.
        \item Oppo data lack China, the US, and the U.K.
        \item A CVE missing on an OEM site is not conclusive proof of unpatched status; OEMs may silently integrate patches.
        \item Samsung support durations are a lower bound (devices may keep receiving updates).
        \item EOL dates inferred from last update, as OEMs do not announce them.
        \item Oppo's last Android version matched from a third-party dataset (lower bound).
        \item No active user base data; end-user behavior and carrier analysis out of scope.
        \item Reasons for late or missing updates cannot be determined (no public OEM explanations).
        \item CVSS 4.0 user-interaction reclassification not possible at scale.
        \item Other OEMs (e.g., Nokia, OnePlus, Realme, Sony) and international carriers left to future work.
    \end{itemize}

    \item \entrypart{My Analysis}

    \textbf{Good}
    \begin{itemize}[leftmargin=*,itemsep=1pt,topsep=2pt]
        \item Public, official sources, with dataset and code released.
        \item Covers the full life-cycle, including the unsupported period.
        \item Compares support lists (what OEMs promise) with actual updates received.
        \item Links update behavior to CVEs and CWEs, with a concrete case study.
        \item Manual verification of firmware SPLs.
        \item Authors are open about limits (e.g., missing CVE is not proof of unpatched).
    \end{itemize}

    \vspace{4pt}
    \textbf{Weak points}
    \begin{itemize}[leftmargin=*,itemsep=1pt,topsep=2pt]
        \item \textbf{Validated vs.\ analyzed:}
        \begin{itemize}[itemsep=0pt]
            \item Observation: only a manual ``spot check'' is described for verification; sample size and error rate are Not Specified.
            \item Unpatched-device CVEs are matched by Android version and chipset. This shows potential exposure, not verified exploitability on the device.
        \end{itemize}
        \item \textbf{Vendor bias:}
        \begin{itemize}[itemsep=0pt]
            \item Samsung is 542592 of the 567728 total updates, so pooled averages are likely dominated by Samsung.
            \item Google has only 26 devices in the unpatched analysis.
            \item Abstract and conclusion say eight OEMs, but detailed analysis covers four.
        \end{itemize}
        \item \textbf{Heuristics and manual steps:}
        \begin{itemize}[itemsep=0pt]
            \item Samsung URLs built from a community-sourced CSC list (314); completeness is unclear.
            \item Xiaomi history comes from a third-party tracker, and chipset and Android-version matching uses third-party sources. This sits uneasily with the ``official data only'' claim (observation).
            \item Support periods inferred from list snapshots.
        \end{itemize}
        \item \textbf{Metric definitions:}
        \begin{itemize}[itemsep=0pt]
            \item Two ``delay'' notions are used (firmware release delay vs.\ CVE publication delay), which can confuse readers (observation).
            \item Average Delay excludes negative differences, which may bias the delay upward (observation).
        \end{itemize}
        \item \textbf{Statistics:} ``Significant'' differences are reported without tests, as far as the text shows (observation).
        \item \textbf{Confounding:} Knox and AER membership may correlate with device tier or price, but this is not controlled (observation).
        \item \textbf{Claim wording:} The abstract pairs ``without user interaction'' with ``low attack complexity''. The body reports the two shares separately, and a joint figure is not shown (observation).
        \item \textbf{Single case:} Unpatched-risk case study uses one device (Galaxy S10 5G, five pairs).
        \item \textbf{Internal number inconsistencies (observation):}
        \begin{itemize}[itemsep=0pt]
            \item Huawei CVEs: 2864 (Sec.~III), 2859 (Sec.~IV, VII).
            \item LG CVEs: 3919, 4035, 4011 in different sections.
            \item Pixel CVEs: 1921 vs.\ 2817.
            \item Oppo pairs: 1137 (Sec.~III) vs.\ 1183 (Sec.~V).
            \item Samsung device counts: 423, 427, 495.
            \item Xiaomi China support duration: 355 vs.\ 358 days.
        \end{itemize}
        \item \textbf{Baselines:} No formal baseline; no comparison with a ground-truth patch audit (e.g., binary-level check) for missing CVEs.
    \end{itemize}

    \vspace{4pt}
    \textbf{Research gaps}
    \begin{itemize}[leftmargin=*,itemsep=1pt,topsep=2pt]
        \item Root causes of regional gaps, support-list discrepancies, and missing CVEs (the paper says it cannot determine them).
        \item Validating whether ``missing'' CVEs are really unpatched, e.g., via firmware or binary analysis.
        \item Checking real exploitability of unpatched-device CVEs; CVSS 4.0 re-scoring (authors' future work).
        \item Weighting risk by active user base (authors note this data is missing).
        \item Fuller coverage of Huawei, Vivo, Motorola, LG, and other OEMs (Nokia, OnePlus, Realme, Sony), and international carriers.
        \item Controlled analysis of confounders such as device tier, price, and chipset vendor.
        \item Measuring end-user receipt of updates; the paper focuses on OEM behavior and leaves users out of scope.
        \item Project Mainline (Google Play system updates) is described as a delivery path but is not measured (observation).
    \end{itemize}

\end{enumerate}

\newpage

\section{Paper 3: Trust, But Verify: A Longitudinal Analysis Of Android OEM Compliance and Customization}

\begin{enumerate}

    \item \entrypart{Paper Details}

    \begin{tabular}{@{}>{\bfseries}p{3.2cm} p{11.5cm}@{}}
        Authors          & Andrea Possemato, Simone Aonzo, Davide Balzarotti, Yanick Fratantonio \\[2pt]
        Publication Year & 2021 (IEEE Symposium on Security and Privacy (SP), 2021) \\[2pt]
        DOI/Link         & \url{https://doi.org/10.1109/SP40001.2021.00074}
    \end{tabular}

    \item \entrypart{Research Problem}

    \begin{itemize}[leftmargin=*,itemsep=1pt,topsep=2pt]
        \item \textbf{Problem:} OEM customizations of AOSP can hurt compatibility and security. Unclear whether ``Android''-branded ROMs actually follow Google's requirements.
        \item \textbf{Why it matters:}
        \begin{itemize}[itemsep=0pt]
            \item More than two billion devices run Android.
            \item Customizations can weaken Google's hardening, increase attack surface, and introduce new vulnerabilities.
            \item Customizations of core components can cause compatibility problems and delays in applying monthly security patches.
            \item Project Zero reported critical bugs in such customizations (e.g., Samsung Skia/Qmage).
        \end{itemize}
        \item \textbf{Google's countermeasures (as described):}
        \begin{itemize}[itemsep=0pt]
            \item Compliance: CDD (natural-language requirements: must / strongly recommended / recommended), CTS (covers only a subset of CDD), VTS (with Project Treble), GTS.
            \item Project Treble (Android 8.0): separate vendor code (/vendor partition) from the core framework to ease patching.
        \end{itemize}
        \item \textbf{Challenges listed by the paper:}
        \begin{itemize}[itemsep=0pt]
            \item Some CDD requirements are hard to express programmatically.
            \item Vendors control the build system and can strip identifying build information.
            \item Kernels are shipped as binaries only, which limits checks.
            \item Missing complete picture of OEM customizations over time across several parts of the OS.
        \end{itemize}
        \item \textbf{Questions asked (stated in Intro):}
        \begin{itemize}[itemsep=0pt]
            \item Are Google's automated compliance checks enough to detect CDD and VTS violations?
            \item Do certified ROMs violate requirements?
            \item Do customizations follow Treble's principle?
            \item Which customizations are most common, and do they hurt the security posture?
            \item How do vendor binaries' hardening compare with AOSP?
        \end{itemize}
    \end{itemize}

    \item \entrypart{Proposed Approach}

    \begin{itemize}[leftmargin=*,itemsep=1pt,topsep=2pt]
        \item \textbf{Main idea:} Fully automated pipeline that unpacks a ROM, finds its AOSP base tag, builds the AOSP reference image, and does a differential analysis of vendor customizations. About 20 minutes per ROM.
        \item \textbf{Stage 1: Unpacking}
        \begin{itemize}[itemsep=0pt]
            \item Auto-detects compression scheme and filesystem type.
        \end{itemize}
        \item \textbf{Stage 2: AOSP tag identification}
        \begin{itemize}[itemsep=0pt]
            \item Build ID in \texttt{ro.build.id}, mapped to AOSP tags using Google's documentation.
            \item Fallbacks: \texttt{ro.build.description}; \texttt{ro.com.google.gmsversion} with \texttt{ro.build.version.sdk}.
            \item Last resort: \texttt{ro.build.version.sdk} plus \texttt{ro.build.date.utc}; pick the tag with nearest creation timestamp.
            \item Clones and compiles the matched AOSP tag as reference.
        \end{itemize}
        \item \textbf{Stage 3: Binary customization}
        \begin{itemize}[itemsep=0pt]
            \item Checks security hardening of binaries (ELF executables, APK native libraries), new functionality, and security posture of new vs.\ AOSP binaries.
            \item Hardening inferred from ELF artifacts: stack canaries, NX, PIE, Full RELRO, Fortify Source, setuid/setgid (Table V); CFI/IntSan.
            \item New functions: exported symbols in shared binaries that appear in no AOSP release.
        \end{itemize}
        \item \textbf{Stage 4: SELinux policy}
        \begin{itemize}[itemsep=0pt]
            \item Extracts vendor-only rules from compiled policies vs.\ AOSP base.
            \item Permissive domains read from compiled policy.
            \item \texttt{neverallow} check: add each vendor rule to the original AOSP policy, recompile, and see if the compile error comes from a \texttt{neverallow}.
        \end{itemize}
        \item \textbf{Stage 5: Init scripts}
        \begin{itemize}[itemsep=0pt]
            \item Finds new services vs.\ AOSP base; focuses on services that start at boot and run as root.
            \item Checks user/group, capabilities, SELinux context.
        \end{itemize}
        \item \textbf{Stage 6: Kernel}
        \begin{itemize}[itemsep=0pt]
            \item Extracts version and banner, then build config.
            \item If \texttt{CONFIG\_IKCONFIG} is on, reads \texttt{.config}.
            \item Otherwise, extracts symbols (modified \texttt{vmlinux-to-elf}) and maps ELF symbols to kernel config options (Table VI).
            \item Marks a kernel non-compliant only if none of the interchangeable options is present.
        \end{itemize}
        \item \textbf{Stage 7: Compliance knowledge base}
        \begin{itemize}[itemsep=0pt]
            \item System-hardening requirements manually extracted from CDD of Android 1.6 to 9 (Table I).
        \end{itemize}
    \end{itemize}

    \item \entrypart{Dataset and Evaluation}

    \begin{itemize}[leftmargin=*,itemsep=1pt,topsep=2pt]
        \item \textbf{Data sources:} Google official site for Google ROMs; firmwarefile.com and stockrom.net for other vendors.
        \item \textbf{Dataset size:} 2,907 ROMs, 42 vendors, 1,403 device models; Android 2.3.3 to 9 (SDK 10 to 28); 2010 to 2020.
        \item \textbf{Composition:} Half from ``big players'' (e.g., Samsung, Huawei, LG, Xiaomi); rest from vendors with market share under 4\%. All ROMs contain Google apps, so the authors assume GMS certification.
        \item \textbf{Reference images:} 326 ROMs compiled from AOSP tags (not counted in statistics).
        \item \textbf{Tag identification source:} \texttt{ro.build.id} for 88\% (2,566); build date + SDK for 9\% (261); \texttt{gmsversion} for 2\% (59); \texttt{ro.build.description} for 23 ROMs.
        \item \textbf{Analysis subsets:}
        \begin{itemize}[itemsep=0pt]
            \item Kernel: 262 ROMs had no kernel binary; 2,645 remaining; 249 had no config or symbols; 2,396 analyzed (561 with textual config, rest via symbols).
            \item SELinux: 1,090 ROMs had no policy; 1,817 with policy; 1,533 target Android $\geq$ 5.
            \item CFI/IntSan: 196 ROMs with SDK $\geq$ 28.
        \end{itemize}
        \item \textbf{Ground truth:} Not Specified (no labeled set; correctness of tag identification and symbol-to-config mapping not quantified).
        \item \textbf{Train/test splits:} Not Specified (measurement study).
        \item \textbf{Metrics:} Share of ROMs violating CDD ``must'' and ``strongly recommended'' rules; counts of permissive domains and \texttt{neverallow} violations; share of binaries with each mitigation; new exported functions; \texttt{setuid} files; new init services.
        \item \textbf{Baseline comparison:} Not Specified. (Vendor binaries and policies are compared with their own AOSP base image as the reference. No comparison with another tool.)
    \end{itemize}

    \item \entrypart{Key Findings}

    \begin{itemize}[leftmargin=*,itemsep=1pt,topsep=2pt]
        \item \textbf{Headline compliance:}
        \begin{itemize}[itemsep=0pt]
            \item 579 of 2,907 ($\sim$20\%) ROMs violate at least one CDD ``must'' rule; 11 of them are branded by Google itself.
            \item 289 ($\sim$10\%) do not implement at least one ``strongly recommended'' suggestion.
        \end{itemize}
        \item \textbf{Kernel:}
        \begin{itemize}[itemsep=0pt]
            \item 162 kernels for Android $\geq$ 8.0 lack \texttt{CONFIG\_IKCONFIG} and \texttt{CONFIG\_MODULES}; authors say these should fail VTS.
            \item 190 of 2,396 (7.9\%) kernels (10 vendors) violate a CDD requirement; 162 of them in Treble ROMs.
            \item Most common: absence of kernel memory protections (\texttt{CONFIG\_DEBUG\_RODATA} or \texttt{CONFIG\_STRICT\_KERNEL\_RWX}), on 150 kernels.
            \item 241 of 2,396 (10\%) lack a strongly recommended feature; 160 did not enable \texttt{CONFIG\_RANDOMIZE\_BASE}. (Table II reports 237 (9.9\%).)
            \item Table II violation rates: SDK 26: 34.5\%; SDK 27: 23.6\%; SDK 28: 39.8\%.
        \end{itemize}
        \item \textbf{SELinux:}
        \begin{itemize}[itemsep=0pt]
            \item Permissive domains: 108 of 1,533 (7.0\%) ROMs (16 vendors) violate CDD (Table III).
            \item \texttt{neverallow} violations: 443 of 1,533 (28.9\%) ROMs, 21 vendors (Table IV).
            \item Peaks of 85.8\% (SDK 26), 79.1\% (SDK 27), 62.2\% (SDK 28).
            \item Authors say the Treble-era problem got worse, not better.
        \end{itemize}
        \item \textbf{Binary hardening (CFI/IntSan, SDK $\geq$ 28, 196 ROMs):}
        \begin{itemize}[itemsep=0pt]
            \item 85 (43.37\%) disable CFI on at least one binary; 104 (53.06\%) disable IntSan.
            \item Six vendors lowered a binary's security vs.\ AOSP.
            \item Among violators, CFI disabled for 38.7\% ($\sigma$=36.5) and IntSan for 35.8\% ($\sigma$=34.9) of binaries.
        \end{itemize}
        \item \textbf{New functions in system libraries:}
        \begin{itemize}[itemsep=0pt]
            \item Almost constant trend: roughly 80 new functions added to 20\% of system libraries.
            \item Vendors also still use functions found only in older AOSP releases.
        \end{itemize}
        \item \textbf{Compile-time hardening (vendor-only binaries):}
        \begin{itemize}[itemsep=0pt]
            \item Binaries shared with AOSP keep AOSP's mitigations.
            \item Vendor-only binaries consistently use fewer mitigations than AOSP binaries (Fig.~3).
            \item Around 40\% of vendor binaries lack stack canaries.
            \item \texttt{setuid}: 319/447 (71\%) of ROMs with SDK $<$ 18 and 371/2453 (15\%) of ROMs with SDK $\geq$ 18 contain at least one. Most frequent: \texttt{su} (18\%), \texttt{procmem} (17\%), \texttt{netcfg} (16\%), \texttt{procrank} (12\%), \texttt{tcpdump} (11\%).
            \item CDD never mentions NX, but CTS tests it (inconsistency noted).
        \end{itemize}
        \item \textbf{Init scripts:}
        \begin{itemize}[itemsep=0pt]
            \item Number of new services grows over time; some ROMs define almost 250 additional services.
            \item SDK 26 example: AOSP averages 59 services, average ROM 90; peak vendor adds 195.
            \item Vendors customize AOSP services by adding root as owner.
            \item \textbf{Case:} 18 ROMs (2 vendors) start \texttt{tcpdump} at boot as root, saving the first 134 bytes of each packet on all interfaces; some use \texttt{tcpdump} 4.9.2 (2017, several CVEs, some with public PoCs); also found in an Android One ROM built in 2019.
        \end{itemize}
        \item \textbf{SELinux customization:}
        \begin{itemize}[itemsep=0pt]
            \item SDK 27: AOSP policy averages 10,000 rules, some vendors over 232,000 (over 20x).
            \item More than 130,000 permission changes to AOSP core policy in SDK 27.
            \item \texttt{isolated\_app}: 58,776 changes (1,375 unique). \texttt{untrusted\_app}: 95,577 changes (4,228 unique).
            \item Examples: SDK 19, 44 ROMs let \texttt{isolated\_app} read/write/ioctl kernel drivers; SDK 23 and 25 rule lets \texttt{isolated\_app} open, read, write app data files (breaks a \texttt{neverallow}); SDK 27 rule lets \texttt{untrusted\_app} read files with the device MAC address.
            \item Motorola: one policy change reverted the \texttt{vold} mitigation, allowing root (cited case).
        \end{itemize}
        \item \textbf{Authors' recommendations:} more automated checks, more accurate existing checks, extend the CDD to discourage risky customizations.
    \end{itemize}

    \item \entrypart{Contribution and Limitations}

    \textbf{Contributions (authors)}
    \begin{itemize}[leftmargin=*,itemsep=1pt,topsep=2pt]
        \item First longitudinal, large-scale analysis of 2,907 Android ROMs from 42 OEMs over 10 Android versions.
        \item In-depth look at two aspects: compliance and security posture.
        \item Found certified ROMs that do not meet security requirements, despite supposedly passing the test suites.
        \item Showed vendor-specific components lag behind AOSP in security posture, and described techniques that weaken AOSP safety nets without strictly violating guidelines.
        \item Automated analysis pipeline; authors state they will release detailed results and the pipeline.
        \item Recommendations for Google on compliance checks and guidelines.
    \end{itemize}

    \vspace{4pt}
    \textbf{Limitations (authors)}
    \begin{itemize}[leftmargin=*,itemsep=1pt,topsep=2pt]
        \item Tag identification may have errors; authors say the analysis is not sensitive to small imprecisions.
        \item Kernel config inference from symbols does not support all kernel config options.
        \item 3 kernel configurations have no symbol mapping (Table VI).
        \item Some ROMs lacked a kernel binary or \texttt{boot.img}, or no config or symbols could be extracted; no measurement possible for them.
        \item 351 ROMs with SELinux kernel support but no policy; authors suppose they are incremental updates.
        \item Stack canary measurement assumes that each binary has at least one protectable buffer.
        \item Violations do not imply maliciousness; authors believe most vendors act in good faith.
    \end{itemize}

    \item \entrypart{My Analysis}

    \textbf{Good}
    \begin{itemize}[leftmargin=*,itemsep=1pt,topsep=2pt]
        \item Large, long-running dataset (2010 to 2020) across many vendors.
        \item Differential analysis against the exact AOSP base is a sound idea; separates AOSP issues from vendor issues.
        \item Covers several layers at once: kernel, SELinux, binaries, init scripts.
        \item Splits results before and after Treble, giving a direct test of Treble's effect.
        \item Concrete, checkable examples (tcpdump at boot, commented-out \texttt{neverallow}, Skia).
        \item Careful wording: conservative kernel-flag rule, and a statement that violations are not proof of malice.
    \end{itemize}

    \vspace{4pt}
    \textbf{Weak points}
    \begin{itemize}[leftmargin=*,itemsep=1pt,topsep=2pt]
        \item \textbf{Validation:}
        \begin{itemize}[itemsep=0pt]
            \item Observation: no accuracy figure for tag identification, symbol-to-config mapping, or \texttt{neverallow} detection; no manual validation sample reported.
            \item Observation: \texttt{neverallow} detection compiles each added vendor rule against AOSP policy, so it flags rule-level conflicts, not a check that the vendor actually removed the rule (the ``commented out'' claim is an inference).
            \item Observation: results come from static analysis of ROMs, not from real devices or exploits.
        \end{itemize}
        \item \textbf{Certification assumption:} All ROMs are assumed GMS certified because they contain Google apps. Observation: this does not prove they passed CTS/VTS for that specific build, so claims about the test suites failing are inferences. The ROMs also come from third-party firmware sites, so provenance is unclear.
        \item \textbf{Dataset bias:}
        \begin{itemize}[itemsep=0pt]
            \item Observation: ROMs come from two third-party sites, so the sample may not represent what ships on devices; per-vendor counts are not given in the text.
            \item Half the data is from big vendors; vendor-level conclusions for small vendors may rest on few ROMs.
            \item Observation: SDK distribution is uneven (Fig.~1), and rates per SDK rest on different sample sizes.
        \end{itemize}
        \item \textbf{Overlap of design and evaluation data:} The same dataset is used to build the heuristics and to measure results. Observation: no held-out set.
        \item \textbf{Coverage:}
        \begin{itemize}[itemsep=0pt]
            \item Large share of ROMs excluded in parts (262 without kernel; 1,090 without SELinux policy).
            \item Observation: percentages are over the analyzed subset, not all 2,907.
            \item Headline ``$\sim$20\% of ROMs'' is over all 2,907, while the component results use smaller denominators (2,396 and 1,533). How these combine to 579 is not broken down in the text.
        \end{itemize}
        \item \textbf{Number inconsistencies (observation):}
        \begin{itemize}[itemsep=0pt]
            \item Strongly recommended kernel violations: 241 (10\%) in text vs.\ 237 (9.9\%) in Table II.
            \item ``160 vendors did not enable \texttt{CONFIG\_RANDOMIZE\_BASE}'' appears to mean 160 kernels.
            \item Dataset described as Android 1.6 to 9.0 (years 2009 to 2020) in the abstract, 2.3 to 9.0 (2010 to 2020) in the Introduction.
        \end{itemize}
        \item \textbf{Heuristics:} Tag identification fallback (nearest build date) and symbol-based kernel config inference are heuristic.
        \item \textbf{Impact claims:} Observation: no evidence is shown that most flagged customizations are exploitable; tcpdump case is shown for 18 ROMs only.
        \item \textbf{Baselines:} No comparison with other analysis tools or with CTS/VTS run on the same ROMs.
        \item \textbf{Age:} Data ends at Android 9 / 2020; Observation: findings may not hold for later Android versions.
    \end{itemize}

    \vspace{4pt}
    \textbf{Research gaps}
    \begin{itemize}[leftmargin=*,itemsep=1pt,topsep=2pt]
        \item Validate the pipeline with a ground-truth set (known tags, known kernel configs, known policy changes).
        \item Run CTS/VTS on the same ROMs to measure which violations the suites actually miss.
        \item Study Android versions after 9, and whether newer Treble-era requirements changed the trend.
        \item Link customizations to real vulnerabilities or exploits instead of only to posture metrics.
        \item Find the causes of vendor behavior (the paper only suggests practicality reasons).
        \item Automate more CDD requirements (authors' own recommendation).
        \item Analyze non-security-hardening CDD parts and other components (e.g., framework, drivers) not covered here.
        \item Combine with update-timeliness data to test whether heavier customization causes patch delays (paper states the link but does not measure it).
    \end{itemize}

\end{enumerate}

\newpage

\section{Paper 4: Trouble Over-The-Air: An Analysis of FOTA Apps in the Android Ecosystem}

\begin{enumerate}

    \item \entrypart{Paper Details}

    \begin{tabular}{@{}>{\bfseries}p{3.2cm} p{11.5cm}@{}}
        Authors          & Eduardo Bl\'{a}zquez, Sergio Pastrana, \'{A}lvaro Feal, Julien Gamba, Platon Kotzias, Narseo Vallina-Rodriguez, Juan Tapiador \\[2pt]
        Publication Year & 2021 (IEEE Symposium on Security and Privacy (SP), 2021) \\[2pt]
        DOI/Link         & \url{https://doi.org/10.1109/SP40001.2021.00095}
    \end{tabular}

    \item \entrypart{Research Problem}

    \begin{itemize}[leftmargin=*,itemsep=1pt,topsep=2pt]
        \item \textbf{Problem:} Android firmware updates are handled by FOTA (Firmware Over-the-Air) apps. These are highly privileged, but the FOTA ecosystem has been largely overlooked by researchers.
        \item \textbf{Why it matters:}
        \begin{itemize}[itemsep=0pt]
            \item Over 2.5 billion active Android devices.
            \item FOTA apps are privileged; an attacker or deceptive FOTA provider could install harmful system-level components over a large user base.
            \item Anecdotal case: Adups, a Chinese update provider, reported as spyware.
            \item Some FOTAs deliver updates to tens of millions of devices.
        \end{itemize}
        \item \textbf{Challenges listed by the paper:}
        \begin{itemize}[itemsep=0pt]
            \item No clear set of rules a FOTA app must follow; Google's recommendations are described as superficial and vague on security practices.
            \item FOTA apps may lack distinctive package names, split functionality across apps, or use custom code.
            \item Updates are often delegated to third- and fourth-party services; stakeholders are hard to attribute (no PKI to verify signing certificates).
            \item Pre-installed apps are not on app stores, so no metadata or privacy policies; dynamic analysis is hard (platform-key signing, hardware dependencies).
        \end{itemize}
        \item \textbf{Scope:} Providers using standard AOSP capabilities to implement FOTA components. Projects Treble and Mainline analysis out of scope (lack of representative data).
    \end{itemize}

    \item \entrypart{Proposed Approach}

    \begin{itemize}[leftmargin=*,itemsep=1pt,topsep=2pt]
        \item \textbf{Main idea:} Large-scale static analysis of pre-installed FOTA apps, plus telemetry from real devices, to study the FOTA supply chain, privacy and security risks, and installation behavior.
        \item \textbf{Stage 1: FOTA Finder (detection)}
        \begin{itemize}[itemsep=0pt]
            \item Built on Androguard; fingerprints from Google's FOTA documentation and manual analysis of 18 well-known FOTA apps.
            \item 4 strong signals (Table I): \texttt{RecoverySystem.verifyPackage} (RV), \texttt{RecoverySystem.installPackage} (RI), strings ``/cache/recovery/command'' and ``--update-package'' (CMD), \texttt{UpdateEngine.applyPayload} (A/B). Any one signal classifies the app as FOTA.
            \item 11 secondary signals for OTA (non-system app install) capabilities (Appendix B).
            \item \textbf{Dextripador:} helper tool the authors reverse-engineered to recover DEX from ODEX files.
        \end{itemize}
        \item \textbf{Stage 2: Ecosystem analysis}
        \begin{itemize}[itemsep=0pt]
            \item Stakeholders inferred from signing certificate (Organization field, email domain) and package name; assumes the signer deploys the app.
            \item Categories via semi-manual snowball sampling and web searches: OEM, MNO, SoC, SFD (Specialized FOTA Developer).
            \item Prevalence by brand and by number of FOTAs per device.
        \end{itemize}
        \item \textbf{Stage 3: Behavior analysis (static)}
        \begin{itemize}[itemsep=0pt]
            \item FlowDroid and Amandroid (taint analysis), modified Androwarn (API usage).
            \item LibRadar (enhanced) for third-party SDKs, only libraries at 100\% confidence.
            \item Cartographer to attribute API calls to app code vs.\ third-party SDK; permission mapping with Axplorer plus custom mappings.
            \item Manifest analysis: shared UIDs, AOSP and custom permissions.
        \end{itemize}
        \item \textbf{Stage 4: Update-process analysis}
        \begin{itemize}[itemsep=0pt]
            \item Checks for install/verify capability using FOTA Finder signals; looks for \texttt{MessageDigest} as sign of custom verification.
            \item Inspects \texttt{otacerts.zip} files from devices.
        \end{itemize}
        \item \textbf{Stage 5: Telemetry analysis (NortonLifeLock)}
        \begin{itemize}[itemsep=0pt]
            \item Matches FOTA package names (and signer certificates) to installer (parent) package in reputation logs.
            \item VirusTotal: malicious if flagged by at least 4 AV engines; AVClass for malware families; Google Play lookup by package name.
        \end{itemize}
    \end{itemize}

    \item \entrypart{Dataset and Evaluation}

    \begin{itemize}[leftmargin=*,itemsep=1pt,topsep=2pt]
        \item \textbf{Dataset 1: Firmware Scanner (Gamba et al.)}
        \begin{itemize}[itemsep=0pt]
            \item 422,121 pre-installed apps, 40,165 users, 184 countries, 12,539 devices (June 2020).
            \item 152,097 apps found only on rooted devices discarded.
            \item 3,311 \texttt{otacerts.zip} files from 1,747 devices.
        \end{itemize}
        \item \textbf{Dataset 2: NortonLifeLock reputation and installation logs}
        \begin{itemize}[itemsep=0pt]
            \item 1.6 B installation events from 19.3 M Android devices, January to December 2019; opt-in customers; metadata only, no APKs.
            \item Limited to installs via PackageInstaller.
            \item FOTA-related subset: 15,961,424 installation events, 20 FOTA packages, 841,344 devices, 199 country codes.
        \end{itemize}
        \item \textbf{Ground truth:} No labeled ground-truth set. FOTA Finder checked for false positives only.
        \item \textbf{Evaluation of FOTA Finder:}
        \begin{itemize}[itemsep=0pt]
            \item Manual review of 50 randomly chosen FOTA apps by one author: 3 potential FPs, judged to be parts of larger multi-app FOTA systems, so not counted as FPs. ``No false positives'' reported.
            \item Google Play check (24 June 2020): 7 FOTA packages found on the store, including Google Play Services; the other six verified as genuine vendor FOTA apps.
            \item False negatives: accepted; not measured.
        \end{itemize}
        \item \textbf{Metrics:} Counts and percentages of FOTA apps, devices, packages; per-category shares; per-behavior and per-permission percentages; installation events, devices, malicious APKs per installer.
        \item \textbf{Train/test splits:} Not Specified (measurement study).
        \item \textbf{Baseline comparison:} Not Specified. (No comparison against another FOTA detector.)
    \end{itemize}

    \item \entrypart{Key Findings}

    \begin{itemize}[leftmargin=*,itemsep=1pt,topsep=2pt]
        \item \textbf{Detection:}
        \begin{itemize}[itemsep=0pt]
            \item 2,013 FOTA apps in 20,924 devices.
            \item FOTA Finder could not process 37\% (154,922) of apps (ODEX parsing failure).
            \item 32\% (647) have no token like update, install, fota in the package name.
            \item 93\% (1,878) use \texttt{RecoverySystem}; 7\% (135) use A/B updates.
            \item 32\% (651) can also install user-level apps.
        \end{itemize}
        \item \textbf{Ecosystem:}
        \begin{itemize}[itemsep=0pt]
            \item 269 unique certificates, 219 subjects, 127 organizations.
            \item Share of 2,013 FOTA apps: OEMs 53\% (77 packages); SoCs 9\% (13); SFDs 9\% (13); MNOs 1.6\% (4); Google 15\% (2 packages); Uncategorized 12\% (24).
            \item 43\% of FOTA apps are developed by a third party; 42\% (47 packages) are first-party OEM FOTAs.
            \item 49 packages (37\%) signed by two or more organizations; \texttt{com.adups.fota.sysoper} signed by 60 organizations; \texttt{com.mediatek.systemupdate.sysoper} by 33 certificates.
            \item 133 packages across 395 brands.
            \item 80\% of devices have Google FOTA components. Excluding those: 71\%, 20\%, 8\%, 0.8\%, 0.2\% of devices have one, two, three, four, five FOTAs. Some devices have as many as 5.
        \end{itemize}
        \item \textbf{Signing and verification issues:}
        \begin{itemize}[itemsep=0pt]
            \item 40 FOTA apps (2\%), 13 packages, devices from 20 brands, signed with AOSP test keys (at least 171 devices).
            \item 5\% of the 1,747 devices with \texttt{otacerts.zip} include an AOSP default key (17 brands; mostly Alps 57\%, Goclever 12\%); verification becomes useless. (Discussion mentions 90 such devices.)
            \item 8\% of non-A/B FOTAs use verification with no evidence of installation; 45\% have install flags without \texttt{verifyPackage}. Of 949 not using Google's verify API, 51\% call \texttt{MessageDigest}, mostly MD5.
        \end{itemize}
        \item \textbf{Behavior (1,716 non-Google FOTA apps, Table III):}
        \begin{itemize}[itemsep=0pt]
            \item IMEI 33.7\% (577); IMSI 31.4\% (538); installed apps 49.2\%; logs 65.3\%.
            \item GPS 0.7\% (12); cell location 4.3\%.
            \item Native code 27.1\%; Linux commands 30.9\%; phone calls 8.5\% (146); audio recording 2.6\% (44); video capture 2.3\% (40).
            \item Some apps send device and user data to update servers (Table IV), e.g., \texttt{com.redstone.ota.ui} sends many fields including IMEI and IMSI.
        \end{itemize}
        \item \textbf{Third-party SDKs:} 8 social, tracking, or ad SDKs in 171 (10\%) apps; Firebase 133, Google Ads 23, Umeng 8, SinaWeibo 6, Tencent 4.
        \item \textbf{Privileges:}
        \begin{itemize}[itemsep=0pt]
            \item Shared UID in 1,737 FOTA apps (80 packages); 60\% use \texttt{android.uid.system}.
            \item Median 25 AOSP permissions and 6 custom permissions per app.
            \item READ\_PHONE\_STATE 70\%; ACCESS\_COARSE\_LOCATION 21\%; SEND\_SMS 16\%; ACCESS\_FINE\_LOCATION 15\%.
            \item Dangerous permissions mostly used by the app itself; tracking SDKs rarely access them.
        \end{itemize}
        \item \textbf{Telemetry:}
        \begin{itemize}[itemsep=0pt]
            \item 101 (63\%) FOTA packages seen in telemetry; only 20 (12\%) show installation activity; results are lower bounds.
            \item 13 of 20 installers install 764 packages that exist on Google Play (50 categories; Tools 20\%, Entertainment 9\%, Communication 9\%, Games 8\%, Shopping 7\%).
            \item 7 of 20 FOTA packages install at least one unwanted app; five have 0.1\% to 2.8\% unwanted.
            \item \texttt{com.android.settings}: 494 malicious APKs (35\%); families include triada and necro.
            \item \texttt{com.qiku.android.ota}: 11 of 12 APKs malicious (92\%); triada, necro, guerilla.
            \item \texttt{com.sonyericsson.updatecenter}: 12.6M events, 381K devices; \texttt{com.samsung.android.app.omcagent}: 3.0M events, 332K devices, 29 malicious APKs (1.5\%).
        \end{itemize}
        \item \textbf{Samsung native code:} 149 devices with FOTA app plus \texttt{libmno\_dmstack.so}; manual analysis shows update by writing to \texttt{/cache/recovery/command}.
    \end{itemize}

    \item \entrypart{Contribution and Limitations}

    \textbf{Contributions (authors)}
    \begin{itemize}[leftmargin=*,itemsep=1pt,topsep=2pt]
        \item First large-scale, systematic study of Android FOTA apps (2,013 apps).
        \item FOTA Finder (detection tool) and Dextripador (ODEX to DEX recovery).
        \item Classification of FOTA stakeholders (OEM, MNO, SoC, SFD) and supply-chain picture.
        \item Static analysis of privacy and security behaviors: sensitive data collection, third-party SDKs, shared UIDs, permissions, weak or missing verification, AOSP test keys.
        \item Telemetry evidence that FOTA apps install non-system apps, including PUPs and malware.
        \item Findings reported to Google and main vendors; tools and aggregated dataset released.
    \end{itemize}

    \vspace{4pt}
    \textbf{Limitations (authors)}
    \begin{itemize}[leftmargin=*,itemsep=1pt,topsep=2pt]
        \item FOTA Finder only finds apps using FOTA-specific APIs or known strings; misses custom implementations, native code.
        \item ODEX parsing fails for 37\% of apps; larger coverage might change the distribution.
        \item Static analysis only: dynamic code loading, reflection, native code, string obfuscation may hide behaviors or cause missed FOTAs.
        \item Findings show potential behaviors, not runtime observations.
        \item Stakeholder classification relies on certificates, package names, and manual web searches; possible mislabels (e.g., an OEM labeled as SoC).
        \item Telemetry limited to installs through PackageInstaller; covers a subset of FOTA installs (lower bound). Package-name matching could be spoofed (partly addressed using certificates).
        \item Dynamic analysis not possible (platform-key signing, device dependencies, avoid production servers); other update steps (e.g., server exchanges) not analyzed.
        \item Low \texttt{otacerts.zip} coverage (Firmware Scanner added the feature only in November 2019; FOTA Finder limits).
        \item Treble and Mainline not analyzed.
        \item Why six FOTA apps were accepted on Google Play is unexplained.
    \end{itemize}

    \item \entrypart{My Analysis}

    \textbf{Good}
    \begin{itemize}[leftmargin=*,itemsep=1pt,topsep=2pt]
        \item Studies an overlooked, high-privilege component at large scale.
        \item Combines static analysis with real-device telemetry to partly cover the dynamic-analysis gap.
        \item Clear, checkable security findings (AOSP test keys in signing and in \texttt{otacerts.zip}).
        \item Open about FOTA Finder's scope and focus on precision over recall.
        \item Tools and aggregated dataset released.
    \end{itemize}

    \vspace{4pt}
    \textbf{Weak points}
    \begin{itemize}[leftmargin=*,itemsep=1pt,topsep=2pt]
        \item \textbf{Validation:}
        \begin{itemize}[itemsep=0pt]
            \item Observation: FP check is 50 apps judged by one author; no recall estimate; ``no false positives'' rests on this small sample and on reinterpreting three borderline cases as not FPs.
            \item Observation: 37\% of the dataset could not be processed, so findings come from the parsed 63\%; a bias toward ODEX-parsable apps is possible.
        \end{itemize}
        \item \textbf{Overlap of design and evaluation data:} Signals were derived from 18 apps found by searching the same dataset for names like update or fota. Observation: name-based seeding may bias toward conventional FOTAs.
        \item \textbf{Dataset bias:}
        \begin{itemize}[itemsep=0pt]
            \item Observation: crowd-sourced volunteer users with a Firmware Scanner app; may not match the global device population. Rooted-device apps discarded.
            \item Telemetry comes from one security vendor's customers, and only 20 packages show installs; two packages dominate events (Sony 12.6M, Samsung 3.0M of 15.96M).
        \end{itemize}
        \item \textbf{Heuristics and manual steps:}
        \begin{itemize}[itemsep=0pt]
            \item Attribution by certificate and package name, with manual web searches; Uncategorized 12\%.
            \item Manual code review for secondary-purpose permission use.
            \item Custom verification judged by a \texttt{MessageDigest} call only; it does not show verification is actually correct (observation).
        \end{itemize}
        \item \textbf{Static analysis only:} Table III and permission results show capability or potential, not actual data collection. The text itself notes uploads are shown for only selected apps, and the purpose is unknown.
        \item \textbf{Malware claims:}
        \begin{itemize}[itemsep=0pt]
            \item Observation: threshold of at least 4 AV engines is aligned with prior work but may include PUPs; \texttt{com.android.settings} is a generic package capturing multiple vendors, so it is not one actor.
            \item Observation: the Qiku figure (11 of 12 APKs) is a tiny sample; the abstract's ``92\%'' headline rests on it.
            \item Observation: the paper says installs ``possibly'' stem from compromised vendor code; the cause is not shown.
        \end{itemize}
        \item \textbf{Number inconsistencies (observation):}
        \begin{itemize}[itemsep=0pt]
            \item AOSP-key devices: 5\% of 1,747 devices (Sec.~VI) vs.\ 90 devices (Sec.~VIII).
            \item \texttt{otacerts.zip}: 598 (6\%) devices said to contain a FOTA app; denominator unclear.
            \item Table III: SIM serial number third-party count 446 exceeds all-apps count 111.
            \item ``Google FOTA'' share (15\%) comes from \texttt{gms} and \texttt{gsf}, which the paper says do not perform system updates; treating them as FOTA inflates counts (observation).
            \item The text says Google packages are 15\% and 80\% of devices have Google FOTA components, so Google-package results mix with OEM results in places.
        \end{itemize}
        \item \textbf{Single-case reporting:} Per-app privacy examples (e.g., Redstone, Motorola) are single apps.
        \item \textbf{Baselines:} No comparison with other detectors or earlier pre-installed-app studies beyond textual references.
        \item \textbf{Age:} Data from 2019 to June 2020; Observation: Android versions, A/B adoption, and Mainline coverage have likely changed.
    \end{itemize}

    \vspace{4pt}
    \textbf{Research gaps}
    \begin{itemize}[leftmargin=*,itemsep=1pt,topsep=2pt]
        \item Detecting FOTAs that use native code, reflection, custom services, or obfuscation (authors leave native/Linux-command behavior for future work).
        \item Dynamic or hybrid analysis of FOTA apps and update-server traffic.
        \item Measuring real exploitability of AOSP test-key signing and \texttt{otacerts.zip} issues on devices.
        \item Understanding why vendors install non-system apps via FOTA (commercial agreements vs.\ shared components); telemetry beyond PackageInstaller.
        \item Impact of Project Treble and Project Mainline on FOTA use (authors leave to future work).
        \item Stronger attribution of FOTA developers and operators across the supply chain.
        \item Link FOTA behavior to update timeliness and support data (e.g., whether third-party FOTAs delay patches); the paper does not measure this (observation).
        \item Better, enforceable Google guidance and checks for FOTA implementation and verification.
    \end{itemize}

\end{enumerate}

\newpage

\section{Paper 5: The Android Platform Security Model}

\begin{enumerate}

    \item \entrypart{Paper Details}

    \begin{tabular}{@{}>{\bfseries}p{3.2cm} p{11.5cm}@{}}
        Authors          & Ren\'{e} Mayrhofer, Jeffrey Vander Stoep, Chad Brubaker, Nick Kralevich \\[2pt]
        Publication Year & 2021 (ACM Transactions on Privacy and Security, Vol.\ 24, No.\ 3, Article 19, April 2021; received May 2020, revised January 2021, accepted January 2021; last updated December 2020 based on Android 11) \\[2pt]
        DOI/Link         & \url{https://doi.org/10.1145/3448609}
    \end{tabular}

    \item \entrypart{Research Problem}

    \begin{itemize}[leftmargin=*,itemsep=1pt,topsep=2pt]
        \item \textbf{Problem:} The Android security model was never formally published. Its design principles only informed architecture, access control, and mitigations implicitly.
        \item \textbf{Why it matters:}
        \begin{itemize}[itemsep=0pt]
            \item More than 2.5 billion monthly active devices; many form factors (phones, tablets, wearables, TV, IoT, cars).
            \item Use cases increasingly security and privacy critical (communication, finance, health, sensors).
            \item Needs to balance security, privacy, and usability for end users, assurances for app developers, and performance under tight hardware constraints.
        \end{itemize}
        \item \textbf{Challenges listed by the paper:}
        \begin{itemize}[itemsep=0pt]
            \item Android is end-user focused; users are not security experts, so the OS must not offload hard decisions to them.
            \item Huge ecosystem: hundreds of OEMs, tens of thousands of devices; no single register of OEMs; API changes take time to reach all devices.
            \item Apps can be written in any language and installed from arbitrary sources, so security cannot rely on compile-time checks; it must rely on runtime protections at the app boundary.
            \item Devices are easily lost or stolen and use untrusted networks.
        \end{itemize}
        \item \textbf{Aims (stated contributions):}
        \begin{itemize}[itemsep=0pt]
            \item Define the model from security principles and context.
            \item Define the threat model and special cases.
            \item Explain how AOSP enforces the model through interacting measures on several layers.
            \item Identify open gaps and future improvements.
        \end{itemize}
        \item \textbf{Scope:} Platform = AOSP components forming a system that passes CTS. GMS, Google Play Protect, Play policy are out of scope.
    \end{itemize}

    \item \entrypart{Proposed Approach}

    \begin{itemize}[leftmargin=*,itemsep=1pt,topsep=2pt]
        \item \textbf{Main idea:} Informal (not formal-notation) security model based on multi-party consent, plus a description of how AOSP implements it with defense in depth and safe-by-default design. Implementation sections refer to Android 11 at initial release.
        \item \textbf{Threat model (Sec.~2.2)}
        \begin{itemize}[itemsep=0pt]
            \item Physical/proximal: T.P1 to T.P4 (e.g., powered-off device under adversary control, locked device, shared unlocked device).
            \item Network: T.N1 (passive eavesdropping, tracking), T.N2 (active manipulation, e.g., on-path TLS attacks).
            \item Untrusted code: T.A1 to T.A7 (API abuse, UI mimicking, input injection, OS bugs).
            \item Untrusted content: T.D1 (unique identifiers), T.D2 (code processing untrusted content).
        \end{itemize}
        \item \textbf{Five model rules (Sec.~3)}
        \begin{itemize}[itemsep=0pt]
            \item \textbf{1 Multi-party consent:} no action unless user, platform, and developer all agree; any party can veto.
            \item \textbf{2 Open ecosystem access:} no central vetting; app-to-app interaction supported.
            \item \textbf{3 Security is a compatibility requirement:} model is part of the CDD, enforced by CTS and other test suites.
            \item \textbf{4 Factory reset restores a safe state:} depends only on integrity-protected system code.
            \item \textbf{5 Applications are security principals:} apps are not fully authorized agents of the user.
        \end{itemize}
        \item \textbf{Implementation layers (Sec.~4)}
        \begin{itemize}[itemsep=0pt]
            \item \textbf{Meaningful consent (4.1):} developer consent via signing; platform consent via Verified Boot and component isolation; user consent principles (avoid over-prompting, understandable prompts, pickers, no offloading hard problems, undo).
            \item \textbf{Authentication (4.2):} tiered lockscreen (primary knowledge factor, secondary biometrics in Class 1/2/3, tertiary environmental); Gatekeeper/Weaver; FIDO2/WebAuthn; Identity Credential.
            \item \textbf{Isolation and containment (4.3):} access control (DAC, MAC/SELinux, Android permissions); UID application sandbox; system process sandboxing; kernel hardening; sandboxing below the kernel (Keymaster, Strongbox, Gatekeeper, Protected Confirmation).
            \item \textbf{Encryption at rest (4.4):} FDE, then FBE with Direct Boot; Adiantum.
            \item \textbf{Encryption in transit (4.5):} TLS by default, network security config, DNS-over-TLS, MAC randomization.
            \item \textbf{Exploit mitigation (4.6):} ASLR, W$\oplus$X, stack protector, sanitizers, CFI, Shadow Call Stack.
            \item \textbf{System integrity (4.7):} Verified Boot, dm-verity, rollback protection, key attestation, AVB VBMeta, key hierarchy, APK signing.
            \item \textbf{Patching (4.8):} monthly security bulletins since August 2015; Project Treble (Android 8.0); Project Mainline / APEX (Android 10).
        \end{itemize}
        \item \textbf{Method of analysis:} Tables 1 to 4 summarize security-relevant AOSP changes from Android 4.x to 11 (about 10 years of code evolution), each mapped to the threats it mitigates.
        \item \textbf{Special cases (Sec.~5):} package listing, VPN apps, backup, enterprise (a DPC as fourth party), Factory Reset Protection.
    \end{itemize}

    \item \entrypart{Dataset and Evaluation}

    \begin{itemize}[leftmargin=*,itemsep=1pt,topsep=2pt]
        \item \textbf{Data sources:} Analysis of security-relevant changes to the whole AOSP code base, Android 4.x to 11 (as stated). Some usage figures come from internal analysis not formally published.
        \item \textbf{Dataset sizes and splits:} Not Specified. This is a conceptual and descriptive paper.
        \item \textbf{Ground truth:} Not Specified.
        \item \textbf{Experimental setup:} Not Specified. No experiments are reported.
        \item \textbf{Metrics:} Not Specified (a few statistics are cited from other sources, see Key Findings).
        \item \textbf{Baseline comparison:} Not Specified. Related work is discussed in prose only (classical models such as Bell-LaPadula, access control lists, capability lists; other mobile OS surveys referenced).
    \end{itemize}

    \item \entrypart{Key Findings}

    \begin{itemize}[leftmargin=*,itemsep=1pt,topsep=2pt]
        \item \textbf{Nature of results:} No empirical evaluation. Results are the model, the threat analysis, and the mapping of mechanisms to threats.
        \item \textbf{Multi-party consent:}
        \begin{itemize}[itemsep=0pt]
            \item Default control of data depends on location: shared storage by users, private app directories by apps, special system locations by the platform.
            \item Example consent sets given for sharing data between apps and for changing MNO configuration.
            \item The platform is treated as multiple components that enforce each other's limits.
        \end{itemize}
        \item \textbf{Cited statistics:}
        \begin{itemize}[itemsep=0pt]
            \item As of August 2020, 77\% of devices with fingerprint sensors have a secure lockscreen enabled, vs.\ 54\% without fingerprints (internal, unpublished).
            \item Historically about 85\% of Android security vulnerabilities result from unsafe memory access.
            \item 59\% of app-reachable kernel vulnerabilities went through \texttt{ioctl()}; 9\% via debugfs; 11\% via \texttt{perf\_event\_open()} (Table 1).
            \item Software codecs: approximately 80\% of critical/high severity media vulnerabilities (Table 2). Missing or incorrect array bounds checks: 34\% of userspace vulnerabilities.
        \end{itemize}
        \item \textbf{Biometric tiers (Android 10):}
        \begin{itemize}[itemsep=0pt]
            \item Class 3: SAR $<$ 7\% and secure pipeline. Class 2: 7\% $<$ SAR $<$ 20\% and secure pipeline. Class 1: SAR $>$ 20\% or insecure pipeline.
            \item All classes: false acceptance rate at most 1/50,000; false rejection rate under 10\%.
            \item Fallback to primary authentication every 72 hours (Class 3) or 24 hours (Class 1 and 2); tertiary modalities after a 4-hour idle period.
        \end{itemize}
        \item \textbf{Key changes by release (Tables 1 to 4, examples):}
        \begin{itemize}[itemsep=0pt]
            \item SELinux for all userspace (5.x), runtime permissions (6.x), seccomp filter for all apps (8.x), scoped external storage (10), TLS by default (9.0), MAC randomization by default (10).
            \item Treble: about 20 HALs moved out of system server, each in its own sandbox (8.x).
            \item CFI (9.0) and Shadow Call Stack (10) in the kernel.
        \end{itemize}
        \item \textbf{Hardware-backed protections:} Keymaster (TEE), Strongbox (tamper-resistant hardware; first on Pixel 3 with Titan M), Insider Attack Resistance, Protected Confirmation.
        \item \textbf{Integrity:} Verified Boot mandated since Android 7 (exemptions up to Android 8 for slow AES devices; none from 9.0); VBMeta digest in key attestation from Android 10 on some devices.
        \item \textbf{Updates:} 2018 Android Enterprise Recommended requirement of 90-day guaranteed security updates.
        \item \textbf{Future work named by authors:} declarative encryption for apps; signing or permission for dynamically loaded code; protection if OEM signing keys leak (e.g., firmware transparency logs); defenses against RAM and hardware attacks.
    \end{itemize}

    \item \entrypart{Contribution and Limitations}

    \textbf{Contributions (authors)}
    \begin{itemize}[leftmargin=*,itemsep=1pt,topsep=2pt]
        \item First published definition of the Android security model (systematizing knowledge that existed in parts before).
        \item Definition of the threat model and how the model addresses it, including special-case handling.
        \item Explanation of how AOSP enforces the model through multiple interacting measures across layers.
        \item Identification of open gaps and room for future improvement.
    \end{itemize}

    \vspace{4pt}
    \textbf{Limitations (authors)}
    \begin{itemize}[leftmargin=*,itemsep=1pt,topsep=2pt]
        \item The model has no simple, consistent formal notation; it evolved from practical experience, not top-down theory, so ``model'' is used in a looser sense.
        \item The five rules are described as the simplest current expression of how different controls combine, and these controls follow their own models that are not necessarily consistent.
        \item Deliberate special cases deviate from the model (package listing, VPN, backup, enterprise, Factory Reset Protection); the list is not comprehensive.
        \item The platform acts both as a consenting party and as enforcer; the authors admit that some components still keep an asymmetry of power.
        \item Out of scope: Google Play Protect and Play policy, replacing apps in transit at first install (signing key trusted on first use), side-channel and hardware attacks (assumes trustworthy hardware), integrity of other CPU firmware (radio, GPU, touch).
        \item Consent is a technical enforcement notion, not legal consent.
        \item Implementation details refer to Android 11 as released.
    \end{itemize}

    \item \entrypart{My Analysis}

    \textbf{Good}
    \begin{itemize}[leftmargin=*,itemsep=1pt,topsep=2pt]
        \item Gives one coherent reference for scattered Android security design; clear five-rule summary.
        \item Threat model is explicit and each mechanism is mapped to threat IDs (Tables 1 to 4).
        \item Written by platform insiders, so implementation detail is likely accurate (observation).
        \item Honest about gaps: special cases, platform-as-party asymmetry, out-of-scope items.
        \item Useful context for the other papers in this review set (e.g., CDD, Treble, SELinux, patching).
    \end{itemize}

    \vspace{4pt}
    \textbf{Weak points}
    \begin{itemize}[leftmargin=*,itemsep=1pt,topsep=2pt]
        \item \textbf{No empirical validation:}
        \begin{itemize}[itemsep=0pt]
            \item Observation: there is no measurement showing that the model's rules hold on real devices or that listed mitigations achieve their claimed effect.
            \item Observation: the model is informal, so claims such as ``whole classes of threats become impractical'' cannot be checked against a formal property.
        \end{itemize}
        \item \textbf{Author and vendor perspective:} Observation: written by Google employees about AOSP; the Pixel and AOSP reference behavior may differ from what OEM devices ship. The paper does not discuss OEM deviation in detail (other papers in this set report such deviation).
        \item \textbf{Unpublished data:} Some figures (lockscreen adoption, vulnerability shares) are from internal or slide-deck sources and cannot be independently checked from the paper.
        \item \textbf{Design and evaluation overlap:} Not applicable, no evaluation. Observation: the tables record changes the authors chose as security-relevant, with no stated selection procedure.
        \item \textbf{Heuristics and manual steps:} Observation: ``security-relevant changes'' were identified by analysis of the code base; method and completeness are not described.
        \item \textbf{Coverage gaps:}
        \begin{itemize}[itemsep=0pt]
            \item Observation: patching and update delivery get a short section, with no data on actual update timeliness.
            \item The model assumes compatibility enforcement (rule 3) works; this is the point other papers in the set challenge (observation).
            \item Observation: rules 3 and 4 depend on CTS, Verified Boot and bootloader trust; non-certified AOSP derivatives are out of scope.
        \end{itemize}
        \item \textbf{Single viewpoint:} Observation: no comparison with other mobile OS security models beyond brief related-work pointers.
        \item \textbf{Age:} Reflects Android 11; later releases are not covered.
    \end{itemize}

    \vspace{4pt}
    \textbf{Research gaps}
    \begin{itemize}[leftmargin=*,itemsep=1pt,topsep=2pt]
        \item Formalizing the five rules so properties can be verified.
        \item Measuring whether shipped OEM ROMs actually satisfy the model (compliance, SELinux, update behavior); partly addressed by the other papers in this review set.
        \item Evaluating how well each mitigation in Tables 1 to 4 works in the field, using exploit or vulnerability data.
        \item Extending the model to post-Android 11 changes.
        \item Handling pre-installed, privileged update components (e.g., FOTA apps) inside the consent model (observation, linked to Paper 4).
        \item Authors' own open items: declarative app-file encryption, code-signing or permission for dynamic code loading, protection against stolen OEM signing keys, hardware-level attack defenses.
        \item Comparative study of Android's model against other mobile OS models (the paper defers this to other surveys).
    \end{itemize}

\end{enumerate}

\newpage

\section{Paper 6: Security Assessment of Customizations in Android Smartwatch Firmware}

\begin{enumerate}

    \item \entrypart{Paper Details}

    \begin{tabular}{@{}>{\bfseries}p{3.2cm} p{11.5cm}@{}}
        Authors          & Yifan Yu, Ruoyan Lin, Shuang Li, Qinsheng Hou, Peng Tang, Wenrui Diao \\[2pt]
        Publication Year & 2024 (IEEE 23rd International Conference on Trust, Security and Privacy in Computing and Communications (TrustCom), 2024) \\[2pt]
        DOI/Link         & \url{https://doi.org/10.1109/TrustCom63139.2024.00162}
    \end{tabular}

    \item \entrypart{Research Problem}

    \begin{itemize}[leftmargin=*,itemsep=1pt,topsep=2pt]
        \item \textbf{Problem:} Vendor customization of Android smartwatch firmware may introduce security risks. Prior firmware studies focused on smartphones, not smartwatches.
        \item \textbf{Why it matters:}
        \begin{itemize}[itemsep=0pt]
            \item Smartwatches are widely used (health monitoring, notifications).
            \item Wear OS is closed-source and needs a Google license; vendors without one build on AOSP.
            \item Vendors change functionality and configuration ``without fully understanding the implications'' (authors' claim).
            \item Vendors may not apply Google's monthly patches promptly, and sometimes change the patch date without fixing vulnerabilities (authors' claim).
            \item Vendors add pre-installed apps without ensuring their security.
        \end{itemize}
        \item \textbf{Research questions:}
        \begin{itemize}[itemsep=0pt]
            \item RQ1: How many inconsistent security configurations are in vendor custom images?
            \item RQ2: Do vendors promptly and reliably follow up on AOSP security patches?
            \item RQ3: How many pre-installed apps are insecure?
        \end{itemize}
        \item \textbf{Challenges listed by the paper:}
        \begin{itemize}[itemsep=0pt]
            \item Firmware is hard to obtain; vendors limit access to images.
            \item Large-scale static analysis cannot determine exact device node operations.
        \end{itemize}
    \end{itemize}

    \item \entrypart{Proposed Approach}

    \begin{itemize}[leftmargin=*,itemsep=1pt,topsep=2pt]
        \item \textbf{Main idea:} Measurement study of smartwatch firmware in three steps: collect firmware, build datasets, run targeted analyses for the three RQs.
        \item \textbf{Stage 1: Firmware collection and unpacking}
        \begin{itemize}[itemsep=0pt]
            \item Sources: official vendor sites and forums (e.g., XDA, needrom, androidfilehost).
            \item Formats: .zip, .img, .ofp; tools: unzip, oppo\_decrypt, brotli, sdat2img, simg2img, extfstools.
        \end{itemize}
        \item \textbf{Stage 2: Dataset construction}
        \begin{itemize}[itemsep=0pt]
            \item System file dataset from /system and /vendor.
            \item Pre-installed APK dataset from /system/app, /system/priv-app, and the vendor partition; 9,754 apps after deduplication.
        \end{itemize}
        \item \textbf{Stage 3: Differential analysis (RQ1)}
        \begin{itemize}[itemsep=0pt]
            \item Builds on Aafer et al. (security configurations) and Zhou et al. (device drivers).
            \item Six features compared across images of the same vendor: permissions (protection levels), protected broadcasts, GIDs, component visibility, component protection, device nodes.
            \item Sources: framework-res.apk manifest, \texttt{/system/etc/permissions/} XML, app manifests, \texttt{ueventd.*.rc}.
            \item Device nodes compared against Google's AOSP-based device (phone images used when no Wear image exists for that version).
        \end{itemize}
        \item \textbf{Stage 4: Patch analysis (RQ2)}
        \begin{itemize}[itemsep=0pt]
            \item \textbf{Patch delay:} gap between AOSP patch release and firmware, from build time and security patch level in \texttt{build.prop}; delay over 30 days counted as significant (per Google guidance).
            \item \textbf{Patch missing:} AndScanner (based on SnoopSnitch); results checked by comparing patch level with patch release date.
        \end{itemize}
        \item \textbf{Stage 5: APK analysis (RQ3)}
        \begin{itemize}[itemsep=0pt]
            \item MobSF static analysis; CVSS and CWE data collected.
            \item CVE matching: CVE/CPE data from CVE and NVD sites matched to app name and version (extracted with Androguard).
        \end{itemize}
    \end{itemize}

    \item \entrypart{Dataset and Evaluation}

    \begin{itemize}[leftmargin=*,itemsep=1pt,topsep=2pt]
        \item \textbf{Data sources:} Official sites and forums; 176 firmware images, 24 vendors, Android 4.4 to 13.
        \item \textbf{Per vendor (Table I):} LEMFO 31; LG 25; Finow 17; ELARI 15; KOSPET 15; Samsung 10; Ticwris 10; Others (17 vendors) 53.
        \item \textbf{Pre-installed apps:} 9,754 after deduplication.
        \item \textbf{Ground truth:} Not Specified. (No labeled vulnerability set; patch-missing results checked only by comparing patch level with release date.)
        \item \textbf{Train/test splits:} Not Specified (measurement study).
        \item \textbf{Setup/metrics:}
        \begin{itemize}[itemsep=0pt]
            \item Count and ratio of inconsistent configurations per feature and vendor (Table II).
            \item Patch delay in days; number of missing patches.
            \item Vulnerable apps per image and security issues per vulnerable app (Table III); CVSS severity; CWE category.
            \item Statistical significance tests: Not Specified.
        \end{itemize}
        \item \textbf{Baseline comparison:} Not Specified. (Only compared with related work in prose.)
    \end{itemize}

    \item \entrypart{Key Findings}

    \begin{itemize}[leftmargin=*,itemsep=1pt,topsep=2pt]
        \item \textbf{RQ1: Configurations (1,684 insecure configurations in total):}
        \begin{itemize}[itemsep=0pt]
            \item \textbf{Permissions:} 22 out of 620 permissions have inconsistent protection levels on average. Mismatch types: normal-dangerous 37\% (128/343), normal-signature 14\% (48/343), dangerous-signature 49\% (167/343).
            \begin{itemize}[itemsep=0pt]
                \item Example: \texttt{WRITE\_SETTINGS} is signature in ELARI 4G Wink but normal in ELARI KidPhone 4GR.
                \item Unspecified protection level: only 16 of 29 permissions are normal on other images; 9 are signature on at least one other image.
            \end{itemize}
            \item \textbf{Protected broadcasts:} on average 21 per vendor image set still used in a receiver's intent-filter without definition; LEMFO 58 of 801.
            \item \textbf{GIDs:} inconsistencies in firmware from 5 vendors; example LEMFO LEM8 vs.\ LEM6 camera GID (signature vs.\ normal).
            \item \textbf{Component visibility:} 1.7\% (781/45,959) of components inconsistent; LEMFO has the most (122). Example: Samsung \texttt{SleepActivity} not exported on Galaxy Watch5 Pro (Android 13), exported on Android 11.
            \item \textbf{Component protection:} 0.4\% (74/17,993) inconsistent. 33 of 74 due to different permission levels; 30 protected in one image but not another.
            \item \textbf{Device nodes:} 98 of 176 images have at least one device node with a lower protection level than AOSP. \texttt{/dev/ion} set to 666 on LEM 10 and ELARI KidPhone4G, vs.\ 664 on AOSP devices.
        \end{itemize}
        \item \textbf{RQ2: Patches:}
        \begin{itemize}[itemsep=0pt]
            \item 118 of 176 images have a delay over 30 days; average delay for those images about 425 days.
            \item Longest delay 1,611 days: ELARI KidPhone MB, built Apr 2024, Android 8.1.0, patch level November 2019.
            \item KOSPET PRIME (Android 7.1.1): ten firmware versions (1.3 built Oct 2019 to 2.6 built Sep 2020) all show patch level 2018-05-05.
            \item Missing patches: 13 in 6 images (one LEMFO, one Microwear, three Mobvoi, one ZTE); 12 Critical, one Moderate.
        \end{itemize}
        \item \textbf{RQ3: Pre-installed apps:}
        \begin{itemize}[itemsep=0pt]
            \item 26.1\% (2,543 of 9,754) apps ``fail security detection''.
            \item 9,616 issues reported; 5,309 high severity (55.2\%).
            \item Top CWE categories: logging sensitive information 2,446; incorrect permission configurations 1,440; cryptographic misuse 1,200.
            \item Samsung: about 227.5 pre-installed apps per image, 169.5 on average flagged (Table III).
            \item Case: \texttt{com.mediatek.bluetooth} on ELARI KidPhone3G (file readable by any app, sensitive logging, IP disclosure, potential SQL injection).
            \item CVE matching: no matches found (more than 1,500 CVE items compared).
        \end{itemize}
    \end{itemize}

    \item \entrypart{Contribution and Limitations}

    \textbf{Contributions (authors)}
    \begin{itemize}[leftmargin=*,itemsep=1pt,topsep=2pt]
        \item Comprehensive measurement of Android smartwatch firmware: 176 images, 24 vendors, 9,754 pre-installed apps.
        \item Analysis of three aspects (security configurations, patches, pre-installed apps) with case investigations, to draw vendors' attention to security.
    \end{itemize}

    \vspace{4pt}
    \textbf{Limitations (authors)}
    \begin{itemize}[leftmargin=*,itemsep=1pt,topsep=2pt]
        \item Limited and uneven firmware collection; some vendors give few images or none; this can cause mistakes in differential analysis.
        \item Static analysis cannot determine exact device node operations; only same-name comparison; Google's smartwatch system does not cover all Android versions, so phone images are used for missing versions.
        \item MobSF results may contain false positives.
    \end{itemize}

    \item \entrypart{My Analysis}

    \textbf{Good}
    \begin{itemize}[leftmargin=*,itemsep=1pt,topsep=2pt]
        \item Studies a new, under-examined device class (smartwatches).
        \item Covers three aspects with a clear RQ-to-method mapping.
        \item Concrete examples (\texttt{/dev/ion}, KOSPET patch level, ELARI permission mismatch).
        \item Reuses established methods (differential analysis, AndScanner, MobSF).
        \item Reports a negative result (no CVE matches) and gives a plausible explanation.
    \end{itemize}

    \vspace{4pt}
    \textbf{Weak points}
    \begin{itemize}[leftmargin=*,itemsep=1pt,topsep=2pt]
        \item \textbf{Validated vs.\ analyzed:}
        \begin{itemize}[itemsep=0pt]
            \item Observation: ``insecure'' means inconsistent with another image of the same vendor. No exploit or dynamic test is reported, so an inconsistency may be a fix rather than a regression.
            \item Observation: the direction (which image is the secure one) is assumed; vendor images that moved from weaker to stronger settings could be counted too.
            \item Observation: 26.1\% is ``fail security detection'' by MobSF, with no manual validation sample or false-positive rate; the paper itself says false positives are likely.
        \end{itemize}
        \item \textbf{Dataset bias:}
        \begin{itemize}[itemsep=0pt]
            \item Observation: uneven vendors (LEMFO 31 vs.\ Samsung 10; 17 vendors share 53 images); Table II covers only some vendors.
            \item Observation: many vendors are small or low-cost brands; results may not generalize to Wear OS or major vendors. Wear OS vs.\ AOSP-based split is not reported.
            \item Observation: firmware from forums may be modified or repacked; provenance is not discussed.
        \end{itemize}
        \item \textbf{Heuristics:}
        \begin{itemize}[itemsep=0pt]
            \item Patch delay uses \texttt{build.prop} build time and patch level, which the authors themselves say can be falsified; no independent check of build dates (observation).
            \item Patch-missing validation compares patch level with release date, which is partly circular (observation).
            \item Device node baseline mixes phone and watch AOSP images.
        \end{itemize}
        \item \textbf{Metric definitions (observation):}
        \begin{itemize}[itemsep=0pt]
            \item The ``1,684'' total is not broken down by feature in the text; Table II shows ratios per vendor, and a sum cannot be checked from it.
            \item Delay ``average 425 days'' covers only images with delay over 30 days; the average over all 176 is not given.
            \item Percent mismatch types use 343 changes, while ``22 out of 620 permissions'' is an average; link between the two is unclear.
            \item Table II device node row uses a different unit (images ratio) than other rows (configuration ratio).
        \end{itemize}
        \item \textbf{Number inconsistencies (observation):}
        \begin{itemize}[itemsep=0pt]
            \item Table I lists 10 Samsung images and Android 4.4 to 13 in text, while the table ranges start at 4.4.2 and 5.0.1.
            \item The abstract says patch delays are ``significant''; the 30-day threshold is borrowed from enterprise guidance, not smartwatch-specific.
        \end{itemize}
        \item \textbf{Single case:} Several risk claims rest on one device or one app (e.g., KOSPET PRIME, \texttt{com.mediatek.bluetooth}).
        \item \textbf{Claim wording (observation):} Statements such as ``vendors may not realize the risks'' or vendors ``just changed the date'' are not shown directly; the data shows only unchanged patch levels.
        \item \textbf{Baselines:} No comparison with smartphone firmware results on the same method; no other tool for APK analysis.
        \item \textbf{Paper length:} Short conference paper; method details (version matching, thresholds) are thin (observation).
    \end{itemize}

    \vspace{4pt}
    \textbf{Research gaps}
    \begin{itemize}[leftmargin=*,itemsep=1pt,topsep=2pt]
        \item Dynamic or exploit-based validation of flagged inconsistencies and pre-installed app issues.
        \item Separate analysis of Wear OS vs.\ AOSP-based watches and of major vs.\ small vendors.
        \item Independent verification of patch levels (e.g., binary-level checks) rather than \texttt{build.prop} metadata.
        \item Larger and more balanced firmware collection, including OTA and over-time analysis per device.
        \item Companion phone apps, Bluetooth and cloud links, which the paper does not cover (observation).
        \item Impact on users: link flagged weaknesses to real data exposure on watches (health data), not just static findings.
        \item Comparison of smartwatch results with smartphone firmware results using the same pipeline.
    \end{itemize}

\end{enumerate}

\newpage

\section{Paper 7: On the security of pre-installed Android apps in low-cost devices}

\begin{enumerate}

    \item \entrypart{Paper Details}

    \begin{tabular}{@{}>{\bfseries}p{3.2cm} p{11.5cm}@{}}
        Authors          & Alioune Diallo, Anta Diop, Abdoul Kader Kabore, Aleksandr Pilgun, Jordan Samhi, Tegawend\'{e} F.\ Bissyand\'{e}, Jacques Klein \\[2pt]
        Publication Year & Not Specified (the paper text gives no venue or year; latest cited references are from 2025) \\[2pt]
        DOI/Link         & Not Specified (code repository: \url{https://github.com/liounea/PiPLAnD})
    \end{tabular}

    \item \entrypart{Research Problem}

    \begin{itemize}[leftmargin=*,itemsep=1pt,topsep=2pt]
        \item \textbf{Problem:} Pre-installed system and vendor apps on low-cost Android devices can run with elevated privileges but get little independent scrutiny.
        \item \textbf{Why it matters:}
        \begin{itemize}[itemsep=0pt]
            \item Low-cost Android phones are widely used in Africa and help reduce the digital divide.
            \item Pre-installation can be abused to embed harmful apps that reach large user populations.
            \item Reports (cited by the paper) describe pre-installed apps exfiltrating data, harvesting financial credentials via clipboard monitoring, doing silent installs, or enrolling users in paid services without consent; some on phones sold in Africa.
            \item Android Go devices typically receive fewer updates and have a shorter support lifecycle (cited).
        \end{itemize}
        \item \textbf{Challenges listed by the paper:}
        \begin{itemize}[itemsep=0pt]
            \item Pre-installed apps are scattered across many folders, not only \texttt{/system/app} and \texttt{/system/priv-app}.
            \item Many pre-installed apps are not on Google Play; some come from unverified third-party vendors.
            \item FlowDroid's default source/sink list misses leaks through custom third-party library methods.
            \item Many steps in the pre-installation supply chain could be an entry point for malicious apps.
        \end{itemize}
        \item \textbf{Gap in prior work (authors' view):}
        \begin{itemize}[itemsep=0pt]
            \item No prior work systematically analyzed pre-installed apps on low-cost devices sold in Africa; Elsabagh et al.\ (FIRMSCOPE) examined Infinix and Tecno but focused on privilege escalation.
            \item Earlier exported-component studies look at manifests only; leak studies focus on Android/Java API methods.
        \end{itemize}
        \item \textbf{Research questions:}
        \begin{itemize}[itemsep=0pt]
            \item RQ1: How much do pre-installed apps leak sensitive data?
            \item RQ2: How much do they show suspicious behaviors (including malware, malicious URLs)?
            \item RQ3: How prevalent are manifest misconfigurations (exported sensitive components)?
        \end{itemize}
    \end{itemize}

    \item \entrypart{Proposed Approach}

    \begin{itemize}[leftmargin=*,itemsep=1pt,topsep=2pt]
        \item \textbf{Main idea:} PiPLAnD, a pipeline that pulls pre-installed APKs from physical devices with ADB and applies static analysis for leaks, suspicious behaviors, and exposed components.
        \item \textbf{Stage 1: Data collection}
        \begin{itemize}[itemsep=0pt]
            \item Seven low-cost devices (Infinix, itel, Tecno), all running Android Go Edition.
            \item APKs extracted automatically over ADB when the device is plugged in.
        \end{itemize}
        \item \textbf{Module 1: Data leak detection}
        \begin{itemize}[itemsep=0pt]
            \item FlowDroid taint analysis with its default sources and sinks.
            \item For apps with \texttt{INTERNET} permission: all methods extracted and given to LLMs for source/sink categorization; the resulting list is added to the default list.
        \end{itemize}
        \item \textbf{Module 2: Behavior analysis}
        \begin{itemize}[itemsep=0pt]
            \item Androguard decompilation; string and pattern matching (Table 2).
            \item Patterns: SMS access/delete/send (\texttt{content://sms}, \texttt{sendTextMessage}, \texttt{Telephony.SMS\_RECEIVED}); dangerous commands (\texttt{am start}, \texttt{chmod}, \texttt{su}, \texttt{sudo}, \texttt{rm -rf}); log collection (\texttt{logcat}); installation (\texttt{pm install}, \texttt{installPackage}).
            \item Findings written to JSON and manually checked.
        \end{itemize}
        \item \textbf{Module 3: Manifest misconfigurations}
        \begin{itemize}[itemsep=0pt]
            \item Extracts exported components without a protecting permission.
            \item SOOT checks each such class for sensitive API calls (list built from FlowDroid sources/sinks and Androwarn); if none, builds a call graph to look for sensitive APIs in called methods.
        \end{itemize}
    \end{itemize}

    \item \entrypart{Dataset and Evaluation}

    \begin{itemize}[leftmargin=*,itemsep=1pt,topsep=2pt]
        \item \textbf{Data sources:} Seven physical devices from Infinix, itel, Tecno (TRANSSION brands, 51\% of the African smartphone market in 2023 per Canalys, as cited); each under 100 US dollars.
        \item \textbf{Dataset size:} 1544 pre-installed APKs (system and third-party apps). Per-device counts: Not Specified.
        \item \textbf{Example signer split (Table 1):} Infinix SMART8: 68\% Infinix, 20\% Google, 4\% Transsion, 2\% Tecno, 1\% Facebook, 1\% default AOSP, 4\% others.
        \item \textbf{Ground truth:} Not Specified. Module 2 output was manually confirmed; sample size Not Specified.
        \item \textbf{Train/test splits:} Not Specified (static measurement study).
        \item \textbf{Metrics:} Number and share of apps per category (Table 3).
        \item \textbf{Baseline comparison:} Not Specified. (Prior tools are compared in prose only; no experimental comparison. VirusTotal was used for one app.)
    \end{itemize}

    \item \entrypart{Key Findings}

    \begin{itemize}[leftmargin=*,itemsep=1pt,topsep=2pt]
        \item \textbf{Overview (Table 3):}
        \begin{itemize}[itemsep=0pt]
            \item Exported sensitive components: 249 (16\%).
            \item Leak of sensitive data: 145 (9\%).
            \item Dangerous commands: 226 (15\%).
            \item Log collection: 10 (0.7\%).
            \item Silent installation behaviors: 33 (2\%).
            \item Access/send/delete SMS: 79 (5\%).
        \end{itemize}
        \item \textbf{Data leaks (RQ1):}
        \begin{itemize}[itemsep=0pt]
            \item Leaked data includes MCC, location (latitude, longitude), device info, IMSI, IMEI.
            \item Leak channels: SharedPreferences, logs, Intents, network.
            \item Case: \texttt{com.transsion.statisticalsales} sends location, IMSI, IMEI, phone version and more to a remote host through third-party library methods (Listing 1). Not flagged by VirusTotal; called malware by some blog posts (cited).
        \end{itemize}
        \item \textbf{Suspicious behaviors (RQ2):}
        \begin{itemize}[itemsep=0pt]
            \item 33 apps with silent installation behavior; 79 apps touching SMS; at least 10 apps able to read the full logcat (\texttt{READ\_LOGS}); 226 apps executing dangerous commands.
        \end{itemize}
        \item \textbf{Manifest misconfigurations (RQ3):}
        \begin{itemize}[itemsep=0pt]
            \item 249 app versions (16\%) export sensitive components without protection.
            \item Example 1: exported activity \texttt{com.transsion.carlcare.WarrantyCardActivity} reaches location data (Listing 2).
            \item Example 2: exported ContentProvider \texttt{com.sprd.providers.photos.SpecialTypesProvider} accesses media files on external storage (Listing 3).
        \end{itemize}
        \item \textbf{Ecosystem observations:}
        \begin{itemize}[itemsep=0pt]
            \item Most pre-installed apps are not on Google Play (Fig.~1).
            \item Apps sit in many folders beyond \texttt{/system/app} and \texttt{/system/priv-app} (Fig.~2).
            \item Palm Store is pre-installed on the examined devices.
        \end{itemize}
    \end{itemize}

    \item \entrypart{Contribution and Limitations}

    \textbf{Contributions (authors)}
    \begin{itemize}[leftmargin=*,itemsep=1pt,topsep=2pt]
        \item PiPLAnD pipeline to inspect pre-installed apps from physical devices (source code released).
        \item Dataset of 1544 pre-installed APKs from seven low-cost devices.
        \item Identification of manifest misconfigurations (249 app versions, 16\%), going beyond manifest-only checks by inspecting code.
        \item Quantification of sensitive-data leakage (145 apps, 9\%), including leaks through third-party libraries.
        \item Documentation of suspicious behaviors in pre-installed apps.
    \end{itemize}

    \vspace{4pt}
    \textbf{Limitations (authors)}
    \begin{itemize}[leftmargin=*,itemsep=1pt,topsep=2pt]
        \item Behavior patterns are only the ``most obvious'' ones from the literature.
        \item Only exported components lacking permissions are examined.
        \item Future work planned: analyze URLs used by apps, native libraries, more suspicious behaviors; extend to more devices; compare with European devices; add dynamic analysis.
    \end{itemize}

    \vspace{4pt}
    \textbf{Note:} Further limitations are not stated by the authors; see My Analysis.

    \item \entrypart{My Analysis}

    \textbf{Good}
    \begin{itemize}[leftmargin=*,itemsep=1pt,topsep=2pt]
        \item Targets an under-studied market and device class (low-cost devices in Africa).
        \item Extracts apps from real devices, avoiding firmware-image gaps and forum provenance issues.
        \item Extends FlowDroid with LLM-categorized custom sources and sinks to catch third-party library leaks.
        \item Code-level check of exported components, not manifest only.
        \item Concrete code listings for a real case (\texttt{statisticalsales}).
        \item Code released.
    \end{itemize}

    \vspace{4pt}
    \textbf{Weak points}
    \begin{itemize}[leftmargin=*,itemsep=1pt,topsep=2pt]
        \item \textbf{Validated vs.\ analyzed:}
        \begin{itemize}[itemsep=0pt]
            \item Observation: the headline counts are tool detections; manual confirmation is mentioned only for Module 2, with no sample size or precision figure.
            \item Observation: ``leak'' means a taint flow from source to sink; no runtime traffic check is shown, and a flow may be intended behavior (e.g., sales or warranty services).
            \item Observation: pattern matches such as \texttt{delete}, \texttt{am start}, \texttt{chmod} are very generic and likely over-count; the paper reports these as ``suspicious behaviors'' without further filtering.
        \end{itemize}
        \item \textbf{Dataset bias:}
        \begin{itemize}[itemsep=0pt]
            \item Observation: seven devices from three brands, all by one manufacturer (TRANSSION); per-device and per-brand results are not given, so conclusions may not cover ``low-cost devices'' in general.
            \item Observation: ``Africa'' is used throughout, but where devices were bought and their regional builds are not described.
            \item Observation: a Google Play lookup is shown only as a chart; counts are not given.
        \end{itemize}
        \item \textbf{Heuristics and manual steps:}
        \begin{itemize}[itemsep=0pt]
            \item LLM-based source/sink categorization: model, prompt, and accuracy are Not Specified (observation).
            \item Sensitive API list built from FlowDroid and Androwarn lists; completeness unclear.
            \item Behavior detection is string-based, so obfuscation, reflection, and native code would be missed (observation).
        \end{itemize}
        \item \textbf{Missing coverage the paper itself raises:} The Background says Android Go devices get fewer updates, but patch levels and update behavior are not measured. Malware detection and malicious URL analysis are listed in RQ2, but no malware scan results are reported besides one VirusTotal mention (observation).
        \item \textbf{Internal inconsistencies (observation):}
        \begin{itemize}[itemsep=0pt]
            \item The conclusion says the pipeline was tested on ``three different low-cost Android devices''; the abstract and Sec.~4 say seven. Answer to RQ1 also says ``three devices''.
            \item Sec.~4 says all devices run Android Go; Android Go itself is described as having less frequent updates.
            \item Contribution list says ``249 app versions'' while Table 3 and the abstract say apps; app and version counts may differ.
            \item Statement that older studies analyzed only manifests is not fully accurate for the cited works (e.g., FIRMSCOPE does code analysis) (observation).
        \end{itemize}
        \item \textbf{Single case:} The vendor exfiltration claim rests mostly on one package (\texttt{com.transsion.statisticalsales}); the abstract says it ``appears to exfiltrate'', so no network capture confirms it.
        \item \textbf{Baselines:} No comparison with FIRMSCOPE, AndScanner, FirmwareDroid, or other tools on the same devices.
        \item \textbf{Paper maturity:} Short paper; no threat model, no false-positive analysis, no statistics (observation).
    \end{itemize}

    \vspace{4pt}
    \textbf{Research gaps}
    \begin{itemize}[leftmargin=*,itemsep=1pt,topsep=2pt]
        \item Dynamic analysis and network traffic capture to confirm leaks and silent installs.
        \item Precision and recall evaluation with a labeled set; ablation of the LLM-extended source/sink list.
        \item More brands, price tiers, and regions; comparison with higher-end devices (authors' future work).
        \item Update and patch behavior on Android Go and low-cost devices (links to Paper 2).
        \item Analysis of FOTA and installer apps on these devices (links to Paper 4), and of customization and compliance (links to Paper 3).
        \item Supply-chain attribution for region-specific pre-installed apps.
        \item Risk to end users and regulatory aspects for apps sending identifiers and location.
    \end{itemize}

\end{enumerate}

\newpage

\section{Paper 8: Large-scale Security Measurements on the Android Firmware Ecosystem}

\begin{enumerate}

    \item \entrypart{Paper Details}

    \begin{tabular}{@{}>{\bfseries}p{3.2cm} p{11.5cm}@{}}
        Authors          & Qinsheng Hou, Wenrui Diao, Yanhao Wang, Xiaofeng Liu, Song Liu, Lingyun Ying, Shanqing Guo, Yuanzhi Li, Meining Nie, Haixin Duan \\[2pt]
        Publication Year & 2022 (44th International Conference on Software Engineering (ICSE '22), May 21--29, 2022, Pittsburgh, PA, USA) \\[2pt]
        DOI/Link         & \url{https://doi.org/10.1145/3510003.3510072}
    \end{tabular}

    \item \entrypart{Research Problem}

    \begin{itemize}[leftmargin=*,itemsep=1pt,topsep=2pt]
        \item \textbf{Problem:} Android fragmentation and vendor customization create security risks. Prior work mostly built analysis tools or studied one aspect of firmware. No panoramic, large-scale view of firmware ecosystem security.
        \item \textbf{Why it matters:}
        \begin{itemize}[itemsep=0pt]
            \item Android has 85\% smartphone market share (as cited).
            \item Many vendors do not follow Google's monthly patches in time; some only change the patch date without fixing vulnerabilities (cited).
            \item Pre-installed apps can be vulnerable (e.g., CVE-2019-15442 in Samsung Themes, CVE-2019-15378 in WfoService from MediaTek on a Panasonic device).
        \end{itemize}
        \item \textbf{Challenges listed by the paper:}
        \begin{itemize}[itemsep=0pt]
            \item Confirming CVEs related to pre-installed apps.
            \item Efficient large-scale security patch detection.
            \item Selecting detection features for pre-installed app security.
        \end{itemize}
        \item \textbf{Research questions:}
        \begin{itemize}[itemsep=0pt]
            \item Q1: Do vendors follow up AOSP security patches in time?
            \item Q2: Are vendors' patching claims reliable?
            \item Q3: How many pre-installed apps contain unfixed CVEs?
            \item Q4: Are ``non-vulnerable'' pre-installed apps secure?
        \end{itemize}
        \item \textbf{Threat model:} Attacker can induce the victim to click a URL, install an app, and even physically touch the device.
    \end{itemize}

    \item \entrypart{Proposed Approach}

    \begin{itemize}[leftmargin=*,itemsep=1pt,topsep=2pt]
        \item \textbf{Main idea:} AndScanner, an automated framework for ROM crawling, ROM parsing, patch analysis, and app analysis, applied to a large firmware dataset.
        \item \textbf{Crawler}
        \begin{itemize}[itemsep=0pt]
            \item ROM images from vendor sites; for vendors without direct links (e.g., Samsung, Huawei), URLs built by changing parameters of a valid link.
            \item Android Dumps (public project) for small vendors; integrity checked against crawled firmware and by reproducing its unpacking tool.
            \item CVE data 2009 to 2020: IDs, dates, descriptions from CVE site; CPE data from NVD; matched to pre-installed app metadata, then manually re-checked to remove duplicates and irrelevant CVEs.
        \end{itemize}
        \item \textbf{ROM Parser}
        \begin{itemize}[itemsep=0pt]
            \item Unpacks images (encrypted ones via dynamic methods such as ADB to get keys; simg2img and other public tools).
            \item Extracts metadata from \texttt{build.prop}: fingerprint, Android version, build time, security patch level, model, vendor.
            \item Files tagged by tuple $<$File MD5, vendor, ROM fingerprint$>$.
        \end{itemize}
        \item \textbf{Patch Analyzer}
        \begin{itemize}[itemsep=0pt]
            \item \textbf{Patch Delay Module (Q1):} delay = build time minus security patch date; delays under 30 days not counted (Google recommendation).
            \item \textbf{Patch Missing Module (Q2):} Python re-implementation of SnoopSnitch's binary signature verification; 411 test cases for 411 CVEs and 520 patch files (Oct 2015 to May 2020, Android 5.0 to 10); multi-process; analysis time per ROM cut by ten times vs.\ the original app.
            \item A patch counts as missing if absent in the binary and released before or in the image's patch level month.
            \item Test cases related to runtime state were discarded, which may cause false positives.
        \end{itemize}
        \item \textbf{App Analyzer (Androguard, CryptoGuard)}
        \begin{itemize}[itemsep=0pt]
            \item \textbf{CVE Matching Module (Q3):} matches package name, version, signature, file MD5 to CVEs; checks vulnerable code or misconfigured attributes; compares CVE disclosure time to ROM build time.
            \item \textbf{Risk Detection Module (Q4):} features chosen because over 90\% (176/191) of pre-installed app CVEs are attribute misconfiguration (168) or cryptographic misuse (8).
            \item Attribute features: F1 exported components, F2 debuggable, F3 networkSecurityConfig, F4 allowBackup, F5 sharedUserId.
            \item Crypto misuse: 8 CryptoGuard high-risk rules (Table 1), e.g., constant keys, trust-all certificates, insecure hashes.
            \item Exploitability checks with a drozer-based scanner and hand-built PoCs on the authors' own devices.
        \end{itemize}
    \end{itemize}

    \item \entrypart{Dataset and Evaluation}

    \begin{itemize}[leftmargin=*,itemsep=1pt,topsep=2pt]
        \item \textbf{Data sources:} Vendor sites, support channels, Android Dumps; collected August to November 2020.
        \item \textbf{Dataset size:} 6,261 firmware images, 153 vendors, 8.6 TB; Android 2.x to 11; build time 2012 to 2020; 443,475 unique pre-installed APKs; 602 CVEs (411 official Android CVEs and 191 related to pre-installed apps).
        \item \textbf{Top vendors (Table 2):} Google 1,211; Huawei 978; OPPO 826; Samsung 707; Xiaomi 524; Motorola 397; OnePlus 173; Vivo 139; Nokia 110; Realme 101; Others (143 vendors) 1,095.
        \item \textbf{Ground truth:} Not Specified as a labeled set. Accuracy check: 50 random patch-missing results manually verified; one suspected false positive.
        \item \textbf{Train/test splits:} Not Specified (measurement study).
        \item \textbf{Metrics:} Patch delay in months; number of missing patches; images with unfixed app CVEs; counts of potential risks per image (F1 to F5, crypto).
        \item \textbf{Implementation:} Python and shell, 12,278 lines of code; three servers.
        \item \textbf{Baseline comparison:} Not Specified. (Compared with partner vs.\ non-partner vendors and supported vs.\ non-supported devices, and with related work in prose; no baseline tool compared.)
    \end{itemize}

    \item \entrypart{Key Findings}

    \begin{itemize}[leftmargin=*,itemsep=1pt,topsep=2pt]
        \item \textbf{Q1 Patch delay:}
        \begin{itemize}[itemsep=0pt]
            \item 1,516 of 6,261 images (24.2\%) have at least one month delay; overall average about 0.8 months; for delayed images about 3.2 months.
            \item Worst case: OPPO A33m latest image built Jan 2019 (Android 5.1.1) with patch date Oct 2015, 40 months.
            \item Delay proportion: Xiaomi 46.8\%, Vivo 69.1\%, OnePlus 39.9\%, OPPO 37.2\%, Samsung 23.3\%, Huawei 2.9\%, Nokia 0.9\% (1 image), Google 0\%.
            \item Treble: of 4,838 images on Android 8.0 or later, 782 (16.2\%) delayed, average 2.8 months; before 8.0, around 51.6\% delayed, average 4.0 months.
            \item Xiaomi models released after 2017 show no delay; older models delayed (Xiaomi 4: 16 months).
            \item Partner vs.\ non-partner vendors: 0.83 vs.\ 1.37 months. Supported vs.\ non-supported devices: 5.18 vs.\ 10.96 months.
            \item OPPO R9s vs.\ A57: A57 delay grows to a peak of 19 months; last two A57 update logs mention security patches, but patch levels did not change.
        \end{itemize}
        \item \textbf{Q2 Patch missing:}
        \begin{itemize}[itemsep=0pt]
            \item 383 of 6,261 images (6.1\%) from 50 vendors have at least one missing patch.
            \item Most extreme: OPPO R9skt with 46 missing patches (built May 2017, patch level Apr 2017).
            \item OPPO: 161 images (19.5\%), average about 1.5 missing; Nokia and Realme: none.
            \item Google: 51 images (4 models) miss CVE-2017-0494 (moderate).
            \item OPPO A59: missing patches grow from 1 (2016 build) to 23 (2018 build).
            \item STS (introduced Aug 2018): average missing patches lower after STS for all top 10 vendors.
            \item Partner vs.\ non-partner: 0.97 vs.\ 1.42 missing; supported vs.\ non-supported: 4.04 vs.\ 11.09.
        \end{itemize}
        \item \textbf{Q3 Unfixed CVEs in pre-installed apps:}
        \begin{itemize}[itemsep=0pt]
            \item 462 images (7.4\%) from 71 vendors (46.4\%) contain unfixed app CVEs.
            \item Realme 89 images, 129 CVE apps; OPPO 53 images, 85 apps; Google, Huawei, OnePlus: none.
            \item \texttt{DropboxChmodService} appears in all affected OPPO and Realme images.
            \item Case: Nokia 6.2 (built Jan 2020) with \texttt{com.google.android.tag} (CVE-2019-9295) and \texttt{com.qualcomm.qti.callenhancement} (CVE-2019-15473).
            \item Scope (Table 7): \texttt{com.google.android.tag} 840 ROMs; \texttt{com.dropboxchmod} 328 ROMs; \texttt{callenhancement} v9 in 251 ROMs with 15 undisclosed vendors.
            \item Five vulnerable apps appear in images built in 2020, at least one year after disclosure.
        \end{itemize}
        \item \textbf{Q4 Other security risks:}
        \begin{itemize}[itemsep=0pt]
            \item 3,595 images (57.4\%) from all 153 vendors contain potentially vulnerable pre-installed apps.
            \item Average potential risks per image: Realme 505 (highest), Xiaomi 408, OnePlus 351, Motorola 339, Huawei 81 (lowest); Google 116.
            \item Sources of potentially vulnerable apps: phone vendors 71.5\% (87,825/122,770); third-party developers 15.7\% (19,330).
            \item Exploitability: 307 exploitable vulnerabilities found and submitted on the authors' devices; 38 confirmed from 15 vendors, 32 with CVE/CNVD IDs.
            \item Cases: debuggable \texttt{com.autonavi.manu.widget} on a OnePlus ROM; hard-coded AES key in Samsung \texttt{smartcallprovider}; Vivo \texttt{WeatherProvider} leaks location (CVE-2021-26279; 24 app versions, 109 Vivo ROMs); a Google settings provider issue reported to nine vendors: four confirmed, four ignored, one (Google) still reviewing.
        \end{itemize}
    \end{itemize}

    \item \entrypart{Contribution and Limitations}

    \textbf{Contributions (authors)}
    \begin{itemize}[leftmargin=*,itemsep=1pt,topsep=2pt]
        \item Largest Android firmware dataset used for security measurement at the time: 6,261 images, 153 vendors, 602 CVEs.
        \item AndScanner: automated framework for ROM crawling, parsing, patch analysis, and app analysis.
        \item Measurements and case studies answering Q1 to Q4; findings reported to vendors; 38 vulnerabilities confirmed, 32 with CVE/CNVD IDs.
        \item Source code and experiment data to be released (GitHub).
    \end{itemize}

    \vspace{4pt}
    \textbf{Limitations (authors)}
    \begin{itemize}[leftmargin=*,itemsep=1pt,topsep=2pt]
        \item CVE reports lack a standard format; missing affected versions prevent matching to apps.
        \item Static analysis cannot handle packed apps (0.23\%, 1,020/443,475, of apps).
        \item Dropping runtime-state patch test cases may cause false positives in patch-missing results.
    \end{itemize}

    \item \entrypart{My Analysis}

    \textbf{Good}
    \begin{itemize}[leftmargin=*,itemsep=1pt,topsep=2pt]
        \item Very large and diverse dataset covering both major and small vendors.
        \item Combines OS-level patch checks with app-level checks; binary-level patch verification goes beyond trusting patch-level strings.
        \item Follows up with exploit PoCs and vendor disclosure, not only static counts.
        \item Looks for causes through longitudinal case studies (OPPO R9s vs.\ A57) and group comparisons (partner, supported).
        \item Reports vendor reactions to disclosure.
    \end{itemize}

    \vspace{4pt}
    \textbf{Weak points}
    \begin{itemize}[leftmargin=*,itemsep=1pt,topsep=2pt]
        \item \textbf{Validated vs.\ analyzed:}
        \begin{itemize}[itemsep=0pt]
            \item Observation: patch-missing accuracy rests on 50 manual checks out of 383 flagged images' results; no recall estimate.
            \item The 57.4\% ``potentially vulnerable'' figure is a static feature count. Only 307 were confirmed exploitable, on the authors' own devices, out of a very large flagged set (observation: exploitability rate over all flagged apps is not given).
        \end{itemize}
        \item \textbf{Overlap of design and evaluation data:} Risk features (F1 to F5, crypto) were chosen from the same 191 CVEs, which are also used for CVE matching (observation); the patch test set (411 CVEs) is limited to SnoopSnitch's coverage, Oct 2015 to May 2020.
        \item \textbf{Dataset and vendor bias:}
        \begin{itemize}[itemsep=0pt]
            \item Observation: Google has 1,211 images (19\%) and is used as a clean reference; Others (143 vendors) hold 1,095 images, so many vendor-level conclusions rest on few images (e.g., Nokia 110, Realme 101).
            \item Observation: images come from crawling and Android Dumps (community); provenance and representativeness of what users actually run are not shown; images are not weighted by device sales.
            \item Observation: the delay measure uses images, not devices or users; an image may never have reached users.
        \end{itemize}
        \item \textbf{Heuristics and metric definitions:}
        \begin{itemize}[itemsep=0pt]
            \item Delay = build time $-$ patch date, rounded down to months; delays under 30 days dropped. Build time can be falsified, and this is not discussed (observation).
            \item The 30-day threshold comes from an enterprise whitepaper (cited), not tested.
            \item Observation: Q3 ``unfixed'' compares CVE disclosure time with ROM build time and app version; apps patched at the module level with unchanged version strings could be mislabeled.
            \item Observation: causal wording (``insufficient security efforts'', ``developer gave up maintaining'', ``excessive trust'') is not tested; Table 4 shows correlations only.
        \end{itemize}
        \item \textbf{Confounding (observation):} Partner/supported status, price tier, and Android version are correlated; no control for them. The Treble comparison (16.2\% vs.\ 51.6\%) also mixes newer, better-maintained devices with older ones; the paper itself says the contribution cannot be quantified.
        \item \textbf{Number inconsistencies (observation):}
        \begin{itemize}[itemsep=0pt]
            \item Delay proportion: Table 2 says Xiaomi 245/524 = 46.8\%, Huawei 29/978 = 2.9\%; Table 4's average months differ in scale (supported 5.18 vs.\ overall about 0.8 months), so the table 4 averages likely cover only delayed images; this is not stated.
            \item Huawei: 2 images with missing patches, but text says Realme ``has not gone into the market before August 2018'' only in a footnote for STS$-$.
            \item Google patch-missing images 51 (4.2\%), while Pixel and Nexus are the reference for patches, which the authors explain by one moderate CVE.
            \item Section 4.1 says images ``could all be unpacked''; no failure rate is given.
        \end{itemize}
        \item \textbf{Single case:} Several explanations (e.g., OPPO treats high-end models better) rest on one pair of models.
        \item \textbf{Baselines:} No comparison with other patch-checking tools or firmware scanners (e.g., FIRMSCOPE) on the same images.
        \item \textbf{Age:} Data collected Aug to Nov 2020; Android up to 11; findings may have changed.
    \end{itemize}

    \vspace{4pt}
    \textbf{Research gaps}
    \begin{itemize}[leftmargin=*,itemsep=1pt,topsep=2pt]
        \item Validate patch-missing detection against ground truth (known patched and unpatched builds); recall not measured.
        \item Separate vendor intent from technical causes of missing patches (the paper offers only possible reasons).
        \item Measure update delivery to real devices and users rather than ROM images (links to Paper 2).
        \item Extend to devices beyond phones, e.g., smartwatches (Paper 6), low-cost devices (Paper 7), and FOTA/update components (Paper 4).
        \item Link ROM customization and compliance (Paper 3) to patch delay and missing patches with controlled analysis.
        \item Assess exploitability at scale rather than a sample of 307 on the authors' own devices.
        \item Patch verification beyond SnoopSnitch's CVE coverage, and for kernel, vendor and non-AOSP components.
        \item Update the analysis for post-Android 11 releases, Mainline modules, and Treble maturity.
    \end{itemize}

\end{enumerate}

\newpage

\section{Paper 9: Keep Me Updated: An Empirical Study of Proprietary Vendor Blobs in Android Firmware}

\begin{enumerate}

    \item \entrypart{Paper Details}

    \begin{tabular}{@{}>{\bfseries}p{3.2cm} p{11.5cm}@{}}
        Authors          & Elliott Wen, Jiaxing Shen, Burkhard Wuensche \\[2pt]
        Publication Year & 2024 (IEEE 30th International Conference on Parallel and Distributed Systems, ICPADS) \\[2pt]
        DOI/Link         & \url{https://doi.org/10.1109/ICPADS63350.2024.00025}
    \end{tabular}

    \item \entrypart{Research Problem}
    \begin{itemize}[leftmargin=*, itemsep=1pt, topsep=2pt]
        \item \textbf{Problem:} proprietary vendor blobs in Android firmware get little security research.
        \begin{itemize}[itemsep=0pt]
            \item Blobs sit in the \texttt{/vendor} partition.
            \item They support hardware such as GPUs, cameras, fingerprint readers.
            \item Many are loaded into apps or system services with high privileges.
        \end{itemize}
        \item \textbf{Why it matters:}
        \begin{itemize}[itemsep=0pt]
            \item Vulnerable blobs can ``significantly impact the security of the entire system.''
            \item GPU blobs are loaded into every app for rendering UI.
            \item GPU blobs are easy to reach: Chrome, Firefox, WebKit can forward arbitrary GPU task inputs to them.
            \item Driver source is given to manufacturers (e.g., Google, Xiaomi) under NDA, so update responsibility rests entirely with manufacturers.
        \end{itemize}
        \item \textbf{Research questions:}
        \begin{itemize}[itemsep=0pt]
            \item RQ1: What is the update pattern of GPU vendor blobs?
            \item RQ2: How vulnerable are these GPU vendor blobs?
        \end{itemize}
        \item \textbf{Challenges listed by the paper:}
        \begin{itemize}[itemsep=0pt]
            \item Closed-source blobs, no documentation, limited scrutiny.
            \item No version metadata inside the binary.
            \item Debug version strings may be missing if aggressive optimizations are enabled.
            \item Existing GPU fuzzer (GraphicsFuzz) needs physical device access.
            \item GraphicsFuzz speed: less than ten shaders per second, too slow for the number of blobs.
            \item Sampler and half-precision float operations are costly or unsupported on x86.
            \item Chrome adds ANGLE, which sanitizes shaders before the vendor blob.
        \end{itemize}
    \end{itemize}

    \item \entrypart{Proposed Approach}
    \begin{itemize}[leftmargin=*, itemsep=1pt, topsep=2pt]
        \item \textbf{Main idea:} automatic pipeline \emph{GPUBlob-Inspector}.
        \begin{itemize}[itemsep=0pt]
            \item Crawl firmware, extract blobs, identify versions, fuzz offline, report results.
            \item Offline fuzzing: capture LLVM IR from the blob, run it on x86 CPUs. No physical device needed.
            \item Manual analysis validates bugs on physical devices.
        \end{itemize}
        \item \textbf{1. Firmware crawler and extractor:}
        \begin{itemize}[itemsep=0pt]
            \item Web crawler for vendor sites (Google, Xiaomi publish all versions).
            \item Third-party subscription services for Samsung and OnePlus: SamMobile, Daxiaamu.
            \item Public repository Android Dump for lesser-known vendors.
            \item Metadata collected: firmware version, release date, supported regions.
            \item Unpacking via an adapted Android firmware tool (srlabs extractor).
            \item GPU blobs in \texttt{/vendor/lib64/egl}: \texttt{libGLES.so} or \texttt{libEGL.so}, \texttt{libGLESv1\_CM.so}, \texttt{libGLESv2.so}.
            \item Dependencies found by recursive \texttt{readelf} on \texttt{DT\_NEEDED}.
        \end{itemize}
        \item \textbf{2. Version identifier (three-part scheme):}
        \begin{itemize}[itemsep=0pt]
            \item Build ID from \texttt{.note.gnu.build-id}. Identifies a blob but gives no order.
            \item Internal blob version from debug logs via regex.
            \begin{itemize}[itemsep=0pt]
                \item Qualcomm: \texttt{EV\{major\}.\{minor\}.\{patch\}}.
                \item ARM: \texttt{r\{major\}p\{patch\}}.
            \end{itemize}
            \item LLVM compiler version, used when internal version is missing.
            \begin{itemize}[itemsep=0pt]
                \item Regex on version string if not stripped.
                \item Stripped binaries: string fingerprints from LLVM logging statements.
                \item Also Binshape [16] function fingerprinting.
            \end{itemize}
        \end{itemize}
        \item \textbf{3. Offline fuzzer (metamorphic testing):}
        \begin{itemize}[itemsep=0pt]
            \item Semantically equivalent shader variants should give the same output. A significant difference signals a potential bug.
            \item Both Qualcomm and ARM GPU blobs are built on LLVM (frontend $\rightarrow$ optimizer $\rightarrow$ backend).
            \item Fuzzer captures IR after optimization and reroutes it to an x86-64 backend instead of the GPU backend.
            \item \textbf{Reference shader database:}
            \begin{itemize}[itemsep=0pt]
                \item List of 145,732 gaming apps from Google Play, APKPure, APKMirror (as of Feb 2024).
                \item 10 genres, 20 games per genre, weighted random sampling by download count.
                \item Patrace used to record GLES calls; \texttt{glShaderSource} calls filtered out.
                \item One researcher played each game for at least 10 minutes: about 2000 minutes of traces, 18,923 shader programs recovered.
            \end{itemize}
            \item \textbf{Variant shader generation:} transformations adapted from GraphicsFuzz, chained randomly and recursively.
            \begin{itemize}[itemsep=0pt]
                \item Statement mutation (e.g., \texttt{mix(origin, unused, c)} with $c=1$).
                \item Control flow mutation (\texttt{if}$\leftrightarrow$\texttt{switch}, \texttt{for}$\leftrightarrow$\texttt{while}, loop wrap, unroll, split).
                \item Code donation (donated code with renamed variables, results stored to unused outputs).
            \end{itemize}
            \item \textbf{LLVM IR capture:}
            \begin{itemize}[itemsep=0pt]
                \item Standard GL path: \texttt{eglCreateContext}, \texttt{glCreateShader}, \texttt{glShaderSource}, \texttt{glCompileShader}.
                \item LLVM option \texttt{-print-after-all} injected via \texttt{cl::ParseCommandLineOptions}; address found by string fingerprint.
                \item Full-system option: QEMU running AOSP, blobs swapped in via \texttt{debug.gles.layers}, plus a shim kernel module that stubs most IOCTL calls.
                \item Qualcomm shortcut: call \texttt{libllvm-glnext.so} directly (\texttt{dlopen}, \texttt{dlsym}) in a user-space ARM emulator (aarch-qemu). Main function: \texttt{QGLCCompileToIRShader}.
            \end{itemize}
            \item \textbf{CPU binary generation:}
            \begin{itemize}[itemsep=0pt]
                \item Toolchain: \texttt{llc}, \texttt{llvm-mc}, \texttt{lld}.
                \item Intrinsics (\texttt{llvm.qgpu.*}) reimplemented in C, compiled with Clang.
                \item Shader I/O: GPU addresses remapped to \texttt{mmap} heap; random inputs with predetermined seeds; output memory hashed.
                \item Math: intrinsics mapped to \texttt{libm}; half-precision promoted to single via an LLVM pass; \texttt{fpext} converted to \texttt{bitcast}.
                \item Sampler: simplified to a hash function; same parameters and seed give the same texture value.
                \item ABI fix: Clang \texttt{ext\_vector\_type} attribute for float vector passing.
            \end{itemize}
        \end{itemize}
        \item \textbf{4. Result reporting:}
        \begin{itemize}[itemsep=0pt]
            \item Outdatedness analysis, inconsistency analysis, physical device case study.
            \item \textbf{Update delay $D$:} for a firmware released at $r_f$ with blob $v_o$:
            \begin{itemize}[itemsep=0pt]
                \item $r_o$ = release date of the oldest firmware with $v_o$.
                \item $v_l$ = most recent blob version, same GPU model, in firmware released before $r_f$.
                \item $r_l$ = release date of the earliest firmware with $v_l$.
                \item $D = r_l - r_o$.
            \end{itemize}
            \item \textbf{Validation on device:} replace reference shader in Patrace retrace file with variant, replay offscreen, compare frames to reference using chi-square test on histograms.
            \item \textbf{Root-cause tools:} stack traces, Valgrind, GDB, \texttt{-print-after-all}, Data Dependency Graph (DDG) comparison.
        \end{itemize}
    \end{itemize}

    \item \entrypart{Dataset and Evaluation}
    \begin{itemize}[leftmargin=*, itemsep=1pt, topsep=2pt]
        \item \textbf{Firmware dataset:}
        \begin{itemize}[itemsep=0pt]
            \item 13,901 firmware images, 38 TBs of disk space.
            \item Released between January 2018 and April 2024.
            \item Vendor count: Not Specified (consistently). Text states 74 phone vendors in the Introduction and 24 phone vendors in Section III-A.
            \item Date range is also inconsistent: ``Jan 2017 to April 2024'' in the Introduction vs.\ ``January 2018 and April 2024'' elsewhere.
            \item Top five vendors (Samsung, Xiaomi, Oppo, OnePlus, Google) contribute 93\% of images.
            \item Nvidia Tegra and MediaTek Dimensity excluded due to minimal market shares.
            \item De-duplicated by Build ID: 792 distinct blob versions.
        \end{itemize}
        \item \textbf{Shader data:} 18,923 shader programs from 200 sampled games (10 genres $\times$ 20).
        \item \textbf{Fuzzing setup:} 200 variants per reference shader.
        \item \textbf{Hardware:} Ubuntu 24.04, two AMD EPYC 7742 64-core CPUs, 1024 GB RAM, 128 TB SSD, 10 Gbps network.
        \item \textbf{Pipeline timing:}
        \begin{itemize}[itemsep=0pt]
            \item Firmware crawling (about 13,000 images): 27 hours.
            \item Firmware extraction: 147 minutes (average 38 seconds per firmware, max 278 seconds).
            \item Version identification: 22 minutes (average 12 seconds per blob, max 19 seconds).
            \item Variant generation: about 491 minutes.
            \item LLVM IR generation: 189 minutes.
            \item x86-64 compilation: 127 minutes.
            \item Execution and comparison: about 7 minutes.
        \end{itemize}
        \item \textbf{Fuzzing target:} one Qualcomm blob only, from Xiaomi Note 13 Pro 5G firmware (released February 23, 2024), internal version \texttt{031.42.23.11}.
        \item \textbf{Metrics:}
        \begin{itemize}[itemsep=0pt]
            \item Share of firmware with outdated blobs; update delay $D$ in days (median, max).
            \item Average firmware and GPU blob update count per device.
            \item LLVM version distribution.
            \item Number of anomalies found; fuzzing speed (shaders per second).
        \end{itemize}
        \item \textbf{Ground truth / validation:} anomalies replayed on physical devices with Patrace; chi-square test on frame histograms.
        \item \textbf{Baseline comparison:} only for fuzzing speed vs.\ GraphicsFuzz (server on AMD 5990x, 128 GB RAM; client on ASUS ROG Phone 8). No baseline for bug-finding count, update-delay method, or version identification.
    \end{itemize}

    \item \entrypart{Key Findings}
    \begin{itemize}[leftmargin=*, itemsep=1pt, topsep=2pt]
        \item \textbf{GPU composition:}
        \begin{itemize}[itemsep=0pt]
            \item 78\% of images contain Qualcomm GPU blobs; only 21\% contain ARM GPU blobs.
            \item 792 distinct versions: 65\% Qualcomm, 35\% ARM.
            \item Table I total: Firmware 78\% / 22\%, Distinct Driver 60\% / 40\%, Device 69\% / 31\% (Qualcomm / ARM).
            \item Qualcomm mostly in mid-range to high-end devices; ARM in entry-level models.
        \end{itemize}
        \item \textbf{Update outdatedness:}
        \begin{itemize}[itemsep=0pt]
            \item 82\% of firmware images contain outdated GPU blobs.
            \item Median delay $D$: 273 days; highest value: 1,128 days.
            \item Abstract says ``up to 1,281 days''; body and conclusion say 1,128 days. Not Specified which is correct.
            \item Introduction says ``on average'' 273 days; Section IV-B says ``median''.
            \item Qualcomm: 85\% of firmware outdated, median delay 281 days.
            \item ARM: 64\% of firmware outdated, median delay 231 days.
            \item Google updates GPU drivers significantly more promptly than others (Fig.\ 6).
        \end{itemize}
        \item \textbf{Update count per device (Fig.\ 5):}
        \begin{itemize}[itemsep=0pt]
            \item Firmware updates far exceed GPU blob updates.
            \item GPU blob update counts per device: Samsung 4.31, Xiaomi 1.52, Oppo 1.74, Google 6.29, OnePlus 3.93.
            \item Xiaomi and Oppo: average GPU driver update count close to one.
        \end{itemize}
        \item \textbf{LLVM outdatedness:}
        \begin{itemize}[itemsep=0pt]
            \item All Qualcomm GPU blobs rely on LLVM 2.8, ``released nearly 15 years ago.''
            \item Majority of ARM blobs on LLVM 9.0, then 10.0, then 11.0; remaining use LLVM 5.0.
            \item LLVM 11.0 released in 2020.
            \item Introduction says ``over 15 years ago''; Section IV says ``nearly 15 years''; conclusion says ``almost 15 years''.
            \item Introduction says the most recent LLVM dependency dates back 5 years.
        \end{itemize}
        \item \textbf{Fuzzing results (one Qualcomm blob):}
        \begin{itemize}[itemsep=0pt]
            \item About 14 hours of fuzzing: 289 instances of irregular program behaviors or incorrect program outputs.
            \item Abstract calls them ``security and behavioral bugs''; paper text confirms many are behavioral and ``do not necessarily lead to security vulnerability.''
            \item \textbf{Memory access violation:} 29 cases, crash at \texttt{glCompileShader}.
            \begin{itemize}[itemsep=0pt]
                \item 5 null pointer references; rest invalid memory access, mostly low addresses (e.g., \texttt{0x7}, \texttt{0x1b}).
                \item 4 addresses are heap memory; Valgrind suggests ``likely'' use-after-free.
                \item Location: 7 in LLVM analysis passes; 17 in function optimization passes; 5 in backend optimization passes.
            \end{itemize}
            \item \textbf{Compilation stalls:} 3 instances, infinite iteration of optimization passes (loop unrolling, CFG simplification, instruction combination).
            \item \textbf{Incorrect program semantics:} described as the most common type; pinpointed flawed passes: instruction combination, dead code elimination, loop simplification. Cause: \texttt{undef} values left after wrongly removed instructions.
            \item Per-category count for incorrect semantics: Not Specified.
        \end{itemize}
        \item \textbf{Fuzzing speed:} GraphicsFuzz about 7 shaders per second; this fuzzer about 75 shaders per second.
        \item \textbf{Across blob versions:} problematic shaders tested on 10 blob versions, evenly sampled by time. They trigger the same anomalous behaviors in all of them.
        \item \textbf{WebGL case study:}
        \begin{itemize}[itemsep=0pt]
            \item Attack path: attacker hosts page with malicious shader; victim visits via system browser or hybrid app (phishing, messages, rogue Wi-Fi).
            \item Cites [18]: 149 of the top 500 most downloaded apps are hybrid.
            \item With ANGLE in the fuzzing loop: 4-hour session found 9 sanitized shaders that crash the Chrome browser.
            \item Working end-to-end DoS or code-execution exploit: Not Specified.
        \end{itemize}
        \item \textbf{Preliminary extensions (Discussion):}
        \begin{itemize}[itemsep=0pt]
            \item 19 vulnerabilities found in Mali GPU blobs (preliminary).
            \item Other blobs (accelerometer, fingerprint) measured preliminarily; results Not Specified.
        \end{itemize}
    \end{itemize}

    \item \entrypart{Contribution and Limitations}

    \textbf{Contributions (authors)}
    \begin{itemize}[leftmargin=*, itemsep=1pt, topsep=2pt]
        \item Large-scale empirical study of vendor blob update patterns and security implications, focused on GPU blobs.
        \item Analysis of over 13,000 Android firmware releases.
        \item Automatic pipeline \emph{GPUBlob-Inspector}: firmware crawler, unpacker, version identifier, offline fuzzer.
        \item Versioning scheme for blobs without metadata (Build ID, internal version, LLVM version).
        \item Performant offline fuzzer needing no physical mobile device (about 75 shaders per second vs.\ about 7 for GraphicsFuzz).
        \item Findings: outdated GPU blobs in 82\% of firmware; obsolete LLVM (2.8) in Qualcomm blobs; 289 bugs and anomalies in the recent Qualcomm blob.
        \item WebGL case study.
        \item Source code of the pipeline and the dataset released publicly.
    \end{itemize}

    \vspace{4pt}
    \textbf{Limitations (authors)}
    \begin{itemize}[leftmargin=*, itemsep=1pt, topsep=2pt]
        \item Fuzzer bypasses the GPU backend (instruction scheduling, register allocation), so it cannot detect bugs there. Authors say reported data is a ``conservative estimate.''
        \item Study focuses on OpenGL ES; Vulkan not supported yet (needs a Vulkan variant shader generator).
        \item Focus on Qualcomm GPUs; ARM results are preliminary, ``further investigation is warranted.''
        \item Other vendor blobs (e.g., fingerprint, accelerometer) only measured preliminarily. Comparison of update patterns is complicated by many hardware manufacturers.
        \item ANGLE sanitization required a changed fuzzing pipeline for Chrome.
        \item Blobs without debug version strings (aggressive optimization) need the LLVM fallback.
    \end{itemize}

    \item \entrypart{My Analysis}

    \textbf{Good}
    \begin{itemize}[leftmargin=*, itemsep=1pt, topsep=2pt]
        \item Large dataset (13,901 images, 792 distinct blob versions) over about six years.
        \item Practical versioning scheme for closed binaries with no metadata.
        \item Offline fuzzing idea (IR capture $\rightarrow$ x86) removes device dependency and scales.
        \item Fuzzing speed compared against GraphicsFuzz (state of the art).
        \item Found anomalies validated on a physical device with a stated method (chi-square on histograms).
        \item Root-cause analysis by category (stack traces, Valgrind, GDB, DDG).
        \item Code and dataset released, which supports reproducibility.
        \item Authors openly state that the fuzzer misses backend bugs.
    \end{itemize}

    \vspace{4pt}
    \textbf{Weak points}
    \begin{itemize}[leftmargin=*, itemsep=1pt, topsep=2pt]
        \item \textbf{Validation scope (observation):} fuzzing run on one Qualcomm blob only (Xiaomi Note 13 Pro 5G, \texttt{031.42.23.11}). Results on one device/blob; generalization to 792 versions is only partly tested (10 versions for repeat triggering).
        \item \textbf{Validated vs.\ analyzed:}
        \begin{itemize}[itemsep=0pt]
            \item 289 instances reported as validated on physical devices, but the paper gives per-category numbers for only 29 memory violations and 3 stalls. The rest are not broken down.
            \item Only 4 heap-address cases are inferred as likely use-after-free, via Valgrind. No confirmed exploit.
            \item ``Security'' count is much smaller than 289; headline of ``289 security and behavioral bugs'' can read stronger than the evidence for security impact.
        \end{itemize}
        \item \textbf{Overclaim risk (observation):} abstract says vulnerabilities ``can be exploited via WebGL'', but the case study only reports 9 sanitized shaders that crash Chrome. No demonstrated code execution reported.
        \item \textbf{Staleness-as-vulnerability:} outdated LLVM implies possible exposure to known LLVM bugs, but the paper does not map specific known LLVM bugs to blob crashes. The paper states this as ``could potentially be susceptible.''
        \item \textbf{Delay metric (observation):} $D$ depends on firmware in the authors' own dataset. If the dataset misses a firmware, $r_o$, $v_l$, $r_l$ can be off. ``Outdated'' means older than a blob in another device's firmware, not older than the vendor's latest release. Compatibility of a newer blob with the target device is not checked.
        \item \textbf{Mean vs.\ median:} text mixes ``on average'' and ``median'' for 273 days. Not clear which is correct.
        \item \textbf{Internal inconsistencies:} 74 vs.\ 24 vendors; Jan 2017 vs.\ Jan 2018; 1,281 vs.\ 1,128 days; ``over 15'' vs.\ ``nearly 15'' years. Reduces trust in the reported numbers.
        \item \textbf{Vendor bias:} data dominated by five vendors (93\%) and Qualcomm (78\% of images). Third-party firmware sources (SamMobile, Daxiaamu) may not represent real OTA rollout. Region and carrier effects are not discussed (observation).
        \item \textbf{Heuristic dependence:}
        \begin{itemize}[itemsep=0pt]
            \item Version identification relies on regex, debug strings, and LLVM string/Binshape fingerprints. Accuracy of version identification: Not Specified.
            \item Reference shaders from one researcher playing 10 minutes per game; selection of 200 games and manual play may bias shader coverage.
            \item Root-cause analysis and device validation are manual; effort and false-positive rate: Not Specified.
            \item Sampler simplified to a hash and half-float promoted to single; effect on false positives/negatives: Not Specified.
        \end{itemize}
        \item \textbf{Design/evaluation overlap:} GraphicsFuzz transformations reused; the same shaders and one blob used to develop the pipeline and to report bugs. Held-out evaluation: Not Specified.
        \item \textbf{Baselines:} no comparison of bug-finding effectiveness with GraphicsFuzz (only speed). Speed comparison uses different hardware (server PC + phone vs.\ 64-core servers), so it is not like-for-like (observation). No ablation of mutation types.
        \item \textbf{Metric definitions:} the 289 count is described as ``instances''; whether they are deduplicated by root cause is Not Specified.
        \item \textbf{Causal claims (observation):} reasons given for Google updating more (more resources, first-party control) and for LLVM not being updated (API breaking changes) are speculation, not tested.
        \item \textbf{Problem not covered:} paper says other blobs (camera, fingerprint) matter, but only GPU is studied in depth. The binderized, high-privilege blobs are left out.
        \item \textbf{Real-world impact:} no responsible disclosure or vendor response reported (Not Specified).
    \end{itemize}

    \vspace{4pt}
    \textbf{Research gaps}
    \begin{itemize}[leftmargin=*, itemsep=1pt, topsep=2pt]
        \item Fuzz GPU backend stages (scheduling, register allocation); authors' GPU instruction interpreter is still in progress.
        \item Extend to Vulkan (SPIR-V frontend) and to ARM Mali beyond the preliminary 19 findings.
        \item Study other vendor blobs (camera, fingerprint, sensors), especially high-privilege binderized blobs.
        \item Fuzz many blob versions and vendors, not one Qualcomm blob, and report per-version bug sets.
        \item Link crashes to specific known LLVM bugs to check whether outdated LLVM actually causes them.
        \item Build a working end-to-end WebGL exploit and test across browsers and hybrid apps, with and without ANGLE.
        \item Measure the real security effect of blob update delay, e.g., by checking whether newer blobs fix the found bugs.
        \item Add comparison with other GPU fuzzers on bug count, not only speed.
        \item Evaluate version-identification accuracy against ground truth.
        \item Study why OEMs do not update blobs (dependency on chip vendor SoC support, API-breaking LLVM changes) with direct evidence.
    \end{itemize}

\end{enumerate}

\newpage

\section{Paper 10: Pleasure or pain? An evaluation of the costs and utilities of bloatware applications in android smartphones}

\begin{enumerate}

    \item \entrypart{Paper Details}

    \begin{tabular}{@{}>{\bfseries}p{3.2cm} p{11.5cm}@{}}
        Authors          & Haroon Elahi, Guojun Wang, Jianer Chen \\[2pt]
        Publication Year & 2020 (Journal of Network and Computer Applications, vol.\ 157, article 102578; received 12 June 2019, accepted 10 February 2020, available online 19 February 2020) \\[2pt]
        DOI/Link         & \url{https://doi.org/10.1016/j.jnca.2020.102578}
    \end{tabular}

    \item \entrypart{Research Problem}
    \begin{itemize}[leftmargin=*, itemsep=1pt, topsep=2pt]
        \item \textbf{Problem:} Android bloatware (pre-installed apps) has privacy, security, and trust concerns, but the paper says there is little empirical support for them.
        \begin{itemize}[itemsep=0pt]
            \item Bloatware definition used: all apps that come pre-installed in a smartphone, following McDaniel (2012).
            \item Users cannot remove these apps by conventional means.
            \item Removal attempts may cause security vulnerabilities, unstable OS, and void warranty.
        \end{itemize}
        \item \textbf{Second problem:} claims that bloatware gives more utility and covers diverse needs of different end-user market segments are not validated by any study.
        \item \textbf{Why it matters:}
        \begin{itemize}[itemsep=0pt]
            \item Users cannot remove the apps, so privacy invasion can last a long time.
            \item The paper cites cases of OEMs, third-party providers, and operators using bloatware to collect data beyond operational needs.
        \end{itemize}
        \item \textbf{Goal:} check whether the indirect costs (privacy, security, trust) are comparable to the utilities.
        \item \textbf{Gaps the paper lists:}
        \begin{itemize}[itemsep=0pt]
            \item Only one published study gives a basic overview of smartphone bloatware (McDaniel, 2012).
            \item Smartphone bloatware work lacks quantitative support.
            \item Desktop bloat studies are not enough: different device nature, data, permission model, sensors.
        \end{itemize}
        \item \textbf{Challenges stated:}
        \begin{itemize}[itemsep=0pt]
            \item Firmware formats vary; open-source tools do not support some image formats.
            \item Huawei and Motorola firmware could not be included for this reason.
        \end{itemize}
    \end{itemize}

    \item \entrypart{Proposed Approach}
    \begin{itemize}[leftmargin=*, itemsep=1pt, topsep=2pt]
        \item \textbf{Main idea:} two empirical studies. (1) Permission analysis of bloatware extracted from firmware. (2) Online user study on utility and coverage claims.
        \item \textbf{Study 1: Firmware and permission analysis}
        \begin{itemize}[itemsep=0pt]
            \item \textbf{Collection:} firmware of 100 models, nine OEMs, from firmwarefile.com; HTC firmware from the HTC official support portal.
            \item \textbf{Extraction tools:} WinZip, Simg2img, Ext2explore to get APK files.
            \item \textbf{Analysis tools:} Androguard API 3.1, Python 3.6 custom scripts, Microsoft Excel, APK Analyzer in Android Studio.
            \item \textbf{Privilege measure:} number and types of permissions (after Dini et al., 2018).
            \begin{itemize}[itemsep=0pt]
                \item Over-privileged = requests more permissions than needed for standard features (combines Felt et al., 2011 and Wang et al., 2016).
                \item Among similar apps, the one with minimum permissions defines the standard features.
            \end{itemize}
            \item \textbf{Manual audit:} partition testing into nine strata by brand; exploratory approach; 35\% of apps audited in each stratum; about 6000 apps audited in total.
            \item \textbf{Permission groups for analysis:}
            \begin{itemize}[itemsep=0pt]
                \item Protection levels: Normal, Dangerous, Signature, Custom.
                \item Six functional groups: personal data access and manipulation, location, telephony, internet and network connectivity, device features, framework and system level operations.
            \end{itemize}
            \item \textbf{Treatment of instances:} the same app in different models is counted as different applications.
        \end{itemize}
        \item \textbf{Study 2: Online user study}
        \begin{itemize}[itemsep=0pt]
            \item Survey built in Google Forms; dichotomous, multiple-choice, and open-ended questions (10 questions, Appendix B).
            \item Shared on Facebook, LinkedIn, WeChat public groups, and student mailing lists of two universities in Pakistan.
            \item Anonymous collection; no ethical approval sought because all participants were adults.
            \item \textbf{Five questions studied:}
            \begin{itemize}[itemsep=0pt]
                \item Do users derive large utilities from bloatware?
                \item Do users make frequent use of bloatware and know it well?
                \item Do users use different bloatware due to different backgrounds?
                \item Does a large amount of bloatware reduce dependence on third-party apps?
                \item Do users want advanced-feature apps shipped with the phone?
            \end{itemize}
            \item \textbf{Open answers:} axial coding (Corbin and Strauss, 1990) linked to Google Play app categories.
        \end{itemize}
    \end{itemize}

    \item \entrypart{Dataset and Evaluation}
    \begin{itemize}[leftmargin=*, itemsep=1pt, topsep=2pt]
        \item \textbf{Firmware data:}
        \begin{itemize}[itemsep=0pt]
            \item 17,179 APK files from 100 smartphone models, nine OEMs; models released between 2012 and 2018.
            \item Permissions extracted for 16,959 applications.
            \item Brands (N models): Haier 10, HTC 11, Lenovo 10, Oppo 10, Panasonic 9, Samsung 20, Vivo 10, Xiaomi 17, ZTE 3.
            \item Android versions: twenty versions, 4.0.3 to 8.1.
            \item Model list in Table A.15.
        \end{itemize}
        \item \textbf{Ground truth / validation:} manual audit with APK Analyzer of about 6000 apps (35\% per brand stratum) to confirm implementations match requested permissions.
        \item \textbf{User study data:}
        \begin{itemize}[itemsep=0pt]
            \item 190 responses received; 5 empty, 5 duplicates; 180 analyzed.
            \item Target audience about 3000 people; the paper says the overall response was very low.
            \item 120 male, 60 female.
            \item Age: 41 aged 18--20; 85 aged 21--25; 23 aged 26--30; 31 above 30.
            \item 81\% Android, 16\% iPhone, 3\% other phones.
            \item Eighteen countries.
            \item Professions: academia (26 teachers and 100 students), IT and software 16, research 8, business 6, engineering 4, household 4, banking 2, law 2, sports 1, fashion 1, psychology 1; 9 did not disclose.
            \item Missing values: 8 no profession; 1 no country; 14 no frequent-app names; only 144 answered the must-have question; 7 no third-party count; only 141 answered the last question.
        \end{itemize}
        \item \textbf{Metrics:}
        \begin{itemize}[itemsep=0pt]
            \item Number of apps per device (mean, std.\ dev., min, max).
            \item Permissions per app; share of apps requesting each protection level.
            \item Average number of apps per phone requesting specific permissions.
            \item Survey response percentages.
        \end{itemize}
        \item \textbf{Baseline comparison:} Not Specified.
    \end{itemize}

    \item \entrypart{Key Findings}
    \begin{itemize}[leftmargin=*, itemsep=1pt, topsep=2pt]
        \item \textbf{Number of apps:}
        \begin{itemize}[itemsep=0pt]
            \item Average 172 bloatware apps per phone (Table 1: N = 100, std.\ dev.\ 72, min 60, max 374).
            \item Samsung: mean 286 (min 192, max 374). Haier and Lenovo: mean 114 each.
            \item Max of 374 is on Samsung Galaxy S8; 2.87 GB storage for the 374 apps.
            \item Samsung and Vivo show an increase in app count with newer Android versions (Figs.\ 1, 2); other brands show little variation.
        \end{itemize}
        \item \textbf{Permission levels (Tables 5, 6):}
        \begin{itemize}[itemsep=0pt]
            \item Mean share of apps requesting: Normal 72\%, Dangerous 60\%, Signature 52\%, Custom 39\%.
            \item Samsung: 66\% Dangerous, 67\% Signature, 68\% Custom.
            \item Some apps in every phone use no permissions (Table 3: 11\% Samsung to 29\% HTC).
        \end{itemize}
        \item \textbf{Top critical permissions (Table 4):} WRITE\_SETTINGS, WRITE\_EXTERNAL\_STORAGE, READ\_PHONE\_STATE, ACCESS\_NETWORK\_STATE, ACCESS\_COARSE\_LOCATION, READ\_EXTERNAL\_STORAGE, WRITE\_MEDIA\_STORAGE, READ\_CONTACTS, GET\_ACCOUNTS, USE\_CREDENTIALS.
        \item \textbf{Examples for Samsung (average apps per phone):}
        \begin{itemize}[itemsep=0pt]
            \item READ\_CONTACTS 59 vs.\ WRITE\_CONTACTS 33; apps that read data outnumber apps that write it.
            \item READ\_PHONE\_STATE 95; INTERNET 122; WRITE\_EXTERNAL\_STORAGE 125.
            \item ACCESS\_FINE\_LOCATION 44; ACCESS\_COARSE\_LOCATION 38.
            \item CAMERA 31; RECORD\_AUDIO 27.
            \item 29 apps can make calls; 17 can make calls without user intervention.
        \end{itemize}
        \item \textbf{Critical operations:} on average per phone, 32 apps get lists of running tasks (87 on Samsung); 10 can delete packages; 10 can force stop apps; 7 can download without notifying the user (14 on Xiaomi).
        \item \textbf{Permission count per app (brand-wise):}
        \begin{itemize}[itemsep=0pt]
            \item Haier 16.27; HTC 12.15; Lenovo 17.51; Oppo 11.42; Panasonic 16.58; Samsung 19.86; Vivo 15.38; Xiaomi 17.51; ZTE 16.12.
            \item Maximum for a single app: 219 permissions, Alipay on Samsung.
        \end{itemize}
        \item \textbf{Examples of apps with permissions that look unrelated to function (paper's own cases):}
        \begin{itemize}[itemsep=0pt]
            \item Haier: ``Lovely Fonts'' accesses location, screen content, device ID, SIM state, subscriber ID.
            \item HTC: Google ``Velvet'' (max permissions, 91); ``VideoHighlights'' reads profile, location, contacts, phone state.
            \item Lenovo: ``LSF-Device-Pad'' known to collect user information; ``Malwarebytes'' detected it as malware.
            \item Oppo: calculator asks for phone state; ``OppoBluetooth'' reads call logs.
            \item Panasonic: Bluetooth apps claim 17 Dangerous permissions.
            \item Samsung: ``ContextProvider'' known to collect user data; ``Wallpaper Cropper'' requests SMS reading.
            \item Vivo: ``BBKNotes'' and ``Bluetooth'' apps request SMS, call, contact permissions.
            \item Xiaomi: ``Calculator'' requests location and phone state; ``CleanMaster'' requests many sensitive permissions.
            \item ZTE: weather app requests running tasks, outgoing calls, manage accounts.
        \end{itemize}
        \item \textbf{User study results (n = 180):}
        \begin{itemize}[itemsep=0pt]
            \item 24\% take services from bloatware as a primary use; 37\% mark third-party apps as a primary use.
            \item Bloatware used: 0--5 apps 40\%, 6--10 apps 35\%, 11--15 apps 16\%, 16--20 apps 4\%, 20+ apps 5\% (Table 13).
            \item Third-party apps installed after purchase: 0--5 apps 31\%, 6--10 apps 35\%, 11--15 apps 19\%, 16--20 apps 8\%, 20+ apps 7\% (Table 13).
            \item Usefulness (Q5): 50\% say some are useful to some extent; 35\% say most; 12\% say all; one respondent disagreed.
            \item Professional life (Q6): 12\% all, 38\% most, 46\% some, 4\% none.
            \item No bloatware (Q9): 25\% totally disappointed, 41\% somewhat disappointed, 26\% does not matter, 8\% happy.
            \item Most-used bloatware categories (Table 14): Communication 64\%, Productivity 35\%, Business 27\%.
            \item No link found between respondent background and bloatware used or wanted.
            \item No advanced-feature demand found in must-have or dream-app answers.
        \end{itemize}
        \item \textbf{Authors' conclusion:} pains (privacy, security, trust) are apparent; pleasures (utilities) are questionable and need further investigation.
        \item \textbf{Abstract vs.\ body numbers:} abstract says 39\% use 0--5 bloatware apps; Table 13 and analysis say 40\%.
    \end{itemize}

    \item \entrypart{Contribution and Limitations}

    \textbf{Contributions (authors)}
    \begin{itemize}[leftmargin=*, itemsep=1pt, topsep=2pt]
        \item Original analysis of bloatware in 100 Android models from nine leading OEMs; average of 172 apps per phone.
        \item Finding that most apps are highly privileged and collaborative; they can affect privacy, resource integrity, and trust, and widen the attack surface.
        \item Online user study (180 participants) testing the utility and coverage claims.
        \item Use of both analyses to compare pains and pleasures of bloatware.
        \item Contribution to a research gap that needs more attention.
        \item Recommendations:
        \begin{itemize}[itemsep=0pt]
            \item Limit bloatware number and abilities in proportion to practical utility.
            \item OEMs should run risk evaluation before including apps.
            \item OEMs could offer feature apps online, secured with product keys shipped with the phone.
            \item Create user-class application profiles.
        \end{itemize}
    \end{itemize}

    \vspace{4pt}
    \textbf{Limitations (authors)}
    \begin{itemize}[leftmargin=*, itemsep=1pt, topsep=2pt]
        \item Firmware mostly from firmwarefile.com; the claim that it hosts official firmware is believed but not verified.
        \item Firmware selection limited by open-source tools; other brands could not be included.
        \item Privacy, security, and trust analysis rests only on permission requests in manifest files; other metrics may give different results.
        \item Survey results can be affected by question wording and over-reporting bias; the authors tried to reduce this by asking the same question in different ways.
        \item Sample and sample size can be questioned; the authors suggest more research with different user samples and sizes.
        \item Claims about utilities need more research (Section 7.4).
    \end{itemize}

    \item \entrypart{My Analysis}

    \textbf{Good}
    \begin{itemize}[leftmargin=*, itemsep=1pt, topsep=2pt]
        \item Combines a static firmware study with a user study; addresses both the cost side and the benefit side.
        \item Large app set (17,179 APKs) over nine brands and many Android versions.
        \item Permission data is broken down by protection level, by function, and by brand, with per-brand examples.
        \item Manual audit of about 6000 apps to check implementations.
        \item Tools used are open source, which the authors say supports replication.
        \item Authors state limitations on firmware source, permission-only metric, and survey bias.
        \item Survey questions are listed in Appendix B and model list in Appendix A.
    \end{itemize}

    \vspace{4pt}
    \textbf{Weak points}
    \begin{itemize}[leftmargin=*, itemsep=1pt, topsep=2pt]
        \item \textbf{Permission as proxy for risk (observation):}
        \begin{itemize}[itemsep=0pt]
            \item Requested permission is not the same as actual data use or leakage. No dynamic or network traffic analysis is reported (Not Specified).
            \item Terms such as ``abuse'', ``violate trust'', and ``advertisement purposes'' (Section 6.3) go beyond what a permission list shows. This reads as stronger than the evidence (observation).
            \item Over-privilege rule: ``standard features'' defined by the app with the fewest permissions in a group of similar apps. Grouping rule and error rate are Not Specified.
        \end{itemize}
        \item \textbf{Validated vs.\ analyzed:}
        \begin{itemize}[itemsep=0pt]
            \item About 6000 of 17,179 apps audited (35\% per stratum); how many audit findings differ from the extracted permissions is Not Specified.
            \item Permissions for 16,959 of 17,179 apps; reason for the 220 missing is Not Specified.
        \end{itemize}
        \item \textbf{Counting method (observation):} the same app in different models is treated as a different application. Counts can be inflated by repeated apps (e.g., one Samsung link-sharing app listed in six models). No count of distinct apps is given.
        \item \textbf{Defining ``bloatware'' very broadly:} every pre-installed app counts, including Google apps, system components, keyboard, dialer, and GmsCore. This mixes core system apps with truly unwanted apps (observation). The paper accepts this as a choice.
        \item \textbf{Vendor and sample bias:}
        \begin{itemize}[itemsep=0pt]
            \item Models per brand differ (3 ZTE to 20 Samsung); ZTE has N = 3.
            \item Huawei and Motorola excluded because of tooling.
            \item Many Haier, Panasonic, and old models with low market share (observation); firmware from a third-party site.
            \item Carrier-specific and region-specific bloatware not discussed (observation).
        \end{itemize}
        \item \textbf{Survey design problems (observation):}
        \begin{itemize}[itemsep=0pt]
            \item Convenience sample; about 3000 reached, 190 responded; response rate not given and described as very low.
            \item 100 of 180 are students and 26 teachers; most aged 18--25; skewed to a university population. Authors still say the sample is representative because all use smartphones.
            \item Only 81\% Android users; 16\% iPhone users counted anyway.
            \item Respondents may not know which apps are pre-installed. Survey does not show how a respondent tells bloatware from other apps; the paper itself says respondents may be unaware of how many are installed.
            \item Several answers use apps that are not usually bloatware (e.g., Facebook, Twitter, Instagram, YouTube listed as frequently used bloatware).
            \item Question 9 shows 66\% disappointed without bloatware; the authors still conclude utility is low (see next point).
        \end{itemize}
        \item \textbf{Interpretation of results (observation):}
        \begin{itemize}[itemsep=0pt]
            \item Q9: 25\% totally disappointed and 41\% somewhat disappointed, i.e., 66\% would be disappointed. Q5: 97\% say some, most, or all are useful (50\% + 35\% + 12\%). The paper says these results ``negate'' the assumption of large utility. Alternative reading is possible.
            \item Comparing the 37\% (third-party) with the 24\% (bloatware) uses different question wording; they are not like-for-like measures. No statistical tests reported (Not Specified).
            \item ``Number of bloatware apps used'' is self-reported and compared with 172 installed apps; this does not by itself show the rest is unused or redundant.
            \item Claim of no link between background and app use is based on visual reading of Table 14, no test; many profession groups have 1--8 respondents.
        \end{itemize}
        \item \textbf{Inconsistencies in text (observation):}
        \begin{itemize}[itemsep=0pt]
            \item Abstract 39\% (0--5 apps) vs.\ 40\% in Table 13.
            \item Section 8 says 37\% report most bloatware useful; Section 5.2.6.2 says 35\% (and 37\% is the third-party primary-use figure).
            \item Section 5.2.4.8 says nineteen Xiaomi phones; Table 1 says N = 17.
            \item Section 5.2.4.3 shows Lenovo permissions-per-app 17.51, same as Xiaomi 17.51.
            \item Overall-permission text says HTC 71\% and Samsung 89\% (derived from Tables 3 and 5; unclear method).
            \item Q5 in Appendix B has four options; analysis text numbers do not map clearly.
        \end{itemize}
        \item \textbf{Design/evaluation overlap:} permission-based risk definition set before analysis and also used to report findings; no held-out check (Not Specified).
        \item \textbf{No baselines:} no comparison with non-bloatware apps from app stores, with other permission-risk scoring, or with user-installed apps (Not Specified). So ``high'' permission numbers have no reference.
        \item \textbf{Causal claims (observation):} link between larger app count and more permission requests, and Android version trend (Figs.\ 1, 2), are shown descriptively. Android version, brand, and release year are mixed.
        \item \textbf{Resource claims:} only storage measured (2.87 GB, one device). Effects on data, power, memory, CPU are stated as ``can be expected'' and not measured.
        \item \textbf{Not covered although described as a real problem:}
        \begin{itemize}[itemsep=0pt]
            \item Real exploitation of bloatware; vulnerabilities in the apps themselves (the paper only argues from the attack surface and past CVEs in Android OS).
            \item Data uploads to servers: stated as ability, not observed.
            \item No vendor response or disclosure reported (Not Specified).
        \end{itemize}
        \item \textbf{Single-source and age:} firmware from 2012--2018 only; results may not hold for later Android versions with changed permission models (observation).
    \end{itemize}

    \vspace{4pt}
    \textbf{Research gaps}
    \begin{itemize}[leftmargin=*, itemsep=1pt, topsep=2pt]
        \item Dynamic and network-traffic analysis to check whether bloatware actually uses its permissions and sends data out.
        \item Compare bloatware permissions against user-installed apps of the same category.
        \item Count distinct apps and separate core system apps, OEM apps, operator apps, and third-party apps.
        \item Larger and more varied user samples, with questions that let respondents identify which apps are pre-installed; add statistical tests.
        \item Measure real resource costs (memory, battery, CPU, data) of bloatware.
        \item Include brands left out for tooling reasons (e.g., Huawei, Motorola) and newer Android versions.
        \item Study vulnerabilities in bloatware code and whether they have been exploited.
        \item Test the authors' proposed model of optional online feature apps with product keys, and user profile--based app selection.
        \item Study why users remove bloatware (creepiness vs.\ low utility); the paper argues for both reasons without direct evidence.
    \end{itemize}

\end{enumerate}

\newpage

\section{Paper 11: Deploying Android Security Updates: an Extensive Study Involving Manufacturers, Carriers, and End Users}

\begin{enumerate}

    \item \entrypart{Paper Details}

    \begin{tabular}{@{}>{\bfseries}p{3.2cm} p{11.5cm}@{}}
        Authors          & Kailani R. Jones, Ting-Fang Yen, Sathya Chandran Sundaramurthy, Alexandru G. Bardas \\[2pt]
        Publication Year & 2020 (2020 ACM SIGSAC Conference on Computer and Communications Security, CCS '20, November 9--13, 2020, Virtual Event, USA) \\[2pt]
        DOI/Link         & \url{https://doi.org/10.1145/3372297.3423346}
    \end{tabular}

    \item \entrypart{Research Problem}
    \begin{itemize}[leftmargin=*, itemsep=1pt, topsep=2pt]
        \item \textbf{Problem:} Android's fragmented ecosystem makes delivery of security updates and OS upgrades ``cumbersome and complex.''
        \begin{itemize}[itemsep=0pt]
            \item Many parties involved: AOSP, chipset vendors, device manufacturers, carriers, carrier partners, third-party testing labs, end users.
            \item Unclear how effective Google projects (Android One, Project Treble, Project Mainline) are at delivering updates to end-user devices.
        \end{itemize}
        \item \textbf{Why it matters:}
        \begin{itemize}[itemsep=0pt]
            \item Devices may be behind the latest Android Security Bulletins ``by months.''
            \item Latency leaves time for adversaries; the paper cites that the majority of exploits are published within two weeks of public release.
            \item Findings may help consumers, policy makers, and software vendors.
        \end{itemize}
        \item \textbf{Goal:} once a security update is available from AOSP, follow its path through manufacturers and carriers to end devices. Measure update frequency and latency. Also track an OS upgrade (Android 10).
        \item \textbf{Gaps the paper lists:}
        \begin{itemize}[itemsep=0pt]
            \item Past work covers limited facets (AOSP patch latency, manufacturers, or end-user side).
            \item FTC report covered a limited dataset (only eight manufacturers).
            \item No prior study of the full rollout process across AOSP, manufacturers, carriers, and end users.
        \end{itemize}
        \item \textbf{Challenges stated:}
        \begin{itemize}[itemsep=0pt]
            \item Carriers and manufacturers post announcements in inconsistent and changing formats, if at all.
            \item No standard reporting mechanism.
            \item Locked vs.\ unlocked devices complicate who controls timing.
            \item Up to 1,500 variations of the same update for one manufacturer.
        \end{itemize}
    \end{itemize}

    \item \entrypart{Proposed Approach}
    \begin{itemize}[leftmargin=*, itemsep=1pt, topsep=2pt]
        \item \textbf{Main idea:} quantitative measurement study. Join four data sources and measure rollout frequency and latency after the AOSP patch level date.
        \item \textbf{1. Android Security Bulletins (AOSP):}
        \begin{itemize}[itemsep=0pt]
            \item 53 monthly bulletins, August 2015 to December 2019.
            \item Baseline = first patch level date in each bulletin.
            \item Collected CVE type, severity, affected components.
        \end{itemize}
        \item \textbf{2. Carrier dataset:}
        \begin{itemize}[itemsep=0pt]
            \item Public update announcements from AT\&T, Verizon, T-Mobile, Sprint.
            \item Custom parser per carrier website: PyQt5, BeautifulSoup, Selenium, regex and HTML tags.
            \item Manual collection of AT\&T URLs and manual validation of collected fields.
            \item Verizon shows only the latest three updates per model; collected three times in 2019 (May, October, December).
            \item Missing bulletin: assume the bulletin of the announcement month (conservative; gives a lower bound on latency).
            \item Model release dates from GSMArena and PhoneArena.
        \end{itemize}
        \item \textbf{3. Manufacturer dataset (Samsung):}
        \begin{itemize}[itemsep=0pt]
            \item Samsung per-model update pages (constructed URLs), 15 Galaxy models.
            \item Locked (T-Mobile, Sprint), unlocked, and unbranded variants.
        \end{itemize}
        \item \textbf{4. End-user dataset:}
        \begin{itemize}[itemsep=0pt]
            \item Anonymized HTTP POST logs from Android devices of a U.S.-based social network app, 2019.
            \item Four fields: request date, hashed account ID, user-agent string, country code.
            \item Build number in user-agent used as proxy for security updates.
            \item ``Device'' = unique pair of hashed user ID and model variant.
            \item ``Active'' device = accessed the app at least once a month for more than 10 months in 2019.
            \item Joined with carrier dataset by model variant and carrier.
        \end{itemize}
        \item \textbf{Metrics:}
        \begin{itemize}[itemsep=0pt]
            \item \textbf{Normalized update frequency:} $N_{act}/N_{pot}$. $N_{pot}$ = months from later of release date and 2015 data start to $T_{end}$ (last update). $N_{act}$ = observed updates.
            \item \textbf{Rollout latency:} days from bulletin publication (usually the 1st of the month) to carrier/manufacturer rollout.
            \item \textbf{Device update latency:} days between carrier announcement and observed build change on device.
        \end{itemize}
        \item \textbf{Statistics:} Spearman correlation between latency/frequency and model age, CVE count, CVE severity, type, classification, skipped bulletins, days since last update.
        \item \textbf{Initiative analysis:} Treble devices (427,291 in end-user data), Android One (1,048 devices), Mainline-supported models (140,720 devices); compared via build numbers, OS version, and latency.
    \end{itemize}

    \item \entrypart{Dataset and Evaluation}
    \begin{itemize}[leftmargin=*, itemsep=1pt, topsep=2pt]
        \item \textbf{Carrier data:}
        \begin{itemize}[itemsep=0pt]
            \item 3,493 raw updates collected; 1,867 ``valid'' security update announcements used (the Introduction also gives 1,953 and Section 4.2 says 1,540 were removed; Appendix A.2 says 1,627 excluded).
            \item 274 unique models, 25 manufacturers.
            \item Table 1 (models): AT\&T 74, Sprint 93, T-Mobile 111, Verizon 149.
            \item Table 1 (announcements): AT\&T 394, Sprint 364, T-Mobile 658, Verizon 537.
            \item Normalization: 9 steps (Table 4): remove no release date, no bulletin indicator, outside August 2015--December 2019, prepaid duplicates, non-security updates, no market release year, higher bulletin levels, and outliers (delay less than $-10$ or greater than 170 days).
        \end{itemize}
        \item \textbf{Samsung data:} 684 announcements for 13 Samsung models across locked, unlocked, unbranded (Section 5.2); Table 2 lists 13 unique models, Section 4.4 says 15 models.
        \item \textbf{End-user data:}
        \begin{itemize}[itemsep=0pt]
            \item 152,156,934 HTTP POST requests from Android devices.
            \item 9,163,277 unique user identifiers; 4,800,228 unique user-agent strings.
            \item 7,247 unique models across 454 device manufacturers.
            \item 1,273,571 active devices used for analysis.
            \item 877,499 (68.90\%) active devices joined with carrier dataset (Samsung 75.90\%, Google 15.94\%, LG 5.39\%, Motorola 2.41\%).
            \item Data gap June 11 to July 23, 2019 (bug).
        \end{itemize}
        \item \textbf{Ground truth / validation:}
        \begin{itemize}[itemsep=0pt]
            \item Manual validation of programmatically collected carrier fields.
            \item Gap sensitivity check: repeat without post-gap builds; first 5 vs.\ last 5 months.
            \item Ground truth for actual device patch level: Not Specified (build number used as proxy).
        \end{itemize}
        \item \textbf{Ethics:} public sites rate-limited (every 2--8 seconds); end-user data anonymized; IRB office determined data is de-identified and not a human subject.
        \item \textbf{Baseline comparison:} Not Specified as a quantitative baseline. Findings are compared in text with FTC report [5] and prior work. Treble vs.\ non-Treble and locked vs.\ unlocked are internal comparisons.
        \item \textbf{Code and data:} collection scripts and public data released on GitHub.
    \end{itemize}

    \item \entrypart{Key Findings}
    \begin{itemize}[leftmargin=*, itemsep=1pt, topsep=2pt]
        \item \textbf{Update frequency:}
        \begin{itemize}[itemsep=0pt]
            \item Average announcements per model: 4.54 ($\pm$4.67).
            \item Average normalized update frequency: 0.35 across all models; T-Mobile highest at 0.40.
            \item Model age vs.\ normalized frequency: correlation $-0.3265$ (p-value $3.3541^{-8}$). T-Mobile: $-0.5054$ (p-value $8.0681^{-7}$).
            \item Updates drop sharply around 36 months of model age; models older than three years are 43.59\% of all models.
            \item LG with T-Mobile: median frequency 0.42; other carriers 0.12 to 0.33. LG frequency capped at 0.75.
            \item Samsung on AT\&T: median 0.54; 0.28 to 0.32 with other carriers (Appendix A.4).
        \end{itemize}
        \item \textbf{Rollout latency (carrier/manufacturer):}
        \begin{itemize}[itemsep=0pt]
            \item Median delay of 24 days from bulletin to rollout; largely unchanged over four years.
            \item AT\&T median 29.5 days; T-Mobile 19 days.
            \item Model age vs.\ latency: correlation 0.0502 (p-value 0.0321), reported as not significant.
            \item 21 models supported by all four carriers: latency differs per carrier. Motorola median from 10 days (T-Mobile) to 37 days (AT\&T).
            \item Locked Samsung models receive updates 4 days earlier on average than unlocked ones; contradicts the FTC report finding.
            \item CVE severity, type, count, and skipped bulletins: mostly no significant effect. Weak positive correlation with days since last update and with CVEs in the ``System'' category (0.45 to 0.61 in three of four carriers).
        \end{itemize}
        \item \textbf{End-user devices (security updates):}
        \begin{itemize}[itemsep=0pt]
            \item Median 11 days (average 15.12 days) from carrier announcement to observed build change.
            \item Devices access app on average every 6.12 days (median six days).
            \item Average 3.44 build numbers per device in 2019.
            \item 70\% of devices updated by day 14; curves flatten around day 14.
            \item Pixel 2 XL on Verizon (30,545 devices): 99.86\% updated after announcement; 30.36\% first seen with new build; 69.50\% seen with old build first.
            \item Paper says user-incurred delay is ``significantly smaller'' than carrier/manufacturer or batch delay.
        \end{itemize}
        \item \textbf{OS upgrade (Android 10, released September 3, 2019):}
        \begin{itemize}[itemsep=0pt]
            \item 19.93\% upgraded in first week; 25.09\% within two weeks; median 20 days to first observation.
            \item Google and Essential devices upgraded very quickly; Samsung, LG, OnePlus, Huawei months later.
            \item No carrier announced Android 10 upgrades until mid-December 2019 (over 100 days after release).
            \item Samsung Galaxy S10, S10+, S10e: 7.67 days on average between carrier release and first Android 10 observation.
        \end{itemize}
        \item \textbf{Android One:}
        \begin{itemize}[itemsep=0pt]
            \item 1,048 devices (under 1\% of end-user data); 75.19\% on Android 9, 24.52\% on Android 10.
            \item More up-to-date than all active devices (54.16\% Android 9, 19.20\% Android 10 in Table 3).
            \item Average 105.48 days from Android 10 release to observation; not faster than other devices.
        \end{itemize}
        \item \textbf{Project Treble:}
        \begin{itemize}[itemsep=0pt]
            \item 48.15\% of Treble devices have more than nine build numbers vs.\ 16.55\% of all active devices.
            \item Treble OS: 35.65\% Android 9, 52.43\% Android 10; 11.91\% still on Android 8.
            \item Security update latency: 19 days average for Treble vs.\ 26 days for non-Treble (7 days faster); median delay of 19 days still.
        \end{itemize}
        \item \textbf{Project Mainline:}
        \begin{itemize}[itemsep=0pt]
            \item 140,720 Mainline-supported devices (11.05\% of active devices); 86.58\% of them are Google Pixel.
            \item Direct effect not measurable (build numbers do not change).
            \item Mainline-compatible devices excluding Pixel: 1.5\%.
        \end{itemize}
        \item \textbf{Internal inconsistencies in text:} abstract says ``median delay of 11 days'' and Treble 7 days faster; Section 5.2 text says Treble ``average latency is 19 days'' and the abstract calls it average, while Section 7 also says median 19 days. Not Specified which is correct.
    \end{itemize}

    \item \entrypart{Contribution and Limitations}

    \textbf{Contributions (authors)}
    \begin{itemize}[leftmargin=*, itemsep=1pt, topsep=2pt]
        \item Extensive quantitative study at scale of manufacturer, carrier, and end-user impact on Android security update and OS upgrade rollout.
        \item Aggregation and correlation of different data sources, giving new measurements on carrier/manufacturer factors and user-incurred delays.
        \item Evaluation of Android One and Project Treble; very early assessment of Project Mainline.
        \item Claim to be the first study covering AOSP, manufacturers, carriers, and end users together.
        \item Public release of collection scripts and data.
        \item Advice: more consistency and transparency in update announcements; initiatives more independent of manufacturers/carriers; continuous tracking of progress.
    \end{itemize}

    \vspace{4pt}
    \textbf{Limitations (authors)}
    \begin{itemize}[leftmargin=*, itemsep=1pt, topsep=2pt]
        \item \textbf{Scope bias:} only the top four U.S. carriers; end-user data mostly U.S. users. Other-country carriers posted limited information.
        \item \textbf{Starting point:} measures only after a patch exists in AOSP. Does not include code-commit delays or chipset vendor delays.
        \item \textbf{Measurement accuracy:}
        \begin{itemize}[itemsep=0pt]
            \item Inconsistent carrier formats; release dates sometimes inferred.
            \item Conservative assumptions: a model could receive more updates than reported.
            \item Update presence inferred from user-agent build numbers; results based on devices where build numbers follow Android convention or the announcement.
        \end{itemize}
        \item \textbf{Model vs.\ device updates:} not all CVEs apply to each model across carriers; ``skipped'' bulletin assumed packaged into a later update.
        \item \textbf{Batch rollout:} no visibility into batch schedule; delay is hard to attribute to user or batch rollout.
        \item \textbf{Black box:} carrier--manufacturer business relationships are not public; findings are empirical observations.
        \item \textbf{User-side data:} timestamps driven by when users open the app; data gap June 11 to July 23, 2019.
    \end{itemize}

    \item \entrypart{My Analysis}

    \textbf{Good}
    \begin{itemize}[leftmargin=*, itemsep=1pt, topsep=2pt]
        \item Very large end-user dataset (9.1M identifiers) combined with carrier and AOSP data; links announcement date to observed device behavior.
        \item Handles hard data problem: inconsistent carrier announcements; processing steps are documented (Table 4, Appendix A.2).
        \item Normalized frequency metric accounts for model age.
        \item Uses significance values for correlations and reports non-significant results.
        \item Honest about conservative assumptions and limitations; checks effect of data gap.
        \item Compares findings with FTC report and prior work, including a contradicting result (locked vs.\ unlocked).
        \item Code and data released.
    \end{itemize}

    \vspace{4pt}
    \textbf{Weak points}
    \begin{itemize}[leftmargin=*, itemsep=1pt, topsep=2pt]
        \item \textbf{Proxy for update presence (observation):}
        \begin{itemize}[itemsep=0pt]
            \item Build number is not the same as patch level for all vendors; Figure 8 shows some manufacturers (e.g., Samsung) do not change build numbers on security updates.
            \item The joined dataset is 75.90\% Samsung and 15.94\% Google. Yet Samsung build numbers are said to change only on OS upgrades; how security updates are detected for Samsung is unclear from the paper (observation).
            \item Validation of the build-number-to-update mapping against ground truth: Not Specified.
        \end{itemize}
        \item \textbf{Validated vs.\ analyzed:}
        \begin{itemize}[itemsep=0pt]
            \item Manual validation of carrier fields is stated for all announcements, but error rate and number of corrections are Not Specified.
            \item 68.90\% of active devices joined with carrier data; remaining devices not analyzed for latency.
        \end{itemize}
        \item \textbf{Data selection effects (observation):}
        \begin{itemize}[itemsep=0pt]
            \item Of 3,493 raw announcements, about 53\% were kept (1,867); dropped items may not be random (missing dates, no bulletin indicator).
            \item Verizon shows only the latest three updates per model and was scraped three times; older updates may be missing, which can bias frequency and latency.
            \item Missing bulletin inferred from announcement date; for latency this lowers the estimate, but can distort per-bulletin comparisons.
            \item Outlier removal (below $-10$ or above 170 days) removed 86 updates; the appendix also mentions a 150-day filter. Inconsistent cutoffs.
        \end{itemize}
        \item \textbf{Activity bias (observation):} only users who open one social app at least monthly for 10+ months. This is a biased device population; the app only reports when opened, so update time is measured as ``first seen'', which gives an upper bound on actual update time.
        \item \textbf{User delay claim:} the paper says user-incurred delay is negligible, but also says it cannot separate user delay from batch rollout. The claim is stronger than the data supports (observation).
        \item \textbf{Causal claims (observation):} several explanations are conjecture: locked-faster result (``we conjecture'' flagship popularity), carrier prioritization, Treble speed-up due to testing procedures. Not tested.
        \item \textbf{Treble comparison (observation):} Treble vs.\ non-Treble differs by device age, manufacturer, and model tier; no control or statistical test is reported for the 7-day difference. Treble OS version includes only Android 8 or later, which biases the comparison with all devices.
        \item \textbf{Frequency definition:} $T_{end}$ is the last update observed; devices that stopped receiving updates earlier, or whose last announcement is missing, may change $N_{pot}$. Effect not discussed (observation).
        \item \textbf{Vendor bias:} U.S. carriers only; Samsung is the only manufacturer dataset (13 or 15 models, flagship Galaxy models). Manufacturer-side delay for other vendors: Not Specified.
        \item \textbf{Inconsistencies in text (observation):}
        \begin{itemize}[itemsep=0pt]
            \item 1,953 vs.\ 1,867 vs.\ 2,621 announcements in different sections.
            \item 15 vs.\ 13 Samsung models.
            \item 11-day median delay in abstract vs.\ Table/Appendix wording; 11 days median and 15.12 average.
            \item Appendix lists ``AT\&T, Sprint, T-Mobile, Sprint'' (typo).
            \item Table 4 total 3,493 minus 1,867 equals 1,626, text says 1,627 and also 1,540.
        \end{itemize}
        \item \textbf{Mainline:} effect could not be measured; the section gives device counts only.
        \item \textbf{Baselines:} no comparison against other measurement methods (e.g., SnoopSnitch-based patch presence) on the same devices; the paper compares only in text.
        \item \textbf{Not covered although described as a real problem:}
        \begin{itemize}[itemsep=0pt]
            \item Chipset vendor delays (the paper names chipset vendors as part of the chain but excludes them).
            \item Non-U.S. markets and unlocked devices outside Samsung.
            \item Actual exploitation risk during the delay.
        \end{itemize}
        \item \textbf{Time range:} 2015--2019 data only; initiatives like Mainline were new at the time of the study.
    \end{itemize}

    \vspace{4pt}
    \textbf{Research gaps}
    \begin{itemize}[leftmargin=*, itemsep=1pt, topsep=2pt]
        \item Measure delays from chipset vendors and OEM code stages, not only after AOSP release.
        \item Extend to non-U.S. carriers, regions, and more manufacturers beyond Samsung.
        \item Validate build-number proxy against real patch level on devices (e.g., device-side collection).
        \item Separate user delay from batch rollout using rollout schedule data.
        \item Control for model age, vendor, and tier when comparing Treble and non-Treble devices; test the 7-day difference statistically.
        \item Measure the real effect of Project Mainline once adoption grows.
        \item Study why carriers prioritize some models (conjectured popularity) with direct evidence.
        \item Standardized, transparent update announcements to allow independent auditing, as the authors suggest.
        \item Link rollout delay to actual exploit risk for specific CVEs.
        \item Continuous tracking of update latency after 2019 to see whether it changes.
    \end{itemize}

\end{enumerate}

\newpage

\section{Paper 12: DEFInit: An Analysis of Exposed Android Init Routines}

\begin{enumerate}

    \item \entrypart{Paper Details}

    \begin{tabular}{@{}>{\bfseries}p{3.2cm} p{11.5cm}@{}}
        Authors          & Yuede Ji, Mohamed Elsabagh, Ryan Johnson, Angelos Stavrou \\[2pt]
        Publication Year & 2021 (30th USENIX Security Symposium, August 11--13, 2021) \\[2pt]
        DOI/Link         & \url{https://www.usenix.org/conference/usenixsecurity21/presentation/ji}
    \end{tabular}

    \item \entrypart{Research Problem}
    \begin{itemize}[leftmargin=*, itemsep=1pt, topsep=2pt]
        \item \textbf{Problem:} vendors add custom Init routines (``actions'' and ``services'' in \texttt{.rc} files) to Android. Privileged pre-installed apps can trigger them by setting system properties.
        \begin{itemize}[itemsep=0pt]
            \item If the app does not protect the call site, the routine's capabilities can leak to unprivileged apps.
            \item This crosses security boundaries set by the system (confused deputy).
        \end{itemize}
        \item \textbf{Why it matters:}
        \begin{itemize}[itemsep=0pt]
            \item Init routines run with higher privileges than regular processes.
            \item Vendor changes to Init, visible at the app layer, were previously unexplored.
            \item Other Android customizations studied before: bootloaders, kernel drivers, memory management, SELinux policies.
        \end{itemize}
        \item \textbf{Goal:} measure prevalence of custom Init routines and their security impact; automate finding exposed ones.
        \item \textbf{Threat model:}
        \begin{itemize}[itemsep=0pt]
            \item Attacker-controlled third-party app installed on the device.
            \item App sends crafted IPC messages to exported components of privileged pre-installed apps.
            \item Only pre-installed apps considered as access vector to Init routines.
            \item Only Android 8.0 and higher.
        \end{itemize}
        \item \textbf{Challenges listed by the paper:}
        \begin{itemize}[itemsep=0pt]
            \item C1: processing \texttt{.rc} files statically while Init is dynamic (imports, overrides, preloaded properties).
            \item C2: determining behavior of actions, scripts, and ELF binaries (undecidable in general; estimated).
            \item C3: modeling trigger conditions and transitive dependencies between actions, services, and programs.
            \item C4: identifying privileged apps that set the properties, and exposed routines.
        \end{itemize}
    \end{itemize}

    \item \entrypart{Proposed Approach}
    \begin{itemize}[leftmargin=*, itemsep=1pt, topsep=2pt]
        \item \textbf{Main idea:} \emph{DEFInit}, a static system that unpacks firmware, extracts custom Init routines, estimates behavior, builds a dependency graph, and flags sensitive routines exposed by privileged apps. It reports to an analyst; it does not decide exploitability.
        \item \textbf{1. Unpack and extract Init routines:}
        \begin{itemize}[itemsep=0pt]
            \item Standard unpacking tools (details omitted).
            \item Parser mimics Init: starts at \texttt{/init.rc}, follows imports depth-first, substitutes variables from \texttt{.prop} and boot files.
            \item Keeps the most-specialized service definition (last with \texttt{override}); merges repeated triggers under the first one.
            \item \textbf{Custom service:} references an executable not in the base AOSP image. \textbf{Custom action:} trigger not in base AOSP.
        \end{itemize}
        \item \textbf{2. Estimate behavior of routines:}
        \begin{itemize}[itemsep=0pt]
            \item Init actions: list of Init commands with hardcoded property values.
            \item Shell scripts: custom tracer that dry-runs in a sandbox using Bash trace mode (\texttt{bash -x}).
            \begin{itemize}[itemsep=0pt]
                \item Taints arguments from the \texttt{.rc} call site; evaluates only conditions depending on them.
                \item Ignores \texttt{sleep}; masks return codes.
            \end{itemize}
            \item ELF binaries: static traces of API calls from the inter-procedural CFG, depth-first over simple paths; inter-procedural Def-Use to resolve constant arguments; also extracts strings that look like system commands.
            \item Annotates routines with behavior categories using a curated list: 53 sensitive commands/APIs in 8 categories (Table 1), shortlisted from 473 commands and 4,259 Bionic APIs (137 commands and 64 APIs after first filter).
            \item 116 pattern-matching rules (one to three calls each) for specific behaviors, e.g., sensitive data command followed by a move to external storage.
        \end{itemize}
        \item \textbf{3. Init Dependency Graph (IDG):}
        \begin{itemize}[itemsep=0pt]
            \item One IDG per firmware (Algorithm 1).
            \item Nodes: trigger, service, executable. Edges: call edges and conjunction edges.
            \item Extra trigger node \texttt{ctl.start=<service>} for each service.
            \item Fall-through edges from executables that set Init properties to matching triggers.
        \end{itemize}
        \item \textbf{4. Identify exposed routines:}
        \begin{itemize}[itemsep=0pt]
            \item Finds app call sites of \texttt{set(key, value)} APIs (e.g., \texttt{android.os.SystemProperties.set}).
            \item ICFG with Def-Use; context- and flow-sensitive extraction of keys and values per call stack (Algorithm 2).
            \item Also tracks properties read by privileged apps and local properties shared between apps with the same UID.
            \item Walks the IDG from written property triggers; conjunction passes only if all conjuncts satisfied; aggregates behaviors of terminal executables.
            \item Control-flow reachability from exported entry points of the privileged app to the call site.
        \end{itemize}
        \item \textbf{Implementation:} 7K SLOC Python on top of BinaryNinja (ELF) and Kryptowire's internal Android static analysis engine (apps).
        \item \textbf{Soundness trade-off:} analysis tuned to avoid false positives at the cost of soundness (e.g., virtual and indirect calls resolved only if one candidate callee).
    \end{itemize}

    \item \entrypart{Dataset and Evaluation}
    \begin{itemize}[leftmargin=*, itemsep=1pt, topsep=2pt]
        \item \textbf{Firmware data:}
        \begin{itemize}[itemsep=0pt]
            \item 259 stock firmware, Android 8.0 to 11.0, from 21 vendors; 64,632 pre-installed apps (average 262 per firmware).
            \item By version: Android 8: 93 firmware; 9: 75; 10: 75; 11: 16.
            \item Per vendor firmware count from 1 to 11 as stated in Table 2 (``total (min;max;avg)'' 259 (1;11;4)); Table A.1 lists 5 to 22 for some vendors.
            \item Sources: vendor sites and third-party sites (Appendix A.2).
        \end{itemize}
        \item \textbf{Ground truth / validation:}
        \begin{itemize}[itemsep=0pt]
            \item Manual verification of code paths: no runtime checks (UID, permission, signature, package name), property set correctly, property triggers routine, executable does the sensitive action.
            \item Dynamic exploit developed for 53 findings on stock devices the authors could obtain; no user interaction other than installing and running the attack app.
        \end{itemize}
        \item \textbf{Setup:} Ubuntu 20.04 server, 8-core Intel Xeon E5-4620 2.20GHz, 512 GiB RAM.
        \item \textbf{Metrics:}
        \begin{itemize}[itemsep=0pt]
            \item Number of total, custom, exposed, exposed sensitive routines (per firmware and per version).
            \item UIDs and SELinux domains of routines; executables used (scripts and binaries).
            \item Exposed sensitive routines by behavior category; verified vulnerabilities by impact.
            \item Runtime: unpacking, Init processing, app analysis, end-to-end.
        \end{itemize}
        \item \textbf{Manual effort stated:} one day for three people to shortlist commands; about four workdays for one person to develop rules; about seven workdays for one person to verify findings.
        \item \textbf{Baseline comparison:} Not Specified. Related work discussed only in text; no tool compared.
    \end{itemize}

    \item \entrypart{Key Findings}
    \begin{itemize}[leftmargin=*, itemsep=1pt, topsep=2pt]
        \item \textbf{Prevalence (Table 3):}
        \begin{itemize}[itemsep=0pt]
            \item 58,523 Init routines in total (223 per firmware); 38,846 (66\%) custom, about 133 per firmware.
            \item AOSP has about 130 Init routines on average.
            \item More than 50\% of firmware had at least 600 IDG nodes, 500 edges, and 300 paths of length three or higher.
        \end{itemize}
        \item \textbf{Exposed routines:}
        \begin{itemize}[itemsep=0pt]
            \item 1,947 exposed routines, all custom, in 101 firmware from 13 vendors (about eight per firmware).
            \item 515 exposed routines run at least one sensitive command/API (``exposed sensitive'').
            \item By version (exposed / exposed sensitive): Android 8: 305 / 81; 9: 600 / 166; 10: 911 / 244; 11: 131 / 24.
            \item Firmware from vendors closest to AOSP (Google) had no exposed Init routines.
        \end{itemize}
        \item \textbf{Privileges of exposed routines (Tables 4, 5):}
        \begin{itemize}[itemsep=0pt]
            \item More than 90\% ran as \texttt{root} or \texttt{system}.
            \item Majority under default SELinux domains (\texttt{init}, \texttt{vendor\_init}).
            \item UID \texttt{root}: 1,342 exposed (354 sensitive); \texttt{system}: 272 (2); \texttt{default (root)}: 188 (113).
            \item Authors say it is unclear why services needed these privileged defaults; it appeared vendors followed the path of least resistance.
        \end{itemize}
        \item \textbf{Executables (Table 6):}
        \begin{itemize}[itemsep=0pt]
            \item 2,685 unique custom scripts, 310 known scripts, 16,772 custom binaries, 18,419 known binaries used by Init routines.
            \item Exposed routines called 1,161 custom scripts, 414 custom binaries, 181 known binaries.
            \item 581 custom scripts and 410 custom binaries invoke a sensitive command.
        \end{itemize}
        \item \textbf{Exposed sensitive behaviors, no user interaction (Table 7):}
        \begin{itemize}[itemsep=0pt]
            \item Sensitive data: 336 routines (298 direct, 38 indirect), 109 apps, 65 firmware, 11 vendors.
            \item Networking: 65; UI interaction: 48; Process management: 30; Device settings: 13.
            \item Package management: 9; Permission control: 9; Power management: 6.
            \item Total (unique): 430 routines, 173 apps, 101 firmware, 13 vendors.
            \item Permission control: \texttt{setenforce 0} found in seven instances in six firmware, one vendor; one without user interaction, six via a button click.
            \item Kernel module loaded from a writeable path in one case (could allow overwrite by a third-party app; authors say ``may'').
        \end{itemize}
        \item \textbf{Verified vulnerabilities (Table 8, no user interaction):}
        \begin{itemize}[itemsep=0pt]
            \item 89 vulnerabilities in 34 unique apps, 30 firmware, 6 vendors, Android 8, 9, 10.
            \item Read system logs 11; read dumpsys 10; read kernel logs 10; load kernel module 9; sniff Wi-Fi traffic 8.
            \item Sniff modem traffic 7; read dumpstate 7; read radio logs 6; record screen 5; reboot device 5.
            \item Read Wi-Fi passwords 3; read bugreport 3; sniff Bluetooth traffic 2; write to node device 2; disable SELinux 1.
            \item Three vendors have so far confirmed findings (49 flaws in 11 firmware); disclosure ongoing.
        \end{itemize}
        \item \textbf{Verified, needs user interaction (Table A.5):} 134 vulnerabilities, 52 apps, 35 firmware, 9 vendors.
        \item \textbf{Routines exposed only via GUI (Tables A.3, A.4):} 435 exposed routines, 130 exposed sensitive.
        \item \textbf{Case studies:}
        \begin{itemize}[itemsep=0pt]
            \item \textbf{Disable SELinux:} custom \texttt{wifitest} service runs a root shell script that calls \texttt{setenforce 0}. One firmware allows background launch without interaction; five firmware (six apps) need a button click.
            \item \textbf{Modem and network traffic:} Android 9 firmware from two vendors; third-party app sends IPC to an exported broadcast receiver; root Init services capture traffic to external storage. Vendors said a common chipset provider introduced it by mistake.
            \item \textbf{Sensitive logs:} three vendors; root script runs \texttt{logcat}; two also expose \texttt{dumpsys}; output on external storage.
            \item \textbf{Screen recording:} one vendor, five firmware; \texttt{screenrecord} for 30 min, file on external storage.
        \end{itemize}
        \item \textbf{Runtime:} unpacking about 5 min on average; Init processing and IDG about 30 min; app analysis about 7 min; 90\% of firmware finished in under 70 min end-to-end.
        \item \textbf{Vendor share (Table A.2, anonymized):} Vendor A has 41\% of verified vulnerabilities with 8.9\% market share.
    \end{itemize}

    \item \entrypart{Contribution and Limitations}

    \textbf{Contributions (authors)}
    \begin{itemize}[leftmargin=*, itemsep=1pt, topsep=2pt]
        \item \textbf{Novel system:} DEFInit, an automated system to process firmware, find Init routines, estimate behavior, find routines exposed by privileged apps, and highlight risky ones.
        \item \textbf{Systematic study:} first comprehensive study of vendor customization of Android Init routines triggerable from privileged apps (259 firmware, Android 8 to 11, 21 vendors).
        \item \textbf{New findings:} prevalence and impact of custom Init routines; multiple exploitable instances with serious security and privacy impact.
        \item Countermeasures proposed: default unprivileged user and SELinux domain for Init children; CDD rules against privileged custom routines triggerable from outside Init; block third-party to pre-installed app interaction by default; SELinux default to stop Init executables writing to external storage; user warning for severe actions.
    \end{itemize}

    \vspace{4pt}
    \textbf{Limitations (authors)}
    \begin{itemize}[leftmargin=*, itemsep=1pt, topsep=2pt]
        \item DEFInit does not decide exploitability; it brings instances to an analyst.
        \item Analysis is conservative and unsound: under-approximated ICFGs; reflection, flows through containers, inter-component communication, and managed/native (JNI) flows not handled; permission-protected components treated as unexported.
        \item No dynamic access control checks (permission, UID, confirmation dialog) modeled; manually checked only for Table 8 findings. Tables 7 and A.6 should be read with this in mind.
        \item Trace collection is static; no firmware emulation.
        \item Multi-app behaviors (one app starts action, another moves files) not detected; SELinux constraints for them left to future work.
        \item Assumes vendors configured SELinux labels and transitions properly for custom routines to work.
        \item Firmware dataset not uniform across vendors and versions; some vendors (e.g., Itel) have small images and few routines; unpacking of proprietary formats may have missed files; vendor differences may not be statistically significant.
        \item Only live-device exploitation of a limited number of findings, due to lack of devices.
        \item Only pre-installed apps as access vector; only Android 8.0+.
        \item Manual effort needed for rule development and verification.
    \end{itemize}

    \item \entrypart{My Analysis}

    \textbf{Good}
    \begin{itemize}[leftmargin=*, itemsep=1pt, topsep=2pt]
        \item Large dataset: 259 firmware, 21 vendors, 64,632 apps, four Android versions.
        \item New attack surface: Init routines reached through privileged apps. Clear, simple motivating example (Figure 1).
        \item Pipeline covers all three executable types (actions, shell scripts, ELF).
        \item IDG captures transitive trigger dependencies and conjunctions.
        \item Context- and flow-sensitive extraction of property keys and values per call stack.
        \item Manual verification checklist is explicit; part of findings confirmed with working exploits on stock devices.
        \item Authors state the unsound analysis choice and the lack of dynamic access-control checks openly.
        \item Concrete countermeasures; some vendor confirmation of findings; fast runtime (90\% under 70 min).
        \item Splits findings by user interaction needed (Tables 8 and A.5).
    \end{itemize}

    \vspace{4pt}
    \textbf{Weak points}
    \begin{itemize}[leftmargin=*, itemsep=1pt, topsep=2pt]
        \item \textbf{Validated vs.\ analyzed (observation):}
        \begin{itemize}[itemsep=0pt]
            \item 1,947 exposed routines and 515 sensitive; only 89 (no interaction) and 134 (with interaction) manually verified. Rest are candidates.
            \item Only 53 of the 89 were dynamically exploited; the other 36 were checked by manual path inspection only. Which 36 and why: Not Specified.
            \item Precision and false-positive rate of the 1,947 and 515 figures: Not Specified. Without dynamic access control checks the numbers may include protected paths.
            \item Recall: Not Specified; the analysis is unsound by design.
        \end{itemize}
        \item \textbf{Design/evaluation overlap:}
        \begin{itemize}[itemsep=0pt]
            \item The 116 rules and 53-command list were written by analysts who also analyzed the findings; same data used to develop rules and report results. No held-out set (Not Specified).
            \item Rule examples target the same behaviors listed in verified results (logcat, tcpdump, screenrecord, setenforce), so ``sensitive'' is defined by the authors' own list.
        \end{itemize}
        \item \textbf{Vendor and dataset bias (observation):}
        \begin{itemize}[itemsep=0pt]
            \item Findings are heavily concentrated: 89 verified in 6 vendors; Vendor A has 41\% of verified vulnerabilities (anonymized). Vendor names withheld.
            \item The ``all seven \texttt{setenforce 0} instances from one vendor'' result and the chipset-provider note suggest a few root causes, not 89 independent bugs. Deduplication by root cause: Not Specified.
            \item Firmware count per vendor is uneven; Table A.1 gives 1 firmware min in Table 2 but vendor totals 4 to 22. Firmware from third-party sites; authenticity not checked (Not Specified).
            \item Only stock firmware; carrier variants not covered.
        \end{itemize}
        \item \textbf{Heuristics and manual steps:}
        \begin{itemize}[itemsep=0pt]
            \item Behavior estimation relies on a hand-picked command list, regex-like rules, and a shell dry-run with masked return codes and ignored \texttt{sleep}. Accuracy of the shell tracer and ELF string matching: Not Specified.
            \item Rules and verification need expert time (about 7 workdays for verification alone).
        \end{itemize}
        \item \textbf{Missing baselines and metrics:}
        \begin{itemize}[itemsep=0pt]
            \item No comparison with other tools (e.g., FirmScope-style capability analysis); no ablation of IDG or Def-Use components.
            \item ``Exposed'' counts routines, not distinct vulnerabilities or apps; the same routine reached via multiple apps or entry points may be counted several times (Table 7 total 430 unique; Table 3 total 1,947 exposed). Relation between these two numbers is unclear (observation).
            \item Number of firmware with each finding not normalized by the number of firmware per vendor.
        \end{itemize}
        \item \textbf{Inconsistencies in text (observation):}
        \begin{itemize}[itemsep=0pt]
            \item Abstract says 13 vendors for exposed routines, 6 vendors for verified; consistent with body. Table 7 gives 430 sensitive unique routines while text says 515 exposed sensitive in 101 firmware.
            \item Table 4 total custom is 32,430 while Table 3 total custom is 38,846.
            \item Table A.3 total row shows ``130 (9;5;0.5)'', where min exceeds max.
            \item Table 3 reports Android 11 with only 16 firmware.
            \item Section 6.3 says seven setenforce instances in six firmware; Table 8 lists 1 verified without interaction; Table A.5 lists 6 with interaction in 5 firmware. Counts of firmware (six vs.\ five) differ.
        \end{itemize}
        \item \textbf{Single case / device:} dynamic exploitation only on devices the authors could buy; the rest of the findings were not run on real hardware.
        \item \textbf{Causal and interpretive claims (observation):}
        \begin{itemize}[itemsep=0pt]
            \item ``Vendors followed the path of least resistance'' and ``introduced by mistake'' are speculation, except for the one vendor note on the chipset case.
            \item Claim that SELinux should not block flows is an assumption, tested only in the dynamically verified cases (none blocked).
            \item Impact statements (e.g., arbitrary code execution through a writeable kernel module path) are possible effects, not shown.
            \item Statement that Google has no exposed routines may reflect the data (Google firmware count 20), not a general rule.
        \end{itemize}
        \item \textbf{Not covered although described as a real problem:}
        \begin{itemize}[itemsep=0pt]
            \item Other access vectors (framework APIs, native services, USB, etc.) are out of scope, though the paper cites them as risks.
            \item Multi-app attack chains and dynamic access-control checks; both listed as future work.
            \item Real-world exploitation or user impact; responsible disclosure not finished (three vendors confirmed).
        \end{itemize}
    \end{itemize}

    \vspace{4pt}
    \textbf{Research gaps}
    \begin{itemize}[leftmargin=*, itemsep=1pt, topsep=2pt]
        \item Automatic modeling of dynamic access control (permission, UID, signature checks, dialogs) on paths to property setters.
        \item Detect multi-app behaviors and model SELinux transitions for them.
        \item Better trace collection: firmware re-hosting or emulation, reflection and JNI handling; measure recall of the static approach.
        \item Automate rule creation (data mining on labeled traces), which the authors suggest.
        \item Other vendor-added access vectors to Init routines (framework APIs, native services, USB, GUI cloaking).
        \item Evaluate precision on a larger dynamically verified sample; deduplicate findings by root cause (vendor, chipset vendor).
        \item Study carrier builds, newer Android versions, and whether proposed countermeasures (default unprivileged Init children, CDD rules) are adopted.
        \item Compare DEFInit against other pre-installed-app and capability-leak tools on the same firmware.
        \item Measure effect of countermeasures on exploitability, e.g., blocking external-storage writes by Init executables.
    \end{itemize}

\end{enumerate}

\end{document}